% Leçon 29 — Expression orale : corruption — Chinois 1ère A — Cours signé OnBuch+
% Programme officiel MINESEC. Compiler avec : tectonic ZHA29-oral-corruption.tex
\newcommand{\DOCMATIERE}{Chinois}
\newcommand{\DOCNIVEAU}{1ère A}
\newcommand{\DOCMODULE}{Module 3 : Citoyenneté et ouverture au monde}
\newcommand{\DOCLECON}{ZHA29}
\newcommand{\DOCTITRE}{Leçon 29 — Expression orale : corruption}
% OnBuch+ — préambule « cours signés » ENRICHI (compilé avec tectonic / XeLaTeX)
% Système visuel « Pop » d'OnBuch+ : fond crème (jamais blanc), bordures encre
% épaisses, ombres « dures » décalées, titres Archivo Black en capitales.
% Version MATHÉMATIQUES : graphes (pgfplots/tikz), boîtes pédagogiques
% (définition, propriété, méthode, attention, le savais-tu, exercice résolu,
% exercice, TP, bilan), page de garde et signature OnBuch+ ancrée en bas.
%
% Macros attendues dans main.tex AVANT \input{preamble} :
%   \DOCMATIERE  \DOCNIVEAU  \DOCMODULE  \DOCLECON (numéro)  \DOCTITRE
\documentclass[11pt]{article}

\usepackage{fontspec}
\usepackage[french]{babel} % français = langue principale ; le chinois via \zh{..} / textebox (xeCJK)
\usepackage{xeCJK}
\usepackage{pifont} % \ding{51} \ding{55} utilisés par les modèles
\usepackage[a4paper,margin=2cm,top=3.4cm,bottom=2.8cm,headheight=1.4cm,footskip=1.3cm]{geometry}
\usepackage{amsmath,amssymb,amsfonts}
\usepackage{mathtools} % \xmapsto, \coloneqq... souvent utilisés par les modèles
\usepackage{cancel} % simplifications barrées (bilans de forces)
% \abs{z} (module, valeur absolue) et \norm{u} (norme), utilisés par les modèles
\providecommand{\abs}{}\let\abs\relax\DeclarePairedDelimiter{\abs}{\lvert}{\rvert}
\providecommand{\norm}{}\let\norm\relax\DeclarePairedDelimiter{\norm}{\lVert}{\rVert}
\usepackage[table]{xcolor} % [table] : fournit aussi \rowcolors (alternance de lignes)
\usepackage{pagecolor}
\usepackage{tikz}
\usetikzlibrary{arrows.meta,positioning,calc,shapes.geometric,shapes.misc,shapes.arrows,decorations.pathmorphing,decorations.pathreplacing,decorations.markings,patterns,shadows,mindmap,trees,fit,backgrounds,matrix,intersections,angles,quotes}
\usepackage{pgfplots}
\usepackage[version=4]{mhchem} % équations nucléaires \ce{^{235}_{92}U}
\usepackage[european,siunitx]{circuitikz} % schémas électriques
\pgfplotsset{compat=1.18}
\usepgfplotslibrary{fillbetween,groupplots} % aires entre courbes (calcul intégral)
\usepackage{enumitem}
\usepackage{fancyhdr}
% [most] charge toutes les bibliothèques tcolorbox (listings, minted, raster,
% documentation...) et gonfle inutilement la mémoire TeX sur un document long
% (~300 boîtes) : on ne charge que ce qui est réellement utilisé.
\usepackage[skins,breakable]{tcolorbox}
\usepackage{graphicx}
\usepackage{colortbl}
\usepackage{booktabs}
\usepackage{tabularx}
\usepackage{diagbox} % case d'en-tête diagonale des tableaux à double entrée
% Colonnes extensibles centrée / gauche / droite (C, L, R), souvent utilisées par les modèles
\newcolumntype{C}{>{\centering\arraybackslash}X}
\newcolumntype{L}{>{\raggedright\arraybackslash}X}
\newcolumntype{R}{>{\raggedleft\arraybackslash}X}
\usepackage{array}
\usepackage{multicol}
\usepackage{multirow}
\usepackage{siunitx}
% parse-numbers=false : accepte aussi \SI{2,5 \times 10^{-3}}{...}
\sisetup{locale=FR,output-decimal-marker={,},group-minimum-digits=4,parse-numbers=false}
\DeclareSIUnit\ohm{\ensuremath{\symup{\Omega}}} % U+2126 absent de la police
\DeclareSIUnit\division{div} % réglages d'oscilloscope (ms/div, V/div)
\DeclareSIUnit\tr{tr} % tours (vitesse de rotation, ex. \SI{1500}{\tr\per\minute})
\DeclareSIUnit\molar{M} % concentration molaire (ex. \SI{3}{\molar}, \SI{1}{\micro\molar})
\DeclareSIUnit\an{an} % année (le modèle confond parfois avec \year, un compteur TeX) — ex. \SI{5}{\kilo\meter\per\an}
\DeclareSIUnit\FCFA{FCFA} % franc CFA (ex. \SI{500}{\FCFA\per\kilo\gram})
\DeclareSIUnit\degreeBrix{\textdegree Brix} % teneur en sucre d'un sirop/jus (réfractométrie)
\DeclareSIUnit\month{mois} % le modèle utilise parfois \month, un compteur TeX, comme unité de durée
\usepackage{qrcode}
% \degree : commande fréquemment utilisée par les modèles pour un angle ou une
% température en dehors de \SI{}{} (non fournie par les paquets chargés).
\providecommand{\degree}{^{\circ}}
% \qed (fin de démonstration), utilisé par les modèles sans amsthm
\providecommand{\qed}{\hfill$\square$}
% \barre : double barre des valeurs interdites dans un tableau de variations (tkz-tab)
\providecommand{\barre}{\ensuremath{\|}}
% environnement proof (amsthm non chargé) utilisé par les modèles
\ifdefined\proof\else\newenvironment{proof}[1][Démonstration]{\par\noindent\textit{#1.}\ }{\qed\par}\fi
\usepackage{lastpage}
\usepackage{hyperref}
\hypersetup{hidelinks}

% ============================================================
% Palette « Pop » OnBuch+ — mêmes hex que la classe P (plus_ui.dart)
% ============================================================
\definecolor{poporange}{HTML}{F59321}
\definecolor{popdark}{HTML}{E07A0C}
\definecolor{popcream}{HTML}{FFF6E8}
\definecolor{popink}{HTML}{1C1714}
\definecolor{poppurple}{HTML}{7A5AE0}
\definecolor{popgreen}{HTML}{2E9E6A}
\definecolor{poppink}{HTML}{E0457B}
\definecolor{popblue}{HTML}{2F7DE1}
\definecolor{popgold}{HTML}{C98A12}
\definecolor{popmuted}{HTML}{8A7F76}
\definecolor{popline}{HTML}{EADFD2}
\definecolor{popgray}{HTML}{8A7F76}
\colorlet{popgrey}{popgray}
% Alias tolérants (le modèle nomme parfois les couleurs différemment)
\colorlet{popwhite}{white}
\colorlet{popblack}{popink}
\colorlet{popred}{poppink}
\colorlet{popyellow}{popgold}
% Teintes claires (fonds de tableaux/encadrés) — le modèle nomme parfois
% les couleurs avec un suffixe L (ex. popblueL) sans les avoir définies.
\colorlet{poporangeL}{poporange!20}
\colorlet{popdarkL}{popdark!20}
\colorlet{poppurpleL}{poppurple!20}
\colorlet{popgreenL}{popgreen!20}
\colorlet{poppinkL}{poppink!20}
\colorlet{popblueL}{popblue!20}
\colorlet{popgoldL}{popgold!20}
\colorlet{popredL}{popred!20}
\colorlet{popgrayL}{popgray!20}
% Teintes claires pour fonds de tableaux / zones
\colorlet{popblueL}{popblue!10!white}
\colorlet{poporangeL}{poporange!14!white}
\colorlet{popgreenL}{popgreen!12!white}
\colorlet{poppurpleL}{poppurple!12!white}
\colorlet{poppinkL}{poppink!12!white}
\colorlet{popgoldL}{popgold!14!white}

\pagecolor{popcream}

% ============================================================
% Typographie
% ============================================================
\defaultfontfeatures{Ligatures=TeX}
\setmainfont{PlusJakartaSans-Medium.ttf}[Path=../../../fonts/,BoldFont=PlusJakartaSans-Bold.ttf]
% Symboles, API et emojis absents de la police principale : repli sur DejaVu Sans (caractères actifs)
\newfontface\symfont{DejaVuSans.ttf}[Path=../../../fonts/]
\makeatletter
\newcommand\symactive[1]{\catcode#1=13 \begingroup\lccode`\~=#1 \lowercase{\endgroup\def~}{{\symfont\char#1}}}
\newcommand\symblank[1]{\catcode#1=13 \begingroup\lccode`\~=#1 \lowercase{\endgroup\def~}{}}
\newcommand\symhyph[1]{\catcode#1=13 \begingroup\lccode`\~=#1 \lowercase{\endgroup\def~}{\mbox{-}}}
\newcount\symc
\newcommand\symrange[2]{\symc=#1 \@whilenum\symc<#2 \do{\expandafter\symactive\expandafter{\the\symc}\advance\symc by 1 }}
\newcommand\symblankrange[2]{\symc=#1 \@whilenum\symc<#2 \do{\expandafter\symblank\expandafter{\the\symc}\advance\symc by 1 }}
\makeatother
\symrange{"0250}{"02FF}\symactive{"2713}\symactive{"2714}\symactive{"2717}\symactive{"2718}\symactive{"2610}\symactive{"2611}\symactive{"2612}
\symrange{"2600}{"27BF}
\catcode"2705=13 \begingroup\lccode`\~="2705 \lowercase{\endgroup\def~}{{\symfont\char"2713}}\symblank{"26BD}\symblank{"26A1}
\symrange{"2460}{"257F}\symactive{"223C}\symactive{"2248}\symactive{"2260}
\symactive{"01F8}\symactive{"01F9}\symactive{"1E3E}\symactive{"1E3F}
\symhyph{"2011}\symhyph{"2E1A}\symblankrange{"1F300}{"1FAFF}\symblank{"FE0F}
% Caractères chinois : WenQuanYi Zen Hei (sous-ensemble GB2312 + pinyin), gras/italique simulés
\setCJKmainfont{WenQuanYiZenHei-subset.ttf}[Path=../../../fonts/,AutoFakeBold=3,AutoFakeSlant=0.2]
\xeCJKsetup{CJKspace=false}
% Pinyin : voyelles à caron (ǎ ǐ ǒ ǔ ǖ ǘ ǚ ǜ) absentes de la police principale -> caractères actifs
\newfontface\pyfont{WenQuanYiZenHei-subset.ttf}[Path=../../../fonts/]
\newcommand\pyactive[1]{\catcode#1=13 \begingroup\lccode`\~=#1 \lowercase{\endgroup\def~}{{\pyfont\char#1}}}
\pyactive{"01CD}\pyactive{"01CE}\pyactive{"01CF}\pyactive{"01D0}\pyactive{"01D1}\pyactive{"01D2}\pyactive{"01D3}\pyactive{"01D4}\pyactive{"01D5}\pyactive{"01D6}\pyactive{"01D7}\pyactive{"01D8}\pyactive{"01D9}\pyactive{"01DA}\pyactive{"01DB}\pyactive{"01DC}
% Maths en sans-serif (Fira Math) avec chiffres et lettres droites de Plus
% Jakarta, pour que les formules se fondent dans le texte.
\usepackage[math-style=ISO,bold-style=ISO]{unicode-math}
\setmathfont{FiraMath-Regular.otf}[Path=../../../fonts/]
\setmathfont{PlusJakartaSans-Medium.ttf}[Path=../../../fonts/,range={up/num,up/{Latin,latin},\mathup}]
\usepackage{icomma} % virgule décimale sans espace en mode math
% Caractères mathématiques Unicode absents des polices de texte (Plus Jakarta,
% Archivo Black) : rendus par leur équivalent mathématique, en texte comme en
% formule — sinon ils s'affichent en carré vide (ex. « Arithmétique dans ℤ »).
\usepackage{newunicodechar}
\newunicodechar{ℝ}{\ensuremath{\mathbb{R}}}
\newunicodechar{ℤ}{\ensuremath{\mathbb{Z}}}
\newunicodechar{ℂ}{\ensuremath{\mathbb{C}}}
\newunicodechar{ℕ}{\ensuremath{\mathbb{N}}}
\newunicodechar{ℚ}{\ensuremath{\mathbb{Q}}}
\newunicodechar{𝔻}{\ensuremath{\mathbb{D}}}
\newunicodechar{α}{\ensuremath{\alpha}}
\newunicodechar{β}{\ensuremath{\beta}}
\newunicodechar{γ}{\ensuremath{\gamma}}
\newunicodechar{δ}{\ensuremath{\delta}}
\newunicodechar{ε}{\ensuremath{\varepsilon}}
\newunicodechar{θ}{\ensuremath{\theta}}
\newunicodechar{λ}{\ensuremath{\lambda}}
\newunicodechar{μ}{\ensuremath{\mu}}
\newunicodechar{σ}{\ensuremath{\sigma}}
\newunicodechar{τ}{\ensuremath{\tau}}
\newunicodechar{φ}{\ensuremath{\varphi}}
\newunicodechar{ψ}{\ensuremath{\psi}}
\newunicodechar{ω}{\ensuremath{\omega}}
\newunicodechar{Σ}{\ensuremath{\Sigma}}
% majuscules produites par \MakeUppercase dans les titres (α-aminés -> Α)
\newunicodechar{Α}{\ensuremath{\alpha}}
\newunicodechar{Β}{\ensuremath{\beta}}
\newunicodechar{Γ}{\ensuremath{\Gamma}}
\newunicodechar{Δ}{\ensuremath{\Delta}}
\newunicodechar{Ε}{\ensuremath{\varepsilon}}
\newunicodechar{Θ}{\ensuremath{\Theta}}
\newunicodechar{Λ}{\ensuremath{\Lambda}}
\newunicodechar{Μ}{\ensuremath{\mu}}
\newunicodechar{Τ}{\ensuremath{\tau}}
\newunicodechar{Φ}{\ensuremath{\Phi}}
\newunicodechar{Ψ}{\ensuremath{\Psi}}
\newunicodechar{Ω}{\ensuremath{\Omega}}
\newunicodechar{↦}{\ensuremath{\mapsto}}
\newunicodechar{∈}{\ensuremath{\in}}
\newunicodechar{∉}{\ensuremath{\notin}}
\newunicodechar{∧}{\ensuremath{\wedge}}
\newunicodechar{⇔}{\ensuremath{\Leftrightarrow}}
\newunicodechar{⟺}{\ensuremath{\Longleftrightarrow}}
\newunicodechar{⇒}{\ensuremath{\Rightarrow}}
\newunicodechar{⟹}{\ensuremath{\Longrightarrow}}
\newunicodechar{∘}{\ensuremath{\circ}}
\newunicodechar{∪}{\ensuremath{\cup}}
\newunicodechar{∩}{\ensuremath{\cap}}
\newunicodechar{≡}{\ensuremath{\equiv}}
\newunicodechar{⊂}{\ensuremath{\subset}}
\newunicodechar{⊥}{\ensuremath{\perp}}
\newunicodechar{∃}{\ensuremath{\exists}}
\newunicodechar{∀}{\ensuremath{\forall}}
\newunicodechar{∛}{\ensuremath{\sqrt[3]{\,}}}
\newunicodechar{∜}{\ensuremath{\sqrt[4]{\,}}}
\newunicodechar{⃗}{\ensuremath{{}^{\to}}}
\newunicodechar{ⁿ}{\ifmmode^{n}\else\textsuperscript{\ensuremath{n}}\fi}
\newunicodechar{ˣ}{\ifmmode^{x}\else\textsuperscript{\ensuremath{x}}\fi}
\newunicodechar{ᵇ}{\ifmmode^{b}\else\textsuperscript{\ensuremath{b}}\fi}
\newunicodechar{ʳ}{\ifmmode^{r}\else\textsuperscript{\ensuremath{r}}\fi}
\newunicodechar{ᵘ}{\ifmmode^{u}\else\textsuperscript{\ensuremath{u}}\fi}
\newunicodechar{ʸ}{\ifmmode^{y}\else\textsuperscript{\ensuremath{y}}\fi}
\newunicodechar{ᵃ}{\ifmmode^{a}\else\textsuperscript{\ensuremath{a}}\fi}
\newunicodechar{ᵉ}{\ifmmode^{e}\else\textsuperscript{\ensuremath{e}}\fi}
\newunicodechar{ᵏ}{\ifmmode^{k}\else\textsuperscript{\ensuremath{k}}\fi}
\newunicodechar{ᵖ}{\ifmmode^{p}\else\textsuperscript{\ensuremath{p}}\fi}
\newunicodechar{ᴺ}{\ifmmode^{N}\else\textsuperscript{\ensuremath{N}}\fi}
\newunicodechar{ᶜ}{\ifmmode^{c}\else\textsuperscript{\ensuremath{c}}\fi}
\newunicodechar{ʰ}{\ifmmode^{h}\else\textsuperscript{\ensuremath{h}}\fi}
\newunicodechar{ᵠ}{\ifmmode^{q}\else\textsuperscript{\ensuremath{q}}\fi}
\newunicodechar{ₙ}{\ifmmode_{n}\else\textsubscript{\ensuremath{n}}\fi}
\newunicodechar{ₐ}{\ifmmode_{a}\else\textsubscript{\ensuremath{a}}\fi}
\newunicodechar{ᵢ}{\ifmmode_{i}\else\textsubscript{\ensuremath{i}}\fi}
\newunicodechar{ᵣ}{\ifmmode_{r}\else\textsubscript{\ensuremath{r}}\fi}
\newunicodechar{ₖ}{\ifmmode_{k}\else\textsubscript{\ensuremath{k}}\fi}
\newunicodechar{ᵦ}{\ifmmode_{\beta}\else\textsubscript{\ensuremath{\beta}}\fi}
\newunicodechar{ₓ}{\ifmmode_{x}\else\textsubscript{\ensuremath{x}}\fi}
\newfontfamily\popdisplay{ArchivoBlack-Regular.ttf}[Path=../../../fonts/]
\newfontfamily\poplogo{SpaceGrotesk-Bold.ttf}[Path=../../../fonts/]
\newfontfamily\popsemi{PlusJakartaSans-SemiBold.ttf}[Path=../../../fonts/]
\newfontfamily\popxbold{PlusJakartaSans-ExtraBold.ttf}[Path=../../../fonts/]
\newfontfamily\popxboldls{PlusJakartaSans-ExtraBold.ttf}[Path=../../../fonts/,LetterSpace=3.0]

\setlist[enumerate]{leftmargin=1.7em,itemsep=5pt,topsep=4pt}
\setlist[itemize]{leftmargin=1.7em,itemsep=4pt,topsep=4pt,label={\color{poporange}\textbullet}}
\setlength{\parindent}{0pt}
\setlength{\parskip}{5pt}
\renewcommand{\arraystretch}{1.25}

% pgfplots : style Pop par défaut
\pgfplotsset{
  every axis/.append style={axis line style={popink,line width=1pt},tick style={popink},
    grid style={popline,line width=0.5pt},label style={font=\small},tick label style={font=\footnotesize}},
  popcurve/.style={poporange,line width=1.8pt,no markers,smooth},
}
% Même style disponible dans un \draw TikZ simple (hors pgfplots)
\tikzset{popcurve/.style={poporange,line width=1.8pt}}
\tikzset{popgrid/.style={popline,very thin}}
\colorlet{popcurve}{poporange} % le nom du style est parfois utilisé comme couleur

% ============================================================
% Pill / squiggle
% ============================================================
\newcommand{\poppill}[3]{%
  \tikz[baseline=-0.55ex]\node[draw=popink,line width=1.2pt,rounded corners=2.6pt,%
    fill=#2,inner xsep=6.5pt,inner ysep=3pt]{%
    \popxboldls\fontsize{7}{7}\selectfont\color{#3}\MakeUppercase{#1}};%
}
\newcommand{\popsquiggle}[1]{%
  \tikz[baseline=-0.4ex]{
    \draw[#1,line width=2.6pt,line cap=round]
      plot [smooth] coordinates {(0,0.09) (0.34,0.01) (0.68,0.21) (1.02,0.01) (1.36,0.10)};
  }%
}
% Mot-clé surligné (marqueur orange sous le texte)
\newcommand{\cle}[1]{{\popxbold\color{popdark}#1}}

% ============================================================
% En-tête / pied
% ============================================================
\pagestyle{fancy}
\fancyhf{}
\renewcommand{\headrulewidth}{0pt}
\renewcommand{\footrulewidth}{0pt}
\lhead{\raisebox{-0.15ex}{\poplogo\fontsize{13}{13}\selectfont\color{popink}OnBuch\color{poporange}+}}
\rhead{\poppill{\DOCMATIERE\ \textbullet\ \DOCNIVEAU\ \textbullet\ Leçon \DOCLECON}{white}{popink}}
\lfoot{\popxbold\fontsize{7.6}{7.6}\selectfont\color{popmuted}\MakeUppercase{Cours signé OnBuch+}\ \textbullet\ {\popsemi\DOCTITRE}}
\rfoot{\tikz[baseline=-0.6ex]\node[rounded corners=8.5pt,fill=popink,minimum height=17pt,minimum width=17pt,inner xsep=5pt,inner sep=0]{\popxbold\fontsize{8}{8}\selectfont\color{popcream}\thepage\,/\,\pageref*{LastPage}};}

% ============================================================
% Titres
% ============================================================
\newcommand{\chaptitre}[2]{%
  \begin{tcolorbox}[enhanced,colback=poporange,colframe=popink,boxrule=2.6pt,arc=11pt,
      drop shadow=popink,left=15pt,right=15pt,top=12pt,bottom=11pt,before skip=3pt,after skip=18pt]
    \poppill{#2}{popink}{popcream}\par\vspace{9pt}
    {\raggedright\hyphenpenalty=10000\exhyphenpenalty=10000\popdisplay\fontsize{22}{24}\selectfont\color{popcream}\MakeUppercase{#1}\par}
  \end{tcolorbox}
}
% Section numérotée : pastille orange avec le numéro + titre Archivo
\newcounter{popsec}
\newcommand{\coursec}[1]{%
  \stepcounter{popsec}\setcounter{popsub}{0}%
  \par\vspace{16pt}\noindent
  \begin{minipage}{\linewidth}
  \tikz[baseline=-0.7ex]\node[draw=popink,line width=1.6pt,fill=poporange,rounded corners=4pt,minimum size=22pt,inner sep=0]{\popdisplay\fontsize{12}{12}\selectfont\color{popcream}\thepopsec};%
  \hspace{8pt}{\popdisplay\fontsize{15}{17}\selectfont\color{popink}\MakeUppercase{#1}}\par
  \vspace{4pt}\noindent{\color{poporange}\rule{36pt}{3.6pt}}
  \end{minipage}\par\nopagebreak\vspace{8pt}
}
\newcounter{popsub}
\newcommand{\courssub}[1]{%
  \stepcounter{popsub}%
  \par\vspace{9pt}\noindent{\popxbold\fontsize{11.5}{13}\selectfont\color{popink}{\color{poporange}\thepopsec.\thepopsub}\hspace{6pt}#1}\par\nopagebreak\vspace{4pt}
}

% ============================================================
% Boîtes pédagogiques — même recette : fond blanc, bordure épaisse colorée,
% ombre dure encre, badge plein. Titre optionnel [#1].
% ============================================================
\tcbset{popbase/.style={breakable,enhanced,colback=white,boxrule=2.3pt,arc=9pt,
  shadow={3pt}{-3pt}{0pt}{popink},left=14pt,right=14pt,top=11pt,bottom=11pt,before skip=11pt,after skip=15pt,
  fontupper=\popsemi\fontsize{10.6}{14.5}\selectfont\color{popink},
  fontlower=\popsemi\fontsize{10.6}{14.5}\selectfont\color{popink}}}
% Coche dessinée (le glyphe ✓ n'existe pas dans les polices du document)
\newcommand{\popcheck}[1]{\tikz[baseline=-0.5ex]\draw[#1,line width=1.6pt,line cap=round,line join=round](-0.13,0.0)--(-0.04,-0.09)--(0.13,0.1);}
\providecommand{\coche}{\popcheck{popgreen}}
\providecommand{\resultat}[1]{\par\smallskip\noindent\fcolorbox{popgreen}{popgreenL}{\parbox{\dimexpr\linewidth-2\fboxsep-2\fboxrule\relax}{\textbf{Résultat.} #1}}\par\smallskip}
\newcommand{\popboxhead}[3]{\poppill{#1}{#2}{white}\ifx\relax#3\relax\else\hspace{6pt}{\popxbold\fontsize{10.6}{12}\selectfont\color{popink}#3}\fi\par\vspace{7pt}}

\newtcolorbox{aretenir}[1][]{popbase,colframe=popgreen,before upper={\popboxhead{À retenir}{popgreen}{#1}}}
% Texte en chinois (document de lecture, dialogue, extrait) : cadre
% d'étude. Un titre optionnel (source, auteur) s'affiche dans l'en-tête.
\newtcolorbox{textebox}[1][]{popbase,colframe=popblue,before upper={\popboxhead{文本 · Texte}{popblue}{#1}},after upper={}}
% Marques ✓ / ✗ (forme correcte / fautive) souvent utilisées par les modèles
\providecommand{\cross}{\textcolor{poppink}{\ensuremath{\times}}}
\providecommand{\xmark}{\cross}\providecommand{\crossmark}{\cross}\providecommand{\cmark}{\ensuremath{\checkmark}}
% Ligne de réponse / justification à compléter (inventée par les modèles)
\providecommand{\justification}[1]{\par\noindent\textit{Justification~:} \dotfill\par}
% \textipa (package tipa non chargé) : la police ne couvre pas l'API -> simple passage
\providecommand{\textipa}[1]{#1}
% Soulignement (ulem non chargé) : \uline, \ul
\providecommand{\uline}[1]{\underline{#1}}\providecommand{\ul}[1]{\underline{#1}}
% \glossaire{\item ...} inventée par les modèles : liste de glossaire
\providecommand{\glossaire}[1]{\begin{itemize}#1\end{itemize}}
% \alert (beamer) inventée par les modèles : texte mis en évidence
\providecommand{\alert}[1]{\textbf{\textcolor{poppink}{#1}}}
% variantes ASCII d'ordinaux écrites en macro
\providecommand{\iere}{\textsuperscript{re}}\providecommand{\ieme}{\textsuperscript{e}}\providecommand{\ere}{\textsuperscript{re}}\providecommand{\eme}{\textsuperscript{e}}\providecommand{\er}{\textsuperscript{er}}\providecommand{\ie}{\textsuperscript{e}}
% Ordinaux français écrits en macro par les modèles : 1\ière, 3\ième
\newcommand{\ière}{\textsuperscript{re}}\newcommand{\ième}{\textsuperscript{e}}
% \exemple (inventée par les modèles) : étiquette « Exemple : »
\providecommand{\exemple}{\textbf{Exemple~:}\ }
% \square utilisable aussi hors mode math (cases à cocher)
\AtBeginDocument{\let\squareMath\square\renewcommand{\square}{\ensuremath{\squareMath}}}
% Mot ou phrase en chinois à l'intérieur d'un paragraphe français
\newcommand{\zh}[1]{\ifmmode\text{#1}\else{#1}\fi}
\let\eng\zh \let\esp\zh \let\all\zh % alias tolérés
% flèches d'intonation inventées par les modèles
\providecommand{\textsearrow}{}\renewcommand{\textsearrow}{\ensuremath{\searrow}}\providecommand{\textnearrow}{}\renewcommand{\textnearrow}{\ensuremath{\nearrow}}
% \fr{..} : mot français (inventé par les modèles)
\providecommand{\fr}[1]{\textit{#1}}
% zhbox : encadré de texte chinois inventé par les modèles
\newenvironment{zhbox}{\begin{quote}}{\end{quote}}
% \courssubsub : titre de niveau 3 inventé par les modèles
\providecommand{\courssubsub}[1]{\par\medskip\noindent\textbf{#1}\par\smallskip}
% \zhbf : texte chinois en gras inventé par les modèles
\providecommand{\zhbf}[1]{\textbf{#1}}
% \vs : « versus » inventé par les modèles
\providecommand{\vs}{\textit{vs}\ }
% \blank : case à compléter inventée par les modèles
\providecommand{\blank}[1][]{\underline{\hspace{1.6em}}}
\newtcolorbox{exemplebox}[1][]{popbase,colframe=poporange,before upper={\popboxhead{Exemple}{poporange}{#1}}}
\newtcolorbox{definition}[1][]{popbase,colframe=popblue,before upper={\popboxhead{Définition}{popblue}{#1}}}
\newtcolorbox{propriete}[1][]{popbase,colframe=poppurple,before upper={\popboxhead{Propriété}{poppurple}{#1}}}
\newtcolorbox{methode}[1][]{popbase,colframe=popgold,before upper={\popboxhead{Méthode}{popgold}{#1}}}
\newtcolorbox{attention}[1][]{popbase,colframe=poppink,before upper={\popboxhead{Attention}{poppink}{#1}}}
\newtcolorbox{experience}[1][]{popbase,colframe=popblue,colback=popblueL,before upper={\popboxhead{Expérience}{popblue}{#1}}}
% « Le savais-tu ? » : carte violette pleine (contexte, culture, Cameroun)
\newtcolorbox{savaistu}[1][]{popbase,colback=poppurple,colframe=popink,
  fontupper=\popsemi\fontsize{10.4}{14.2}\selectfont\color{white},
  before upper={\poppill{Le savais-tu\,?}{white}{poppurple}\ifx\relax#1\relax\else\hspace{6pt}{\popxbold\color{white}#1}\fi\par\vspace{7pt}}}
% Exercice résolu : énoncé en haut, \tcblower puis la solution sur fond crème
\newcounter{popexo}\newcounter{popexr}
\newtcolorbox{exoresolu}[1][]{popbase,colframe=poporange,colbacklower=popcream,
  segmentation style={popink,line width=1.2pt,dashed},
  before upper={\stepcounter{popexr}\popboxhead{Exercice résolu \thepopexr}{poporange}{#1}},
  before lower={\poppill{Solution}{popink}{popcream}\par\vspace{6pt}}}
% Exercice (non résolu en place) — la correction est dans la partie Corrigés
\newtcolorbox{exercice}[1][]{popbase,colframe=popink,
  before upper={\stepcounter{popexo}\popboxhead{Exercice \thepopexo}{popink}{#1}}}
% Corrigé
\newtcolorbox{corrige}[1][]{popbase,colframe=popgreen,colback=popgreenL,
  before upper={\popboxhead{Corrigé}{popgreen}{#1}}}
% Objectifs / prérequis / bilan
\newtcolorbox{objectifs}{popbase,colframe=popink,colback=poporangeL,
  before upper={\popboxhead{Objectifs de la leçon}{popink}{}}}
\newtcolorbox{prerequis}{popbase,colframe=popink,colback=popblueL,
  before upper={\popboxhead{Prérequis}{popblue}{}}}
\newtcolorbox{bilan}{popbase,colframe=popink,colback=white,boxrule=2.8pt,
  before upper={\popboxhead{Fiche bilan}{poporange}{L'essentiel à connaître pour l'examen}}}
% Figure encadrée avec légende
\newtcolorbox{popfigure}[1][]{enhanced,breakable=false,colback=white,colframe=popline,boxrule=1.4pt,arc=8pt,
  left=10pt,right=10pt,top=10pt,bottom=8pt,before skip=10pt,after skip=12pt,halign=center,
  lower separated=false,halign lower=center,
  fontlower=\popsemi\fontsize{9}{11.5}\selectfont\color{popmuted},
  #1}
\newcommand{\legende}[1]{\par\vspace{6pt}{\popsemi\fontsize{9}{11.5}\selectfont\color{popmuted}#1}}

% ============================================================
% Signature OnBuch+
% ============================================================
\newcommand{\poplogoinline}[1]{{\poplogo\fontsize{#1}{#1}\selectfont\color{popink}OnBuch\color{poporange}+}}
\newcommand{\signatureonbuch}{%
  \begin{tcolorbox}[enhanced,colback=popink,colframe=popink,boxrule=2.2pt,arc=10pt,
      drop shadow=poporange,left=14pt,right=14pt,top=10pt,bottom=10pt,before skip=0pt,after skip=0pt]
    \tikz[baseline=-0.7ex]\node[circle,fill=popgreen,draw=popcream,line width=1.3pt,minimum size=22pt,inner sep=0]{\popcheck{white}};%
    \hspace{9pt}\begin{minipage}[c]{0.62\linewidth}
      {\popxbold\fontsize{12}{13}\selectfont\color{popcream}Signé\ \textbullet\ {\poplogo OnBuch\color{poporange}+}}\par\vspace{2pt}
      {\raggedright\popsemi\fontsize{8}{10.5}\selectfont\color{popline}\MakeUppercase{Équipe pédagogique OnBuch+}\\\MakeUppercase{Conforme au programme officiel MINESEC}\par}
    \end{minipage}\hfill
    {\poplogo\fontsize{22}{22}\selectfont\color{popcream}OnBuch\color{poporange}+}
  \end{tcolorbox}%
}

% ============================================================
% Page de garde
% #1 titre · #2 matière · #3 niveau · #4 module · #5 n° leçon · #6 durée ·
% #7 description / objectifs (texte ou itemize)
% ============================================================
\newcommand{\pagegarde}[7]{%
  \thispagestyle{empty}
  \newgeometry{a4paper,margin=1.7cm,top=1.6cm,bottom=1.4cm}%
  \begin{tikzpicture}[remember picture,overlay]
    \fill[poporange!18] ([xshift=-3.2cm,yshift=-3.6cm]current page.north east) circle (4.6cm);
    \fill[poppurple!14] ([xshift=2.4cm,yshift=7.6cm]current page.south west) circle (3.8cm);
  \end{tikzpicture}%
  {\poplogo\fontsize{30}{30}\selectfont\color{popink}OnBuch\color{poporange}+}\hfill
  \raisebox{0.6ex}{\poppill{Cours signé \textbullet\ Premium}{popink}{popcream}}\par
  \vspace{1pt}\popsquiggle{poppurple}\par
  \vspace{8pt}
  \begin{tcolorbox}[enhanced,colback=poporange,colframe=popink,boxrule=2.8pt,arc=13pt,
      drop shadow=popink,left=18pt,right=18pt,top=12pt,bottom=13pt,after skip=10pt]
    \poppill{#2 \textbullet\ #3}{popink}{popcream}\hspace{5pt}\poppill{#4}{white}{popink}\par\vspace{8pt}
    {\popdisplay\fontsize{14}{14}\selectfont\color{popink}LEÇON #5}\par\vspace{4pt}
    {\raggedright\hyphenpenalty=10000\exhyphenpenalty=10000\popdisplay\fontsize{25}{27}\selectfont\color{popcream}\MakeUppercase{#1}\par}
    \vspace{8pt}
    {\popxbold\fontsize{9.5}{9.5}\selectfont\color{popink}\MakeUppercase{Durée conseillée : #6}}
  \end{tcolorbox}
  \begin{tcolorbox}[enhanced,colback=white,colframe=popink,boxrule=2.2pt,arc=10pt,
      drop shadow=popink,left=16pt,right=16pt,top=10pt,bottom=10pt,after skip=8pt]
    \poppill{Dans cette leçon}{poppurple}{white}\par\vspace{5pt}
    \popsemi\fontsize{9.3}{12.6}\selectfont\color{popink}#7
  \end{tcolorbox}
  \noindent\begin{minipage}[t]{0.487\linewidth}
    \begin{tcolorbox}[enhanced,colback=popgreen,colframe=popink,boxrule=2.2pt,arc=10pt,
        drop shadow=popink,left=11pt,right=11pt,top=10pt,bottom=10pt,height=3.1cm,valign=center]
      \tcbox[colback=white,colframe=popink,boxrule=1.3pt,arc=5pt,boxsep=0pt,nobeforeafter,
        left=3pt,right=3pt,top=3pt,bottom=3pt,valign=center]{\qrcode[height=1.9cm]{https://play.google.com/store/apps/details?id=com.luvvix.onbuch&hl=fr}}%
      \hspace{8pt}\begin{minipage}[c]{0.5\linewidth}
        {\popdisplay\fontsize{11}{12.5}\selectfont\color{white}\MakeUppercase{Télécharge OnBuch}}\par\vspace{4pt}
        {\raggedright\popsemi\fontsize{8.4}{11}\selectfont\color{white}Gratuit \textbullet\ Google Play\\Tuteur IA, annales, concours, résultats\par}
      \end{minipage}
    \end{tcolorbox}
  \end{minipage}\hfill
  \begin{minipage}[t]{0.487\linewidth}
    \begin{tcolorbox}[enhanced,colback=poppurple,colframe=popink,boxrule=2.2pt,arc=10pt,
        drop shadow=popink,left=11pt,right=11pt,top=10pt,bottom=10pt,height=3.1cm,valign=center]
      \tikz[baseline=-0.7ex]\node[circle,fill=white,draw=popink,line width=1.3pt,minimum size=1.6cm,inner sep=0]{\popdisplay\fontsize{24}{24}\selectfont\color{poppurple}+};%
      \hspace{8pt}\begin{minipage}[c]{0.6\linewidth}
        {\popdisplay\fontsize{11}{12.5}\selectfont\color{white}\MakeUppercase{Passe à OnBuch+}}\par\vspace{4pt}
        {\raggedright\popsemi\fontsize{8.4}{11}\selectfont\color{white}Tous les cours signés, Tuteur IA illimité, fiches PDF — onglet OnBuch+ de l'app\par}
      \end{minipage}
    \end{tcolorbox}
  \end{minipage}
  \vspace{8pt}
  \popcontenu
  \vfill
  \signatureonbuch
  \restoregeometry
  \newpage
}

% Rangée « Ce cours contient » (page de garde)
\newcommand{\poptile}[3]{% #1 couleur #2 gros texte #3 légende
  \begin{tcolorbox}[enhanced,colback=#1,colframe=popink,boxrule=2pt,arc=9pt,shadow={2.5pt}{-2.5pt}{0pt}{popink},
      left=6pt,right=6pt,top=9pt,bottom=9pt,halign=center,valign=center,height=1.75cm,width=\linewidth,nobeforeafter]
    {\popdisplay\fontsize{12.5}{13}\selectfont\color{white}\MakeUppercase{#2}}\par\vspace{3pt}
    {\centering\popsemi\fontsize{7.8}{9.5}\selectfont\color{white}#3\par}
  \end{tcolorbox}}
\newcommand{\popcontenu}{%
  \par\noindent{\popxboldls\fontsize{7.5}{7.5}\selectfont\color{popmuted}CE COURS SIGNÉ CONTIENT}\par\vspace{7pt}
  \noindent\begin{minipage}[t]{0.235\linewidth}\poptile{popblue}{Cours}{complet, illustré, pas à pas}\end{minipage}\hfill
  \begin{minipage}[t]{0.235\linewidth}\poptile{poporange}{Méthodes}{et exercices résolus}\end{minipage}\hfill
  \begin{minipage}[t]{0.235\linewidth}\poptile{poppink}{Exercices}{type examen + corrigés}\end{minipage}\hfill
  \begin{minipage}[t]{0.235\linewidth}\poptile{popgreen}{Fiche bilan}{l'essentiel à retenir}\end{minipage}\par}

% Sommaire
\newtcolorbox{sommaire}{enhanced,colback=white,colframe=popink,boxrule=2.2pt,arc=10pt,
  drop shadow=popink,left=18pt,right=18pt,top=13pt,bottom=13pt,before skip=6pt,after skip=20pt,
  before upper={\poppill{Sommaire}{poppurple}{white}\par\vspace{9pt}\popsemi\fontsize{10.6}{14.5}\selectfont\color{popink}}}

% Signature de fin de cours — poussée tout en bas de la dernière page
\newcommand{\findecours}{%
  \par\vfill\nopagebreak
  \begin{center}\popsquiggle{poporange}\end{center}\vspace{4pt}
  \signatureonbuch
}
\AtBeginDocument{\def\llbracket{\mathopen{[\mkern-3.5mu[}}\def\rrbracket{\mathclose{]\mkern-3.5mu]}}} % intervalles d'entiers (glyphe ⟦ absent de la police)
% Glyphes absents de Fira Math : redéfinis à partir de glyphes disponibles
\AtBeginDocument{%
  \def\setminus{\mathbin{\backslash}}\let\smallsetminus\setminus
  \def\vdots{\mathord{\vbox{\baselineskip3.2pt\lineskiplimit0pt\kern4pt\hbox{.}\hbox{.}\hbox{.}}}}%
  \def\ddots{\mathinner{\mkern1mu\raise6.5pt\hbox{.}\mkern2mu\raise3.6pt\hbox{.}\mkern2mu\raise0.7pt\hbox{.}\mkern1mu}}%
  \def\triangle{\mathord{\scalebox{0.9}{$\symup{\Delta}$}}}%
  \def\nabla{\mathord{\raisebox{\depth}{\scalebox{1}[-1]{$\symup{\Delta}$}}}}%
  \def\checkmark{\popcheck{popdark}}%
  \def\star{\mathbin{\ast}}\def\bigstar{\mathord{\ast}}%
  \def\diamond{\mathbin{\scalebox{0.6}{$\Diamond$}}}%
  \def\bigcup{\mathop{\mathchoice{\raisebox{-0.3em}{\scalebox{1.7}{$\cup$}}}{\scalebox{1.25}{$\cup$}}{\cup}{\cup}}\displaylimits}%
  \def\bigcap{\mathop{\mathchoice{\raisebox{-0.3em}{\scalebox{1.7}{$\cap$}}}{\scalebox{1.25}{$\cap$}}{\cap}{\cap}}\displaylimits}%
  \def\lesseqqgtr{\mathrel{\lessgtr}}\def\bigoplus{\mathop{\oplus}\displaylimits}%
}
\providecommand{\ssi}{si et seulement si} % abréviation fréquente
\AtBeginDocument{\providecommand{\wideparen}{\overparen}} % arc de cercle

%% FIGFIX-BEGIN v5 — correctifs globaux des figures (docs/onbuch-generation/tools/figfix)
\usepackage{adjustbox}\usepackage{etoolbox}\tcbuselibrary{raster}
% toute figure TikZ/pgfplots plus large que la ligne est réduite à la largeur disponible
\BeforeBeginEnvironment{tikzpicture}{\begin{adjustbox}{max width=\linewidth}}
\AfterEndEnvironment{tikzpicture}{\end{adjustbox}}
% légendes pgfplots lisibles par-dessus les courbes
\pgfplotsset{every axis/.append style={legend style={fill=white,fill opacity=0.92,text opacity=1,draw=black!15,inner sep=3pt}}}
% étiquettes d'arêtes (syntaxe edge["..."]) avec fond blanc
\tikzset{every edge quotes/.append style={fill=white,inner sep=1.2pt}}
%% FIGFIX-END
\begin{document}

\pagegarde{Leçon 29 — Expression orale : corruption}{Chinois}{1ère A}{Module 3 : Citoyenneté et ouverture au monde}{ZHA29}{2 h}{Leçon d'expression orale consacrée au thème de la corruption : lexique essentiel, dialogues modèles avec pinyin, formules pour présenter un exposé, donner son avis et relancer, puis jeux de rôle et débat guidé en classe.
\vspace{4pt}
\begin{multicols}{2}\small
\begin{itemize}
\item À la fin de cette leçon, je sais utiliser le vocabulaire de base sur la corruption (腐败, 贪污, 受贿, 诚实, 法律...)
\item je sais présenter un exposé simple et structuré en chinois (introduction, arguments, conclusion)
\item je sais donner mon avis avec les formules 我觉得 / 我认为 / 在我看来
\item je sais exprimer l'accord et le désaccord poliment (我同意 / 我不同意 / 你说得对, 但是...)
\item je sais relancer un camarade avec des questions (你呢？你怎么看？为什么？)
\item je sais prononcer correctement les tons des mots clés du thème
\end{itemize}\end{multicols}}

\begin{sommaire}
\begin{multicols}{2}
\begin{enumerate}
\item Lexique du thème : la corruption
\item Dialogue modèle 1 : donner son avis
\item Dialogue modèle 2 : accord, désaccord, relance
\item Présenter un exposé simple
\item Travail des tons et de la prononciation
\item Jeux de rôle et débat guidé
\item Activité d'intégration
\item Méthodes et exercices résolus
\item Exercices
\item Corrigés des exercices
\item Fiche bilan
\end{enumerate}
\end{multicols}
\end{sommaire}

\begin{objectifs}
\begin{itemize}
\item À la fin de cette leçon, je sais utiliser le vocabulaire de base sur la corruption (腐败, 贪污, 受贿, 诚实, 法律...)
\item je sais présenter un exposé simple et structuré en chinois (introduction, arguments, conclusion)
\item je sais donner mon avis avec les formules 我觉得 / 我认为 / 在我看来
\item je sais exprimer l'accord et le désaccord poliment (我同意 / 我不同意 / 你说得对, 但是...)
\item je sais relancer un camarade avec des questions (你呢？你怎么看？为什么？)
\item je sais prononcer correctement les tons des mots clés du thème
\item je sais participer à un jeu de rôle et à un débat guidé sur la corruption au Cameroun et en Chine
\item je sais comparer brièvement la lutte contre la corruption au Cameroun et en Chine
\end{itemize}
\end{objectifs}

\begin{prerequis}
\begin{itemize}
\item Structures de base de la phrase chinoise (sujet + verbe + complément) vues en Seconde
\item Pronoms personnels et verbes courants (是, 有, 说, 看, 想)
\item Formules d'opinion simples (我觉得, 我喜欢 / 我不喜欢)
\item Lecture du pinyin et connaissance des quatre tons
\item Lexique de la vie scolaire et de la société (学校, 国家, 钱, 工作)
\end{itemize}
\end{prerequis}

\newpage

\coursec{Lexique du thème : la corruption}

\courssub{Situation d'accroche : un mot qu'on entend partout}

À Yaoundé, un élève raconte à ses camarades que son père a dû payer un « pot-de-vin » pour obtenir un document administratif. Tout le monde comprend, tout le monde a une opinion : certains sont choqués, d'autres disent que « c'est comme ça ». En Chine, la lutte contre la corruption, \zh{反腐败} (fǎnfǔbài), est un grand sujet de société dont on parle dans les médias, à l'école et dans les campagnes officielles.

\textbf{Problématique :} comment dire en chinois « Je pense que la corruption nuit au développement » ? Comment réagir, donner son avis, relancer un camarade dans un débat ? Pour cela, il faut d'abord posséder le \cle{vocabulaire} du thème. C'est l'objectif de cette première section : apprendre les mots clés, savoir les prononcer avec les bons tons et les utiliser dans des phrases simples.

\begin{textebox}[Phrase d'ouverture du débat]
腐败是一个大问题。\\
Fǔbài shì yí ge dà wèntí.\\
« La corruption est un grand problème. »
\end{textebox}

Observez cette phrase : \zh{腐败} est le sujet, \zh{是} est le verbe « être », \zh{一个} est le classificateur général \zh{个} précédé du chiffre un, et \zh{大问题} signifie « grand problème ». La structure est exactement la même qu'en français : Sujet + être + nom. C'est une excellente phrase d'ouverture pour un exposé oral. Notez aussi la règle phonétique : devant un 4\up{e} ton, \zh{一} (yī) se prononce \emph{yí} ; on écrit donc \zh{一个} mais on dit \emph{yí ge}.

\courssub{Le mot central : 腐败 fǔbài}

\begin{definition}[腐败 fǔbài]
\zh{腐败} (fǔbài) : nom et adjectif. Sens premier : « pourri, corrompu » (on dit d'un aliment pourri \zh{腐败的食物}). Sens social : « la corruption », c'est-à-dire l'abus d'un pouvoir pour obtenir de l'argent ou des avantages personnels. Le caractère \zh{腐} porte l'idée de « pourriture » et \zh{败} l'idée d'« échec, défaite » : littéralement, « pourri et vaincu ».
\end{definition}

De ce mot central dérivent deux expressions indispensables :

\begin{itemize}
\item \zh{反腐败} (fǎnfǔbài) : « lutte contre la corruption ». Le préfixe \zh{反} (fǎn) signifie « contre, anti- », comme dans \zh{反对} (fǎnduì, « s'opposer »).
\item \zh{腐败问题} (fǔbài wèntí) : « le problème de la corruption ».
\end{itemize}

\begin{exemplebox}[Premières phrases avec 腐败]
\zh{我们反对腐败。} \emph{Wǒmen fǎnduì fǔbài.} « Nous luttons contre la corruption. »\\
\zh{反腐败很重要。} \emph{Fǎnfǔbài hěn zhòngyào.} « La lutte contre la corruption est très importante. »\\
\zh{腐败影响国家的发展。} \emph{Fǔbài yǐngxiǎng guójiā de fāzhǎn.} « La corruption nuit au développement du pays. »
\end{exemplebox}

\courssub{Les actes : détourner, recevoir, donner}

Pour parler de la corruption, il faut distinguer trois verbes précis. En français, on utilise souvent le mot vague « corrompre » ; le chinois, lui, nomme chaque geste exactement.

\begin{definition}[贿赂 huìlù]
\zh{贿赂} (huìlù) : nom. Argent ou cadeau donné à quelqu'un pour obtenir un avantage illégal : le « pot-de-vin ». Attention : \zh{贿赂} désigne l'objet (l'argent, le cadeau), pas l'action.
\end{definition}

\begin{center}
\begin{tabularx}{\linewidth}{|X|X|X|X|}
\hline
\rowcolor{popblueL} Caractères & Pinyin & Nature & Traduction \\
\hline
贪污 & tānwū & verbe & détourner de l'argent (public) \\
\hline
受贿 & shòuhuì & verbe & accepter un pot-de-vin \\
\hline
行贿 & xínghuì & verbe & donner un pot-de-vin \\
\hline
贿赂 & huìlù & nom & pot-de-vin \\
\hline
钱 & qián & nom & argent \\
\hline
礼物 & lǐwù & nom & cadeau \\
\hline
\end{tabularx}
\end{center}

\begin{attention}[Ne pas confondre 受贿 et 行贿]
Erreur classique des francophones : inverser \zh{受贿} (shòuhuì) et \zh{行贿} (xínghuì). Retenez le sens des caractères : \zh{受} (shòu) = \cle{recevoir}, donc \zh{受贿} = « recevoir un pot-de-vin » (c'est celui qui accepte) ; \zh{行} (xíng) = \cle{faire, effectuer}, donc \zh{行贿} = « donner un pot-de-vin » (c'est celui qui paie). Astuce mnémotechnique : \zh{受} évoque une main qui reçoit ; \zh{行} évoque l'action, le geste de donner. Les deux actes sont condamnés par la loi : \zh{受贿和行贿都是违法的。} \emph{Shòuhuì hé xínghuì dōu shì wéifǎ de.} « Recevoir et donner des pots-de-vin sont tous les deux illégaux. »
\end{attention}

\begin{exemplebox}[Les trois actes en contexte]
\zh{他贪污了很多钱。} \emph{Tā tānwū le hěn duō qián.} « Il a détourné beaucoup d'argent. »\\
\zh{那个官员受贿了。} \emph{Nà ge guānyuán shòuhuì le.} « Ce fonctionnaire a accepté un pot-de-vin. »\\
\zh{有人行贿，有人受贿。} \emph{Yǒu rén xínghuì, yǒu rén shòuhuì.} « Certains donnent des pots-de-vin, d'autres en acceptent. »
\end{exemplebox}

Remarquez l'emploi de \zh{了} après le verbe pour marquer une action accomplie, et le classificateur \zh{个} dans \zh{那个官员} (« ce fonctionnaire »). La structure \zh{有人……，有人……} (« certains..., d'autres... ») est très utile pour décrire une situation dans un exposé.

\courssub{Les valeurs : honnêteté, intégrité, justice}

Un débat sur la corruption ne se limite pas aux actes illégaux : il faut aussi pouvoir parler des valeurs positives. Trois adjectifs sont essentiels.

\begin{center}
\begin{tabularx}{\linewidth}{|X|X|X|X|}
\hline
\rowcolor{popblueL} Caractères & Pinyin & Nature & Traduction \\
\hline
诚实 & chéngshí & adjectif & honnête, sincère \\
\hline
正直 & zhèngzhí & adjectif & intègre, droit \\
\hline
公平 & gōngpíng & adjectif & juste, équitable \\
\hline
\end{tabularx}
\end{center}

\begin{propriete}[Adjectif chinois = adjectif-verbe]
En chinois, l'adjectif se comporte comme un verbe : on ne dit pas « il \emph{est} honnête » avec \zh{是}, mais \zh{他很诚实。} \emph{Tā hěn chéngshí.} « Il est honnête. » Le mot \zh{很} (hěn, « très ») sert ici de simple lien et ne marque pas forcément l'intensité. Ne dites jamais \zh{他是诚实} : c'est une erreur typique des francophones qui calquent « il est honnête ». En revanche, \zh{是} est obligatoire quand l'adjectif qualifie un nom grâce à \zh{的} : \zh{她是一个正直的人。} « C'est une personne intègre. »
\end{propriete}

\begin{exemplebox}[Les valeurs en phrases]
\zh{他很诚实。} \emph{Tā hěn chéngshí.} « Il est honnête. »\\
\zh{她是一个正直的人。} \emph{Tā shì yí ge zhèngzhí de rén.} « C'est une personne intègre. » (Ici \zh{是} est correct, car \zh{正直} qualifie le nom \zh{人} grâce à \zh{的}.)\\
\zh{我们要公平的社会。} \emph{Wǒmen yào gōngpíng de shèhuì.} « Nous voulons une société juste. »
\end{exemplebox}

\courssub{Les institutions et les conséquences}

Pour compléter le lexique du débat, deux derniers groupes de mots : les institutions qui luttent contre la corruption, et les mots qui expriment ses conséquences.

\begin{center}
\begin{tabularx}{\linewidth}{|X|X|X|X|}
\hline
\rowcolor{popblueL} Caractères & Pinyin & Nature & Traduction \\
\hline
法律 & fǎlǜ & nom & loi, droit \\
\hline
警察 & jǐngchá & nom & police, policier \\
\hline
政府 & zhèngfǔ & nom & gouvernement \\
\hline
问题 & wèntí & nom & problème, question \\
\hline
影响 & yǐngxiǎng & nom / verbe & influence, impact / influencer, nuire à \\
\hline
发展 & fāzhǎn & nom / verbe & développement / se développer \\
\hline
\end{tabularx}
\end{center}

\begin{exemplebox}[Institutions et conséquences]
\zh{政府应该反腐败。} \emph{Zhèngfǔ yīnggāi fǎnfǔbài.} « Le gouvernement doit lutter contre la corruption. »\\
\zh{警察抓贪污的人。} \emph{Jǐngchá zhuā tānwū de rén.} « La police arrête les personnes qui détournent de l'argent. »\\
\zh{腐败影响国家的发展。} \emph{Fǔbài yǐngxiǎng guójiā de fāzhǎn.} « La corruption nuit au développement du pays. »\\
\zh{这是一个大问题。} \emph{Zhè shì yí ge dà wèntí.} « C'est un grand problème. »
\end{exemplebox}

Notez le verbe modal \zh{应该} (yīnggāi, « devoir »), très utile pour donner son avis : Sujet + \zh{应该} + verbe. Notez aussi le verbe \zh{抓} (zhuā, « attraper, arrêter »), imagé et fréquent à l'oral.

\courssub{Carte mentale du lexique}

\begin{popfigure}
\begin{center}\tcbox[colback=popblue,colframe=popblue,coltext=white,arc=9pt,boxsep=4pt,left=6pt,right=6pt]{\bfseries\small 腐败 fǔbài}\end{center}
\vspace{-2pt}
\begin{tcbraster}[raster columns=2,raster equal height=rows,raster column skip=6pt,raster row skip=6pt,enhanced,blankest]
\begin{tcolorbox}[enhanced,colback=white,colframe=poporange,boxrule=0.9pt,arc=3pt,title={actes 行为},fonttitle=\bfseries\scriptsize,coltitle=white,colbacktitle=poporange,left=4pt,right=4pt,top=2pt,bottom=2pt,boxsep=1pt,toptitle=1.5pt,bottomtitle=1.5pt,halign=left]
\begin{itemize}[nosep,leftmargin=*,itemsep=1pt,font=\scriptsize]
\item 贪污 tānwū
\item 受贿 shòuhuì
\item 行贿 xínghuì
\end{itemize}
\end{tcolorbox}
\begin{tcolorbox}[enhanced,colback=white,colframe=popgreen,boxrule=0.9pt,arc=3pt,title={valeurs 价值},fonttitle=\bfseries\scriptsize,coltitle=white,colbacktitle=popgreen,left=4pt,right=4pt,top=2pt,bottom=2pt,boxsep=1pt,toptitle=1.5pt,bottomtitle=1.5pt,halign=left]
\begin{itemize}[nosep,leftmargin=*,itemsep=1pt,font=\scriptsize]
\item 诚实 chéngshí
\item 正直 zhèngzhí
\item 公平 gōngpíng
\end{itemize}
\end{tcolorbox}
\begin{tcolorbox}[enhanced,colback=white,colframe=poppurple,boxrule=0.9pt,arc=3pt,title={institutions 机构},fonttitle=\bfseries\scriptsize,coltitle=white,colbacktitle=poppurple,left=4pt,right=4pt,top=2pt,bottom=2pt,boxsep=1pt,toptitle=1.5pt,bottomtitle=1.5pt,halign=left]
\begin{itemize}[nosep,leftmargin=*,itemsep=1pt,font=\scriptsize]
\item 法律 fǎlǜ
\item 警察 jǐngchá
\item 政府 zhèngfǔ
\end{itemize}
\end{tcolorbox}
\begin{tcolorbox}[enhanced,colback=white,colframe=poppink,boxrule=0.9pt,arc=3pt,title={conséquences 后果},fonttitle=\bfseries\scriptsize,coltitle=white,colbacktitle=poppink,left=4pt,right=4pt,top=2pt,bottom=2pt,boxsep=1pt,toptitle=1.5pt,bottomtitle=1.5pt,halign=left]
\begin{itemize}[nosep,leftmargin=*,itemsep=1pt,font=\scriptsize]
\item 问题 wèntí
\item 影响 yǐngxiǎng
\item 发展 fāzhǎn
\end{itemize}
\end{tcolorbox}
\end{tcbraster}

\legende{Carte mentale du lexique : autour du mot central 腐败, quatre familles de mots à mémoriser ensemble.}
\end{popfigure}

\begin{methode}[Mémoriser le lexique par familles]
\begin{enumerate}
\item Apprenez d'abord le mot central \zh{腐败} et sa famille : \zh{反腐败}, \zh{腐败问题}.
\item Associez chaque famille à une couleur de la carte mentale et récitez les trois mots de chaque branche à voix haute, en marquant bien les tons.
\item Construisez une phrase simple par famille : un acte (\zh{他受贿了。}), une valeur (\zh{她很正直。}), une institution (\zh{政府反腐败。}), une conséquence (\zh{腐败影响发展。}).
\item Vérifiez-vous en binôme : un élève dit le français, l'autre répond en chinois avec le pinyin et les tons.
\end{enumerate}
\end{methode}

\courssub{Prononciation guidée : les tons des mots clés}

La justesse des tons est indispensable à l'oral : un mot mal toné peut devenir incompréhensible. Voici le profil tonal des mots clés, à travailler d'abord en répétition chorale (toute la classe répète après le professeur), puis individuellement.

\begin{center}
\begin{tabularx}{\linewidth}{|X|X|X|}
\hline
\rowcolor{popblueL} Mot & Tons & Conseil de prononciation \\
\hline
腐败 fǔbài & 3 + 4 & Le 3\up{e} ton descend bas ; le 4\up{e} ton tombe net, comme un ordre. \\
\hline
贪污 tānwū & 1 + 1 & Deux tons hauts et plats : gardez la voix haute et stable. \\
\hline
受贿 shòuhuì & 4 + 4 & Deux tons descendants : frappez chaque syllabe fermement. \\
\hline
行贿 xínghuì & 2 + 4 & Le 2\up{e} ton monte comme une question, puis chute sur \zh{贿}. \\
\hline
诚实 chéngshí & 2 + 2 & Deux tons montants : la voix s'élève deux fois. \\
\hline
公平 gōngpíng & 1 + 2 & Ton haut plat, puis montée. \\
\hline
法律 fǎlǜ & 3 + 4 & Attention à \zh{律} : voyelle \zh{ü}, lèvres arrondies et avancées. \\
\hline
影响 yǐngxiǎng & 3 + 3 & Deux 3\up{e} tons : le premier se prononce comme un 2\up{e} ton (\emph{yíngxiǎng}) par règle de sandhi. \\
\hline
发展 fāzhǎn & 1 + 3 & Ton haut plat, puis ton creux. \\
\hline
\end{tabularx}
\end{center}

\begin{attention}[Le sandhi du troisième ton]
Quand deux 3\up{e} tons se suivent, le premier devient un 2\up{e} ton à l'oral : \zh{影响} s'écrit yǐngxiǎng mais se prononce \emph{yíngxiǎng}. De même, \zh{很好} (hěn hǎo) se prononce \emph{hén hǎo}. On écrit toujours le ton d'origine dans le pinyin, mais on change la prononciation. À l'oral, un 3\up{e} ton suivi d'un autre ton se réduit souvent à sa partie basse (demi-3\up{e} ton), comme dans \zh{腐败} fǔbài : ne remontez pas la voix sur \zh{腐}.
\end{attention}

\begin{exoresolu}[Associer le mot chinois à sa traduction]
Énoncé : reliez chaque mot chinois à sa traduction française. 1. 行贿 \quad 2. 受贿 \quad 3. 贪污 \quad 4. 公平 \quad 5. 法律 \quad a. loi \quad b. donner un pot-de-vin \quad c. détourner de l'argent \quad d. juste, équitable \quad e. accepter un pot-de-vin
\tcblower
Solution détaillée : 1-b : \zh{行} (xíng) = faire, donner, donc \zh{行贿} = donner un pot-de-vin. 2-e : \zh{受} (shòu) = recevoir, donc \zh{受贿} = accepter un pot-de-vin. 3-c : \zh{贪污} (tānwū) = détourner de l'argent public. 4-d : \zh{公平} (gōngpíng) = juste, équitable ; \zh{公} = public, commun, \zh{平} = plat, égal. 5-a : \zh{法律} (fǎlǜ) = loi ; \zh{法} = loi, \zh{律} = règle. Vérification par des phrases : \zh{行贿和受贿都不好。} \emph{Xínghuì hé shòuhuì dōu bù hǎo.} « Donner et accepter des pots-de-vin, ce n'est pas bien. » \zh{法律面前，人人平等。} \emph{Fǎlǜ miànqián, rénrén píngděng.} « Devant la loi, tous sont égaux. »
\end{exoresolu}

\begin{exercice}[À vous de parler]
En binôme, prononcez à tour de rôle les six phrases des boîtes d'exemples de cette section, en veillant aux tons. Puis chacun construit une phrase originale avec deux mots du lexique (par exemple \zh{诚实} + \zh{发展}). Le camarade vérifie les tons et la place des mots.
\end{exercice}

\begin{corrige}[Exercice « À vous de parler »]
Exemples de phrases originales attendues : \zh{诚实的人帮助国家的发展。} \emph{Chéngshí de rén bāngzhù guójiā de fāzhǎn.} « Les personnes honnêtes aident le développement du pays. » \zh{公平的社会发展得很快。} \emph{Gōngpíng de shèhuì fāzhǎn de hěn kuài.} « Une société juste se développe vite. » Critères de validation : tons corrects, adjectif précédé de \zh{很} ou suivi de \zh{的} + nom, ordre Sujet + Verbe + Complément respecté.
\end{corrige}

\begin{savaistu}[« Tigres et mouches » : la campagne anticorruption en Chine]
Depuis les années 2010, la Chine mène une vaste campagne \zh{反腐败} (fǎnfǔbài) résumée par une formule célèbre : \zh{老虎苍蝇一起打} (lǎohǔ cāngying yìqǐ dǎ), littéralement « frapper ensemble les tigres et les mouches ». Les « tigres » sont les grands dirigeants corrompus, les « mouches » les petits fonctionnaires malhonnêtes : la campagne vise tous les niveaux. Au Cameroun aussi, la lutte contre la corruption fait l'objet d'institutions et de campagnes publiques : c'est un excellent sujet de comparaison pour votre exposé oral. Vous pouvez dire : \zh{中国和喀麦隆都反对腐败。} \emph{Zhōngguó hé Kāmàilóng dōu fǎnduì fǔbài.} « La Chine et le Cameroun luttent tous les deux contre la corruption. »
\end{savaistu}

\begin{aretenir}[L'essentiel du lexique]
\begin{itemize}
\item Mot central : \zh{腐败} fǔbài (corruption) ; action de lutte : \zh{反腐败} fǎnfǔbài.
\item Trois actes à ne pas confondre : \zh{贪污} (détourner), \zh{受贿} (recevoir un pot-de-vin), \zh{行贿} (donner un pot-de-vin).
\item Trois valeurs : \zh{诚实} (honnête), \zh{正直} (intègre), \zh{公平} (juste) — adjectifs-verbes précédés de \zh{很}, jamais de \zh{是} devant l'adjectif seul.
\item Institutions : \zh{法律}, \zh{警察}, \zh{政府} ; conséquences : \zh{问题}, \zh{影响}, \zh{发展}.
\item Phrase-clé du débat : \zh{腐败影响国家的发展。} « La corruption nuit au développement du pays. »
\end{itemize}
\end{aretenir}

\coursec{Dialogue modèle 1 : donner son avis}

Dans la section précédente, nous avons construit le \cle{lexique de la corruption} : \zh{腐败} (fǔbài, « corruption »), \zh{严重} (yánzhòng, « grave »), \zh{影响} (yǐngxiǎng, « influencer, nuire à »), \zh{发展} (fāzhǎn, « développement »), \zh{诚实} (chéngshí, « honnêteté »). Posséder des mots ne suffit pas : pour participer à un débat ou présenter un exposé, il faut savoir \textbf{donner son avis}. C'est exactement l'objectif de cette section : apprendre à dire « je pense que… » en chinois, avec plusieurs formules de registres différents, et à prononcer correctement les tons de ces formules.

\courssub{Observation : un premier échange}

Lisons d'abord ce court dialogue entre deux élèves d'une classe de Première à Yaoundé, qui préparent un exposé sur la vie citoyenne.

\begin{textebox}[Dialogue 1 — Donner son avis]
A : 你怎么看？\\
\emph{Nǐ zěnme kàn ?}\\
« Qu'en penses-tu ? »\\[4pt]
B : 我觉得腐败是一个严重的问题。\\
\emph{Wǒ juéde fǔbài shì yí ge yánzhòng de wèntí.}\\
« Je pense que la corruption est un problème grave. »\\[4pt]
A : 你觉得腐败是一个严重的问题吗？\\
\emph{Nǐ juéde fǔbài shì yí ge yánzhòng de wèntí ma ?}\\
« Penses-tu que la corruption est un problème grave ? »\\[4pt]
B : 我觉得是。腐败影响国家的发展。\\
\emph{Wǒ juéde shì. Fǔbài yǐngxiǎng guójiā de fāzhǎn.}\\
« Je pense que oui. Elle nuit au développement du pays. »
\end{textebox}

\textbf{Observons.} Les deux questions sont différentes : \zh{你怎么看？} est une question ouverte (« quel est ton avis ? »), tandis que \zh{你觉得…吗？} est une question fermée qui attend « oui » ou « non ». Dans les deux cas, la réponse commence par la même formule : \zh{我觉得} « je pense ». Remarquons aussi la réponse courte \zh{我觉得是} : le chinois n'a pas de mot unique pour « oui » ; on reprend le verbe ou l'adjectif de la question, ici \zh{是} « être ».

\begin{center}
\begin{tabularx}{\linewidth}{|X|X|X|X|}
\hline
\rowcolor{popblueL} Caractères & Pinyin & Nature & Traduction \\
\hline
觉得 & juéde & verbe & penser, trouver, avoir le sentiment que \\
\hline
认为 & rènwéi & verbe & considérer, estimer (formel) \\
\hline
在我看来 & zài wǒ kànlái & locution & à mon avis, selon moi \\
\hline
严重 & yánzhòng & adjectif & grave, sérieux \\
\hline
问题 & wèntí & nom & problème, question \\
\hline
影响 & yǐngxiǎng & verbe / nom & influencer, nuire à ; influence \\
\hline
国家 & guójiā & nom & pays, État \\
\hline
发展 & fāzhǎn & verbe / nom & (se) développer ; développement \\
\hline
诚实 & chéngshí & adjectif / nom & honnête ; honnêteté \\
\hline
重要 & zhòngyào & adjectif & important \\
\hline
\end{tabularx}
\end{center}

\courssub{La règle : les verbes d'opinion 觉得 et 认为}

\begin{propriete}[Exprimer une opinion : 觉得 et 认为]
Les verbes \cle{觉得} (juéde) et \cle{认为} (rènwéi) introduisent une opinion personnelle, comme « je pense que » en français. Ils sont suivis d'une \textbf{proposition complète} (sujet + verbe), jamais d'un nom seul.\\
Structure : \textbf{Sujet + 觉得 / 认为 + proposition (sujet + verbe/adjectif)}\\
\zh{我觉得腐败很严重。} \emph{Wǒ juéde fǔbài hěn yánzhòng.} « Je pense que la corruption est très grave. »\\
Différence de registre : \zh{觉得} est courant et oral ; \zh{认为} est plus \textbf{formel}, idéal pour un exposé ou une rédaction. La locution \zh{在我看来} (zài wǒ kànlái, « à mon avis ») se place en tête de phrase et insiste sur le point de vue personnel.
\end{propriete}

Comparaison avec le français : le français exige « je pense \emph{que}… » avec une conjonction ; le chinois juxtapose directement les deux propositions, sans mot de liaison : \zh{我觉得} + \zh{腐败很严重}. Il ne faut donc \textbf{rien} ajouter entre les deux. Autre différence : en chinois, on ne conjugue rien — le verbe d'opinion reste identique quelle que soit la personne (\zh{我觉得}, \zh{他觉得} tā juéde « il pense », \zh{我们认为} wǒmen rènwéi « nous estimons »).

\begin{exemplebox}[Exemples traduits]
\zh{我认为诚实很重要。} \emph{Wǒ rènwéi chéngshí hěn zhòngyào.} « Je considère que l'honnêteté est très importante. »\\
\zh{我觉得腐败很不好。} \emph{Wǒ juéde fǔbài hěn bù hǎo.} « Je pense que la corruption est très mauvaise. »\\
\zh{他认为腐败影响国家的发展。} \emph{Tā rènwéi fǔbài yǐngxiǎng guójiā de fāzhǎn.} « Il estime que la corruption nuit au développement du pays. »\\
\zh{在我看来，这是一个很大的问题。} \emph{Zài wǒ kànlái, zhè shì yí ge hěn dà de wèntí.} « À mon avis, c'est un très grand problème. »\\
\zh{我们都觉得诚实的人值得尊重。} \emph{Wǒmen dōu juéde chéngshí de rén zhídé zūnzhòng.} « Nous pensons tous que les gens honnêtes méritent le respect. »
\end{exemplebox}

\begin{popfigure}
\begin{tikzpicture}[node distance=0.9cm]
\node[draw=popgreen, fill=popgreenL, rounded corners, thick, align=center, minimum width=3.4cm] (juede) {\zh{我觉得}\\ \emph{wǒ juéde}\\ je pense (oral)};
\node[draw=popblue, fill=popblueL, rounded corners, thick, align=center, minimum width=3.4cm, right=1.1cm of juede] (renwei) {\zh{我认为}\\ \emph{wǒ rènwéi}\\ je considère (formel)};
\node[draw=poppurple, fill=poppurpleL, rounded corners, thick, align=center, minimum width=3.4cm, right=1.1cm of renwei] (kanlai) {\zh{在我看来}\\ \emph{zài wǒ kànlái}\\ à mon avis (point de vue)};
\draw[-{Stealth}, very thick, popdark] (juede) -- (renwei);
\draw[-{Stealth}, very thick, popdark] (renwei) -- (kanlai);
\node[below=0.25cm of renwei, text=popmuted, align=center] {\small registre de plus en plus soutenu};
\end{tikzpicture}
\legende{L'échelle d'opinion : trois formules pour dire « je pense », du plus courant au plus formel.}
\end{popfigure}

\courssub{La méthode : construire sa phrase d'opinion en trois étapes}

\begin{methode}[Donner son avis en débat]
\begin{enumerate}
\item \textbf{Choisir la formule d'opinion} adaptée au contexte : \zh{我觉得} à l'oral entre camarades, \zh{我认为} ou \zh{在我看来} pour un exposé devant la classe.
\item \textbf{Annoncer le sujet} : \zh{腐败} (la corruption), \zh{诚实} (l'honnêteté), \zh{这个问题} (ce problème).
\item \textbf{Compléter} avec \zh{是} + groupe nominal (\zh{是一个严重的问题} « est un problème grave ») ou \zh{很} + adjectif (\zh{很严重} « est très grave »), ou une phrase complète avec verbe d'action (\zh{腐败影响国家的发展}).
\end{enumerate}
Schéma récapitulatif : \textbf{我觉得 / 我认为 + 腐败 + 是 + un problème grave} ou \textbf{+ 很 + adjectif} ou \textbf{+ verbe + complément}.
\end{methode}

\begin{center}
\begin{tabularx}{\linewidth}{|X|X|X|}
\hline
\rowcolor{popblueL} Structure & Exemple (caractères + pinyin) & Traduction \\
\hline
我觉得 + 是 + GN & 我觉得腐败是一个严重的问题。\newline \emph{Wǒ juéde fǔbài shì yí ge yánzhòng de wèntí.} & Je pense que la corruption est un problème grave. \\
\hline
我觉得 + 很 + Adj. & 我觉得腐败很严重。\newline \emph{Wǒ juéde fǔbài hěn yánzhòng.} & Je pense que la corruption est très grave. \\
\hline
我认为 + phrase & 我认为腐败很不好。\newline \emph{Wǒ rènwéi fǔbài hěn bù hǎo.} & Je considère que la corruption est très mauvaise. \\
\hline
在我看来，+ phrase & 在我看来，这是一个很大的问题。\newline \emph{Zài wǒ kànlái, zhè shì yí ge hěn dà de wèntí.} & À mon avis, c'est un très grand problème. \\
\hline
\end{tabularx}
\end{center}

\courssub{Travail des tons : prononcer les formules d'opinion}

Un débat réussi commence par une prononciation claire. Voici les quatre mots-clés de cette section, avec leurs tons à soigner :

\begin{itemize}
\item \zh{觉得} \emph{juéde} : 2\up{e} ton montant sur \emph{jué}, puis \textbf{ton neutre} léger et court sur \emph{de}. Erreur typique des francophones : appuyer sur la deuxième syllabe. Dites « JUÉ-de », le \emph{de} presque soufflé.
\item \zh{认为} \emph{rènwéi} : 4\up{e} ton descendant franc sur \emph{rèn}, puis 2\up{e} ton montant sur \emph{wéi}. Le contraste « tombe puis remonte » doit être net.
\item \zh{严重} \emph{yánzhòng} : 2\up{e} ton puis 4\up{e} ton. Attention à ne pas inverser : \emph{yán} monte, \emph{zhòng} descend.
\item \zh{影响} \emph{yǐngxiǎng} : écrit avec deux 3\up{e} tons (3+3), mais par \cle{sandhi tonal}, le premier 3\up{e} ton devient un 2\up{e} ton à l'oral : on prononce donc \emph{yíngxiǎng} (2+3). C'est la même règle que \zh{你好} \emph{nǐhǎo} prononcé \emph{níhǎo}.
\end{itemize}

\begin{attention}[Pièges fréquents des francophones]
\begin{itemize}
\item Dans \zh{你怎么看？}, le verbe \zh{看} se prononce \emph{kàn} (4\up{e} ton), jamais \emph{kān}. Ici il ne signifie pas « regarder » mais « envisager, juger ».
\item Pour demander un avis \textbf{argumenté} en débat, on dit \zh{你怎么看？} « qu'en penses-tu ? ». N'utilisez pas \zh{你怎么想？} (\emph{nǐ zěnme xiǎng}), qui demanderait plutôt « qu'as-tu en tête, qu'imagines-tu ? » : \zh{怎么看} est la formule consacrée du débat d'opinion.
\item Ne traduisez pas « que » : après \zh{觉得} ou \zh{认为}, on n'ajoute aucun mot de liaison. \zh{我觉得是} suffit pour « je pense que oui » — inutile de chercher un équivalent de « oui ».
\item \zh{认为} ne s'emploie pas pour les sensations : pour « je trouve le ndolé bon », dites \zh{我觉得} (perception), pas \zh{我认为} (jugement réfléchi).
\end{itemize}
\end{attention}

\begin{exoresolu}[À vous de donner votre avis]
Énoncé : votre camarade vous demande \zh{你怎么看？} au sujet de l'honnêteté à l'école. Répondez en une phrase complète avec \zh{觉得}, puis reformulez avec \zh{认为} pour un exposé.
\tcblower
Solution détaillée :\\
Étape 1 — formule orale + sujet + \zh{很} + adjectif :\\
\zh{我觉得诚实很重要。} \emph{Wǒ juéde chéngshí hěn zhòngyào.} « Je pense que l'honnêteté est très importante. »\\
Étape 2 — reformulation formelle pour l'exposé :\\
\zh{我认为诚实是一个很重要的问题。} \emph{Wǒ rènwéi chéngshí shì yí ge hěn zhòngyào de wèntí.} « Je considère que l'honnêteté est une question très importante. »\\
Vérification : aucune conjonction entre les deux propositions ; \zh{很} précède l'adjectif ; les tons de \emph{juéde} (2 + neutre) et de \emph{rènwéi} (4+2) sont respectés.
\end{exoresolu}

\begin{exercice}[Entraînement rapide]
Traduisez en chinois, en choisissant la formule adaptée au registre indiqué :
\begin{enumerate}
\item « Je pense que la corruption est un problème grave. » (oral, entre amis)
\item « À mon avis, l'honnêteté est très importante. » (exposé)
\item « Il estime que la corruption nuit au développement du pays. » (exposé)
\end{enumerate}
\end{exercice}

\begin{corrige}[Exercice « Entraînement rapide »]
\begin{enumerate}
\item \zh{我觉得腐败是一个严重的问题。} \emph{Wǒ juéde fǔbài shì yí ge yánzhòng de wèntí.}
\item \zh{在我看来，诚实很重要。} \emph{Zài wǒ kànlái, chéngshí hěn zhòngyào.}
\item \zh{他认为腐败影响国家的发展。} \emph{Tā rènwéi fǔbài yǐngxiǎng guójiā de fāzhǎn.}
\end{enumerate}
\end{corrige}

\begin{aretenir}[L'essentiel de la section]
Pour donner son avis : \zh{我觉得} (oral) et \zh{认为} (formel) + proposition complète, sans conjonction ; \zh{在我看来} en tête de phrase pour marquer son point de vue. Question ouverte du débat : \zh{你怎么看？} (\emph{kàn}, 4\up{e} ton). Réponse courte affirmative : \zh{我觉得是}. Tons à maîtriser : \emph{juéde} (2 + neutre), \emph{rènwéi} (4+2), \emph{yánzhòng} (2+4), \emph{yǐngxiǎng} prononcé (2+3) par sandhi.
\end{aretenir}

\coursec{Dialogue modèle 2 : accord, désaccord, relance}

Dans la section précédente, nous avons appris à \cle{donner notre avis} avec des formules comme \zh{我认为……} (Wǒ rènwéi..., « je pense que... ») et \zh{我觉得……} (Wǒ juéde..., « je trouve que... »). Mais un débat ne s'arrête pas à une opinion : il faut aussi \cle{réagir} à ce que dit l'autre, dire si l'on est d'accord ou non, \cle{relancer} la conversation, puis \cle{conclure}. C'est exactement ce que nous allons construire ici, pas à pas, à partir d'un dialogue authentique entre deux élèves d'un lycée de Yaoundé.

\courssub{Observation : un mini-débat sur l'honnêteté}

\begin{textebox}[Dialogue modèle 2 — 诚实很重要]
A : 我们应该诚实。\\
\emph{Wǒmen yīnggāi chéngshí.}\\
« Nous devons être honnêtes. »\\
B : 我同意！诚实很重要。\\
\emph{Wǒ tóngyì ! Chéngshí hěn zhòngyào.}\\
« Je suis d'accord ! L'honnêteté est très importante. »\\
A : 但是，有时候诚实很难。\\
\emph{Dànshì, yǒu shíhou chéngshí hěn nán.}\\
« Mais parfois, être honnête est difficile. »\\
B : 你说得有道理，但是诚实总是对的。\\
\emph{Nǐ shuō de yǒu dàolǐ, dànshì chéngshí zǒngshì duì de.}\\
« Ce que tu dis est sensé, mais l'honnêteté a toujours raison. »\\
A : 你能举个例子吗？\\
\emph{Nǐ néng jǔ ge lìzi ma ?}\\
« Peux-tu donner un exemple ? »\\
B : 比如，考试的时候不能作弊。\\
\emph{Bǐrú, kǎoshì de shíhou bù néng zuòbì.}\\
« Par exemple, à l'examen, on ne doit pas tricher. »\\
A : 你说得对！所以，我们应该反对腐败。\\
\emph{Nǐ shuō de duì ! Suǒyǐ, wǒmen yīnggāi fǎnduì fǔbài.}\\
« Tu as raison ! Donc, nous devons lutter contre la corruption. »
\end{textebox}

Observez ce dialogue. Chaque réplique joue un rôle précis : B \emph{approuve}, A \emph{nuance}, B \emph{reconnaît puis objecte}, A \emph{demande un exemple}, puis A \emph{conclut}. C'est le squelette de tout débat scolaire en chinois. Nous allons maintenant démonter chaque pièce.

\courssub{Exprimer l'accord}

\begin{definition}[Les formules d'accord]
Pour dire que l'on partage l'avis de l'autre, le chinois utilise principalement :
\begin{itemize}
\item \zh{我同意。} \emph{Wǒ tóngyì.} — « Je suis d'accord. » (formule neutre, toujours correcte)
\item \zh{我同意你的看法。} \emph{Wǒ tóngyì nǐ de kànfǎ.} — « Je suis d'accord avec ton point de vue. » (plus complet, plus formel)
\item \zh{你说得对！} \emph{Nǐ shuō de duì !} — « Tu as raison ! » (chaleureux, très fréquent à l'oral)
\end{itemize}
\end{definition}

\begin{propriete}[Structure de 你说得对]
La formule \zh{你说得对} se construit ainsi : Sujet (\zh{你}) + verbe (\zh{说}, dire) + \zh{得} + adjectif (\zh{对}, juste). Le mot-outil \cle{得} introduit ici un \emph{complément de manière ou de résultat} : littéralement « tu parles de façon juste ». On peut remplacer l'adjectif : \zh{你说得很好。} \emph{Nǐ shuō de hěn hǎo.} — « Tu parles très bien (le chinois). »
\end{propriete}

\begin{attention}[« D'accord » ne se dit pas 好 dans un débat]
Erreur classique du francophone : traduire « je suis d'accord » par \zh{我好} ou \zh{好}. En chinois, \zh{好} (hǎo) signifie seulement « OK, d'accord » pour accepter une \emph{proposition d'action} (« On va au marché ? — 好！ »). Dans un débat d'\emph{idées}, il faut obligatoirement \zh{同意} (tóngyì) avec son sujet : \zh{我同意}. De même, on ne dit jamais \zh{我是同意} : \zh{同意} est un verbe, pas un adjectif, il ne prend pas \zh{是}.
\end{attention}

\courssub{Exprimer un désaccord poli}

En chinois comme en français, dire brutalement « je ne suis pas d'accord » peut sembler sec. La culture chinoise valorise l'harmonie : on \emph{reconnaît d'abord} le mérite de l'argument adverse, \emph{puis} on objecte avec \zh{但是}.

\begin{definition}[Les formules de désaccord]
\begin{itemize}
\item \zh{我不同意。} \emph{Wǒ bù tóngyì.} — « Je ne suis pas d'accord. » (direct, à utiliser avec précaution)
\item \zh{你说得有道理，但是……} \emph{Nǐ shuō de yǒu dàolǐ, dànshì...} — « Ce que tu dis est sensé, mais... » (formule polie par excellence)
\end{itemize}
\zh{有道理} \emph{yǒu dàolǐ} signifie littéralement « avoir de la raison, être sensé » : c'est le tampon poli qui adoucit le désaccord.
\end{definition}

\begin{propriete}[但是 dànshì : la concession polie]
\zh{但是} (dànshì, « mais ») permet un désaccord poli après avoir reconnu l'argument de l'autre. Structure : \textbf{reconnaissance + 但是 + objection}.\\
\zh{你说得对，但是……} \emph{Nǐ shuō de duì, dànshì...} — « Tu as raison, mais... »\\
Variante plus légère : \zh{可是} (kěshì), identique en sens, un peu plus orale. Attention à la place : \zh{但是} se met \emph{en tête} de la seconde proposition, jamais à la fin comme « mais » ne peut d'ailleurs pas l'être en français.
\end{propriete}

\begin{exemplebox}[Désaccord poli en contexte]
\zh{你说得有道理，但是法律也很重要。}\\
\emph{Nǐ shuō de yǒu dàolǐ, dànshì fǎlǜ yě hěn zhòngyào.}\\
« Ce que tu dis est sensé, mais la loi est aussi importante. »\\[4pt]
\zh{腐败不好，但是反腐败很难。}\\
\emph{Fǔbài bù hǎo, dànshì fǎn fǔbài hěn nán.}\\
« La corruption est mauvaise, mais lutter contre elle est difficile. »\\[4pt]
\zh{我不同意，因为诚实总是可能的。}\\
\emph{Wǒ bù tóngyì, yīnwèi chéngshí zǒngshì kěnéng de.}\\
« Je ne suis pas d'accord, parce que l'honnêteté est toujours possible. »
\end{exemplebox}

\courssub{Relancer l'échange}

Un bon débatteur ne monologue pas : il \cle{relance} son partenaire. Le chinois dispose de trois outils simples et très efficaces.

\begin{definition}[Les formules de relance]
\begin{itemize}
\item \zh{你呢？} \emph{Nǐ ne ?} — « Et toi ? » (renvoie la même question à l'autre)
\item \zh{你怎么看？} \emph{Nǐ zěnme kàn ?} — « Qu'en penses-tu ? » (demande un avis)
\item \zh{为什么？} \emph{Wèishénme ?} — « Pourquoi ? » (demande une justification)
\item \zh{你能举个例子吗？} \emph{Nǐ néng jǔ ge lìzi ma ?} — « Peux-tu donner un exemple ? » (demande une preuve concrète)
\end{itemize}
\end{definition}

\begin{methode}[Relancer en deux gestes]
\begin{enumerate}
\item \textbf{Renvoyer la question} : ajoutez \zh{呢} (ne) après le pronom. \zh{我觉得腐败很危险。你呢？} \emph{Wǒ juéde fǔbài hěn wēixiǎn. Nǐ ne ?} — « Je trouve la corruption très dangereuse. Et toi ? » La particule \zh{呢} reprend la question précédente sans la répéter : économique et élégant.
\item \textbf{Demander la justification} : utilisez \zh{为什么} (wèishénme). \zh{你为什么这么说？} \emph{Nǐ wèishénme zhème shuō ?} — « Pourquoi dis-tu cela ? » La réponse attendra alors \zh{因为} (yīnwèi, « parce que »).
\end{enumerate}
\end{methode}

\courssub{Conclure un échange}

Après l'échange, on tire une conclusion commune avec \zh{所以} (suǒyǐ, « donc, c'est pourquoi »), souvent suivi d'une proposition d'action avec \zh{应该} (yīnggāi, « devoir, il faut »).

\begin{exemplebox}[Formules de conclusion]
\zh{所以，我们应该反对腐败。}\\
\emph{Suǒyǐ, wǒmen yīnggāi fǎnduì fǔbài.}\\
« Donc, nous devons lutter contre la corruption. »\\[4pt]
\zh{所以，诚实比钱更重要。}\\
\emph{Suǒyǐ, chéngshí bǐ qián gèng zhòngyào.}\\
« Donc, l'honnêteté est plus importante que l'argent. »
\end{exemplebox}

\courssub{Tableau récapitulatif des formules du débat}

\begin{center}
\begin{tabularx}{\linewidth}{|X|X|X|X|}
\hline
\rowcolor{popblueL} Fonction & Caractères & Pinyin & Traduction \\
\hline
Accord & 我同意。 & Wǒ tóngyì. & Je suis d'accord. \\
\hline
Accord + avis & 我同意你的看法。 & Wǒ tóngyì nǐ de kànfǎ. & Je partage ton point de vue. \\
\hline
Accord chaleureux & 你说得对！ & Nǐ shuō de duì ! & Tu as raison ! \\
\hline
Désaccord direct & 我不同意。 & Wǒ bù tóngyì. & Je ne suis pas d'accord. \\
\hline
Désaccord poli & 你说得有道理，但是…… & Nǐ shuō de yǒu dàolǐ, dànshì... & C'est sensé, mais... \\
\hline
Relance & 你呢？ & Nǐ ne ? & Et toi ? \\
\hline
Demande d'avis & 你怎么看？ & Nǐ zěnme kàn ? & Qu'en penses-tu ? \\
\hline
Justification & 为什么？ & Wèishénme ? & Pourquoi ? \\
\hline
Exemple & 你能举个例子吗？ & Nǐ néng jǔ ge lìzi ma ? & Peux-tu donner un exemple ? \\
\hline
Conclusion & 所以，我们应该…… & Suǒyǐ, wǒmen yīnggāi... & Donc, nous devons... \\
\hline
\end{tabularx}
\end{center}

\courssub{Schéma du débat : de l'opinion à la conclusion}

\begin{popfigure}
\begin{tikzpicture}[
node distance=0.8cm and 1.5cm,
etape/.style={rectangle, rounded corners=3pt, draw=popblue, fill=popblueL, thick, align=center, font=\small, inner sep=5pt},
choix/.style={rectangle, rounded corners=3pt, draw=poporange, fill=poporangeL, thick, align=center, font=\small, inner sep=5pt},
fleche/.style={-{Stealth[length=2.5mm]}, thick, popdark}
]
\node[etape, align=center] (op){Opinion\\我认为…… / 我觉得……};
\node[choix, below left=0.8cm and 1.0cm of op, align=center] (acc){Accord\\我同意！\\你说得对！};
\node[choix, below right=0.8cm and 1.0cm of op, align=center] (des){Désaccord poli\\有道理，但是……};
\node[etape, below=2.2cm of op, align=center] (rel){Relance\\你呢？ 为什么？\\举个例子？};
\node[etape, below=0.8cm of rel, draw=popgreen, fill=popgreenL, align=center] (conc){Conclusion\\所以，我们应该……};
\draw[fleche] (op) -- (acc);
\draw[fleche] (op) -- (des);
\draw[fleche] (acc) -- (rel);
\draw[fleche] (des) -- (rel);
\draw[fleche] (rel) -- (conc);
\end{tikzpicture}
\legende{Organigramme du mini-débat : opinion, réaction, relance, conclusion.}
\end{popfigure}

\courssub{Travail des tons : quatre mots à musique précise}

La prononciation des tons est décisive dans un débat : un ton faux peut rendre une formule incompréhensible. Travaillez ces quatre mots-clés à voix haute.

\begin{center}
\begin{tabularx}{\linewidth}{|X|X|X|X|}
\hline
\rowcolor{popblueL} Mot & Tons & Profil mélodique & Piège \\
\hline
同意 tóngyì & 2 + 4 & monte puis tombe & ne pas dire tōngyì (1+4) \\
\hline
但是 dànshì & 4 + 4 & deux chutes sèches & ne pas adoucir le 2e ton \\
\hline
为什么 wèishénme & 4 + 2 + neutre & tombe, monte, léger & me final sans ton, très bref \\
\hline
应该 yīnggāi & 1 + 1 & deux tons hauts et plats & garder les deux tons hauts \\
\hline
\end{tabularx}
\end{center}

\begin{attention}[Le ton neutre de 呢 et de me]
Dans \zh{你呢} (Nǐ ne) et \zh{为什么} (wèishénme), la dernière syllabe est au \cle{ton neutre} : courte, légère, sans mélodie. Les francophones ont tendance à lui donner un ton descendant (comme un accent français final) : résistez, laissez la syllabe « tomber en douceur ».
\end{attention}

\courssub{Exercice résolu et entraînement}

\begin{exoresolu}[Compléter un débat]
Énoncé — Complétez ce dialogue avec les formules étudiées : 我同意 / 但是 / 你呢 / 所以 / 举个例子.\\
A : 我觉得腐败对喀麦隆很危险。 \emph{Wǒ juéde fǔbài duì Kāmàilóng hěn wēixiǎn.}\\
B : \_\_\_\_\_\_\_\_！你说得很有道理。\\
A : \_\_\_\_\_\_\_\_？你觉得呢？\\
B : 你说得有道理，\_\_\_\_\_\_\_\_ 反腐败不容易。\\
A : 你能 \_\_\_\_\_\_\_\_ 吗？\\
B : 比如，考试作弊也是一种腐败。\_\_\_\_\_\_\_\_，我们应该从小诚实。
\tcblower
Solution détaillée :\\
B : \textbf{我同意}！你说得很有道理。 (accord après une opinion : 我同意)\\
A : \textbf{你呢}？你觉得呢？ (relance avec 呢 après le pronom)\\
B : 你说得有道理，\textbf{但是}反腐败不容易。 (désaccord poli : reconnaissance + 但是)\\
A : 你能\textbf{举个例子}吗？ (demande d'exemple : 举个例子, « donner un exemple »)\\
B : 比如…… \textbf{所以}，我们应该从小诚实。 (conclusion avec 所以 + 应该)\\
Traduction complète : « A : Je trouve la corruption très dangereuse pour le Cameroun. B : Je suis d'accord ! Ce que tu dis est très sensé. A : Et toi, qu'en penses-tu ? B : C'est sensé, mais lutter contre la corruption n'est pas facile. A : Peux-tu donner un exemple ? B : Par exemple, tricher à l'examen est aussi une forme de corruption. Donc, nous devons être honnêtes dès le plus jeune âge. »
\end{exoresolu}

\begin{exercice}[À vous de débattre]
Par binôme, construisez un mini-débat de six répliques sur le sujet : \zh{学生应该诚实吗？} (Xuésheng yīnggāi chéngshí ma ? — « Les élèves doivent-ils être honnêtes ? »). Votre dialogue doit contenir obligatoirement : une opinion (\zh{我认为} ou \zh{我觉得}), un accord (\zh{我同意} ou \zh{你说得对}), un désaccord poli (\zh{有道理，但是……}), une relance (\zh{你呢？} ou \zh{为什么？}), une demande d'exemple (\zh{举个例子}) et une conclusion (\zh{所以……}). Attention aux tons de \zh{同意} (2+4) et \zh{但是} (4+4) lors de la présentation orale.
\end{exercice}

\begin{corrige}[Exercice « À vous de débattre »]
Modèle de production attendue (une solution parmi d'autres) :\\
A : 我认为学生应该诚实。\emph{Wǒ rènwéi xuésheng yīnggāi chéngshí.} « Je pense que les élèves doivent être honnêtes. »\\
B : 我同意！诚实很重要。\emph{Wǒ tóngyì ! Chéngshí hěn zhòngyào.} « D'accord ! L'honnêteté est très importante. »\\
A : 你说得有道理，但是有时候很难。\emph{Nǐ shuō de yǒu dàolǐ, dànshì yǒu shíhou hěn nán.} « C'est sensé, mais parfois c'est difficile. »\\
B : 为什么难？\emph{Wèishénme nán ?} « Pourquoi est-ce difficile ? »\\
A : 比如，考试的时候有人作弊。\emph{Bǐrú, kǎoshì de shíhou yǒu rén zuòbì.} « Par exemple, à l'examen, certains trichent. »\\
B : 你说得对！所以，我们应该反对作弊。\emph{Nǐ shuō de duì ! Suǒyǐ, wǒmen yīnggāi fǎnduì zuòbì.} « Tu as raison ! Donc, nous devons lutter contre la triche. »\\
Critères d'évaluation : présence des six formules, exactitude des tons, fluidité de l'échange.
\end{corrige}

\begin{aretenir}[L'essentiel de la section]
\begin{itemize}
\item Accord : \zh{我同意} / \zh{你说得对} — jamais \zh{好} pour une idée.
\item Désaccord poli : reconnaissance + \zh{但是} + objection (\zh{你说得有道理，但是……}).
\item Relance : \zh{你呢？} / \zh{你怎么看？} / \zh{为什么？} / \zh{举个例子吗？}
\item Conclusion : \zh{所以，我们应该……}
\item Tons à soigner : tóngyì (2+4), dànshì (4+4), wèishénme (4+2+neutre), yīnggāi (1+1).
\end{itemize}
\end{aretenir}

\coursec{Présenter un exposé simple}

Dans la section précédente, nous avons appris à réagir dans un débat : dire que l'on est d'accord (\zh{我同意} Wǒ tóngyì), exprimer un désaccord poli (\zh{我不太同意} Wǒ bú tài tóngyì) et relancer la parole (\zh{你怎么看？} Nǐ zěnme kàn ?). Mais un débat ne commence jamais de nulle part : il faut d'abord que quelqu'un prenne la parole pour \cle{présenter le sujet}. C'est exactement le rôle de l'\cle{exposé}. Dans cette section, nous allons apprendre à construire, pas à pas, un petit exposé oral en chinois sur le thème de la corruption, avec une structure claire que l'examinateur reconnaîtra immédiatement.

\courssub{Observer d'abord : un mini-exposé modèle}

Avant toute règle, écoutons (et lisons) un exposé complet, très simple, tel qu'un élève de Première peut le produire après cette leçon.

\begin{textebox}[Un exposé modèle : 腐败的问题]
大家好！今天我要谈一谈腐败的问题。\\
Dàjiā hǎo ! Jīntiān wǒ yào tán yi tán fǔbài de wèntí.\\
« Bonjour à tous ! Aujourd'hui, je vais parler un peu du problème de la corruption. »\\
首先，腐败是一个大问题。\\
Shǒuxiān, fǔbài shì yí ge dà wèntí.\\
« D'abord, la corruption est un grand problème. »\\
然后，它影响国家的发展。\\
Ránhòu, tā yǐngxiǎng guójiā de fāzhǎn.\\
« Ensuite, elle nuit au développement du pays. »\\
最后，我们应该诚实，反对腐败。\\
Zuìhòu, wǒmen yīnggāi chéngshí, fǎnduì fǔbài.\\
« Enfin, nous devons être honnêtes et lutter contre la corruption. »\\
谢谢大家！\\
Xièxie dàjiā !\\
« Merci à tous ! »
\end{textebox}

\textbf{Glossaire rapide}

\begin{center}
\begin{tabularx}{\linewidth}{|X|X|X|X|}
\hline
\rowcolor{popblueL} Caractères & Pinyin & Nature & Traduction \\
\hline
腐败 & fǔbài & nom / adj. & la corruption ; corrompu \\
\hline
问题 & wèntí & nom & problème, question \\
\hline
影响 & yǐngxiǎng & verbe / nom & influencer, affecter ; l'influence \\
\hline
国家 & guójiā & nom & le pays, l'État \\
\hline
发展 & fāzhǎn & nom / verbe & le développement ; se développer \\
\hline
诚实 & chéngshí & adj. & honnête \\
\hline
反对 & fǎnduì & verbe & s'opposer à, lutter contre \\
\hline
\end{tabularx}
\end{center}

\textbf{Questions de compréhension (oral) :}
\begin{enumerate}
\item \zh{他要谈什么？} Tā yào tán shénme ? — De quoi va-t-il parler ?
\item \zh{腐败影响什么？} Fǔbài yǐngxiǎng shénme ? — Qu'est-ce que la corruption affecte ?
\item \zh{我们应该怎么做？} Wǒmen yīnggāi zěnme zuò ? — Que devons-nous faire ?
\end{enumerate}

\courssub{La structure en trois parties : 开头、中间、结尾}

Observez le mini-exposé : il ne commence pas directement par les arguments, et il ne se termine pas brutalement. Tout exposé chinois, comme en français, suit un plan en \cle{trois parties} :

\begin{itemize}
\item \zh{开头} kāitóu : l'\cle{introduction} — on salue et on annonce le sujet ;
\item \zh{中间} zhōngjiān : le \cle{développement} — on donne ses idées, une par une ;
\item \zh{结尾} jiéwěi : la \cle{conclusion} — on résume, on donne son avis et on remercie.
\end{itemize}

\begin{popfigure}
\begin{tikzpicture}[node distance=0.35cm]
\tikzstyle{etape}=[rectangle, rounded corners=3pt, draw=popblue, fill=popblueL, align=center, minimum height=0.9cm, font=\small]
\tikzstyle{fleche}=[-{Stealth[length=2.5mm]}, thick, poporange]
\node[etape, fill=poppurpleL, draw=poppurple, align=center] (n1){开头\\ kāitóu\\ 大家好！\\ 今天我要谈……};
\node[etape, right=of n1, align=center] (n2){首先\\ shǒuxiān\\ 腐败是……};
\node[etape, right=of n2, align=center] (n3){然后\\ ránhòu\\ 它影响……};
\node[etape, right=of n3, align=center] (n4){最后\\ zuìhòu\\ 我们应该……};
\node[etape, fill=poppurpleL, draw=poppurple, right=of n4, align=center] (n5){结尾\\ jiéwěi\\ 所以……\\ 谢谢大家！};
\draw[fleche] (n1) -- (n2);
\draw[fleche] (n2) -- (n3);
\draw[fleche] (n3) -- (n4);
\draw[fleche] (n4) -- (n5);
\node[above=0.15cm of n2, font=\footnotesize\itshape, popmuted] {中间 zhōngjiān : le développement};
\end{tikzpicture}
\legende{Frise de l'exposé : de l'introduction (开头) à la conclusion (结尾), en passant par les trois connecteurs du développement.}
\end{popfigure}

\begin{aretenir}[Le squelette à mémoriser]
Un exposé simple = \cle{salutation} + \cle{annonce du sujet} + \cle{deux ou trois arguments reliés par des connecteurs} + \cle{conclusion avec 所以} + \cle{remerciement}. Si vous connaissez ce squelette par cœur, vous pouvez parler de n'importe quel sujet en changeant seulement le contenu.
\end{aretenir}

\courssub{L'introduction : saluer et annoncer le sujet}

\textbf{Étape 1 : saluer.} La formule standard est \zh{大家好！} Dàjiā hǎo ! — « Bonjour à tous ! ». Devant des professeurs ou un jury, on préfère \zh{老师好，同学们好！} Lǎoshī hǎo, tóngxuémen hǎo ! — « Bonjour professeur, bonjour camarades ! ».

\textbf{Étape 2 : annoncer le sujet.} La formule clé de cette leçon est :

\zh{今天我要谈一谈腐败的问题。} \emph{Jīntiān wǒ yào tán yi tán fǔbài de wèntí.} « Aujourd'hui, je vais parler un peu du problème de la corruption. »

Cette phrase contient deux outils grammaticaux essentiels.

\begin{propriete}[要 + verbe : l'intention ; 谈一谈 : la réduplication]
\textbf{1. 要 + verbe} exprime l'\cle{intention} ou le futur proche : \zh{我要谈} Wǒ yào tán = « je vais parler / j'ai l'intention de parler ». Structure : Sujet + 要 + verbe (+ complément).\\
\emph{Exemples :} \zh{我要喝水。} Wǒ yào hē shuǐ. « Je vais boire de l'eau. » / \zh{明天我们要考试。} Míngtiān wǒmen yào kǎoshì. « Demain, nous allons passer un examen. »\\
\textbf{2. La réduplication du verbe} (谈 → 谈一谈, ou 谈谈) \cle{adoucit} l'action : elle signifie « faire un peu, brièvement ». \zh{谈一谈} tán yi tán = « parler un peu ». C'est une marque de modestie très appréciée en chinois : l'orateur ne prétend pas tout dire.\\
\emph{Autres exemples :} \zh{看一看} kàn yi kàn « jeter un coup d'œil » ; \zh{想一想} xiǎng yi xiǎng « réfléchir un moment » ; \zh{介绍一下} jièshào yíxià « présenter brièvement ».\\
\textbf{Attention :} le 一 de la réduplication est souvent omis à l'oral : \zh{谈谈} = \zh{谈一谈}.
\end{propriete}

\begin{attention}[要 ou 想 ?]
\zh{要} yào exprime une intention ferme, presque une décision (« je vais le faire »), tandis que \zh{想} xiǎng exprime un simple désir (« j'aimerais »). Pour annoncer un exposé, on dit bien \zh{我要谈一谈} Wǒ yào tán yi tán (« je vais parler de… ») : c'est un engagement devant l'auditoire. \zh{我想谈一谈} Wǒ xiǎng tán yi tán reste correct, mais plus hésitant.
\end{attention}

Pour annoncer n'importe quel sujet, il suffit de remplacer le nom après \zh{谈一谈} : \zh{今天我要谈一谈足球。} Jīntiān wǒ yào tán yi tán zúqiú. « Aujourd'hui, je vais parler un peu du football. »

\courssub{Le développement : les connecteurs 首先、然后、最后}

Un exposé sans connecteurs ressemble à une liste de phrases jetées en désordre. Le chinois, comme le français, utilise des \cle{mots de liaison chronologiques} pour guider l'auditeur :

\begin{center}
\begin{tabularx}{\linewidth}{|X|X|X|X|}
\hline
\rowcolor{popblueL} Caractères & Pinyin & Fonction & Équivalent français \\
\hline
首先 & shǒuxiān & premier argument & d'abord, premièrement \\
\hline
然后 & ránhòu & argument suivant & ensuite, puis \\
\hline
最后 & zuìhòu & dernier point & enfin, pour finir \\
\hline
所以 & suǒyǐ & conclusion logique & donc, c'est pourquoi \\
\hline
\end{tabularx}
\end{center}

\begin{propriete}[Place de 首先、然后、最后]
Ces connecteurs se placent \cle{en début de phrase}, suivis d'une virgule (，). Deux positions sont possibles :\\
— \cle{devant le sujet} : \zh{首先，腐败是一个大问题。} Shǒuxiān, fǔbài shì yí ge dà wèntí. « D'abord, la corruption est un grand problème. »\\
— \cle{après le sujet} : \zh{我们首先看腐败的原因。} Wǒmen shǒuxiān kàn fǔbài de yuányīn. « Nous examinons d'abord les causes de la corruption. »\\
Les deux sont correctes ; la position initiale est la plus fréquente à l'oral.
\end{propriete}

\begin{attention}[Erreurs fréquentes des francophones]
\textbf{1.} Ne traduisez pas « d'abord » par \zh{第一} dans un exposé parlé : \zh{第一} dì-yī sert à énumérer de façon plus écrite et soutenue (第一、第二、第三). À l'oral simple, utilisez \zh{首先}.\\
\textbf{2.} Ne placez jamais le connecteur en fin de phrase comme on le fait parfois en français (« la corruption est un grand problème, d'abord ») : en chinois, il ouvre obligatoirement l'énoncé.\\
\textbf{3.} \zh{然后} ránhòu signifie « ensuite » (ordre), pas « donc » (conséquence). Pour « donc », c'est \zh{所以} suǒyǐ. Confondre les deux est la faute la plus fréquente.\\
\textbf{4.} \zh{影响} yǐngxiǎng est un verbe \cle{transitif direct} : on dit \zh{影响国家的发展} yǐngxiǎng guójiā de fāzhǎn « affecter le développement du pays », sans préposition. N'ajoutez pas \zh{对} après 影响 sur le modèle français « avoir une influence \emph{sur} ».
\end{attention}

\textbf{Le contenu du développement.} Pour un exposé de niveau Première, deux idées suffisent, et le thème de la corruption s'y prête naturellement :

\begin{itemize}
\item \cle{une cause} : \zh{腐败的原因} fǔbài de yuányīn — « les causes de la corruption ». Exemple : \zh{首先，腐败的原因很多：有的人想要快钱。} Shǒuxiān, fǔbài de yuányīn hěn duō : yǒu de rén xiǎng yào kuài qián. « D'abord, les causes de la corruption sont nombreuses : certaines personnes veulent de l'argent facile. » (\zh{快钱} kuài qián, littéralement « l'argent rapide », désigne l'argent gagné facilement et malhonnêtement.)
\item \cle{une conséquence} : \zh{腐败的影响} fǔbài de yǐngxiǎng — « les effets de la corruption ». Exemple : \zh{然后，腐败影响国家的发展，也影响我们的生活。} Ránhòu, fǔbài yǐngxiǎng guójiā de fāzhǎn, yě yǐngxiǎng wǒmen de shēnghuó. « Ensuite, la corruption nuit au développement du pays et aussi à notre vie. »
\end{itemize}

Remarquez le mot-outil \zh{的} de : il relie le thème à son aspect, exactement comme « de » en français : \zh{腐败的原因} = « les causes \emph{de} la corruption ». Remarquez aussi \zh{也} yě « aussi », placé \cle{avant le verbe} : \zh{也影响} yě yǐngxiǎng « affecte aussi ».

\courssub{La conclusion : 所以 et les remerciements}

La conclusion fait deux choses : elle \cle{tire une leçon} du développement, puis elle \cle{remercie} l'auditoire.

\textbf{1. Tirer une leçon avec 所以.} \zh{所以} suǒyǐ « donc, c'est pourquoi » introduit la conséquence logique :

\zh{所以，我们应该诚实，反对腐败。} \emph{Suǒyǐ, wǒmen yīnggāi chéngshí, fǎnduì fǔbài.} « C'est pourquoi nous devons être honnêtes et lutter contre la corruption. »

Le verbe modal \zh{应该} yīnggāi « devoir (il faut) » est ici parfait pour formuler un devoir moral : Sujet + 应该 + verbe. Exemple : \zh{学生应该努力学习。} Xuésheng yīnggāi nǔlì xuéxí. « Les élèves doivent étudier avec effort. »

\textbf{2. Remercier.} La formule de clôture obligatoire est \zh{谢谢大家！} Xièxie dàjiā ! — « Merci à tous ! ». On peut ajouter \zh{我的报告结束了。} Wǒ de bàogào jiéshù le. « Mon exposé est terminé. » (\zh{结束} jiéshù, « se terminer », suivi de \zh{了} le qui marque l'accomplissement.)

\begin{methode}[Structurer un exposé simple en cinq gestes]
\begin{enumerate}
\item \textbf{Saluer} : \zh{大家好！} Dàjiā hǎo !
\item \textbf{Annoncer le sujet} : \zh{今天我要谈一谈……} Jīntiān wǒ yào tán yi tán … « Aujourd'hui, je vais parler un peu de… »
\item \textbf{Développer deux arguments} reliés par \zh{首先} shǒuxiān (d'abord) et \zh{然后} ránhòu (ensuite) ; ajoutez \zh{最后} zuìhòu (enfin) pour un troisième point.
\item \textbf{Conclure} avec \zh{所以} suǒyǐ + \zh{我们应该……} wǒmen yīnggāi … « donc, nous devons… »
\item \textbf{Remercier} : \zh{谢谢大家！} Xièxie dàjiā !
\end{enumerate}
Avant de parler, écrivez ce squelette au brouillon et glissez-y vos mots de vocabulaire : vous n'aurez plus qu'à le dire.
\end{methode}

\courssub{Conseils de prise de parole}

Un bon plan ne suffit pas : la \cle{forme orale} compte autant que le fond.

\begin{itemize}
\item \cle{Parlez lentement.} En chinois, la vitesse tue les tons. Un débit lent et régulier est perçu comme maîtrisé, jamais comme faible.
\item \cle{Marquez les tons}, surtout sur les mots-clés du thème : \zh{腐败} fǔbài (3\textsuperscript{e} + 4\textsuperscript{e} ton, descendez bien sur bài), \zh{影响} yǐngxiǎng (3\textsuperscript{e} + 3\textsuperscript{e}, le premier devient 2\textsuperscript{e} : yíngxiǎng), \zh{诚实} chéngshí (2\textsuperscript{e} + 2\textsuperscript{e}, montez deux fois).
\item \cle{Regardez l'auditoire}, pas votre feuille : mémorisez le squelette (salut → sujet → 首先 → 然后 → 所以 → merci) et ne gardez en main que les mots difficiles.
\item \cle{Respirez} entre les parties : une petite pause après \zh{首先，…} donne du poids à chaque argument.
\end{itemize}

\begin{savaistu}[À l'examen, la structure paie]
Au baccalauréat et dans les évaluations MINESEC d'expression orale, les correcteurs disposent d'une grille où la \cle{cohérence du propos} compte autant que la richesse du vocabulaire. Un exposé simple mais bien structuré — avec \zh{首先}、\zh{然后}、\zh{最后}、\zh{所以} clairement prononcés — obtient une meilleure note qu'un exposé riche en mots savants mais désorganisé. En Chine même, les élèves apprennent dès le primaire à « ouvrir » (开头) et « fermer » (结尾) leurs prises de parole : c'est une marque d'éducation autant que de langue.
\end{savaistu}

\courssub{Entraînement guidé}

\begin{exoresolu}[Compléter le squelette d'un exposé]
\textbf{Énoncé.} Complétez cet exposé sur le sport avec les connecteurs et formules étudiés : \\
\zh{……！今天我要……一谈足球。……，足球很有意思。……，它对身体很好。……，我们应该多运动。……大家！}
\tcblower
\textbf{Solution détaillée.} On suit le squelette en cinq gestes :\\
1. Salutation : \zh{大家好！} Dàjiā hǎo !\\
2. Annonce du sujet : \zh{今天我要谈一谈足球。} Jīntiān wǒ yào tán yi tán zúqiú. (要 + verbe, réduplication 谈一谈)\\
3. Premier argument : \zh{首先，足球很有意思。} Shǒuxiān, zúqiú hěn yǒu yìsi. « D'abord, le football est très intéressant. »\\
4. Deuxième argument : \zh{然后，它对身体很好。} Ránhòu, tā duì shēntǐ hěn hǎo. « Ensuite, il est bon pour la santé. » (ici, \zh{对} duì « pour » est nécessaire : \zh{对身体好} = « bon \emph{pour} le corps » ; ne pas le confondre avec 影响, qui lui est transitif direct.)\\
5. Conclusion + remerciement : \zh{所以，我们应该多运动。谢谢大家！} Suǒyǐ, wǒmen yīnggāi duō yùndòng. Xièxie dàjiā ! « Donc, nous devons faire plus de sport. Merci à tous ! »\\
\textbf{Vérification :} chaque connecteur est en début de phrase ; 所以 introduit bien la conclusion, pas un simple « ensuite ».
\end{exoresolu}

\begin{exercice}[À vous de parler]
En vous inspirant du modèle, préparez par écrit (5 à 7 phrases) puis présentez oralement un mini-exposé sur l'un des sujets suivants : \zh{腐败} fǔbài (la corruption), \zh{诚实} chéngshí (l'honnêteté) ou \zh{学习汉语} xuéxí Hànyǔ (l'apprentissage du chinois). Obligatoire : \zh{大家好}、\zh{我要谈一谈}、\zh{首先}、\zh{然后}、\zh{所以}、\zh{谢谢大家}. Votre voisin écoute et vérifie que les cinq formules sont présentes et que les tons sont marqués.
\end{exercice}

Vous disposez maintenant d'un squelette complet et de formules prêtes à l'emploi. Dans la prochaine section, nous affinerons la \cle{prononciation} et le travail des tons pour que cet exposé sonne juste, avant de passer aux jeux de rôle et au débat guidé.

\coursec{Travail des tons et de la prononciation}

Vous savez maintenant présenter un petit exposé sur la corruption : vous avez les mots, les structures et un plan. Mais un exposé en chinois ne convainc que si la \cle{prononciation} est correcte. En chinois, chaque syllabe porte un \cle{ton}, c'est-à-dire une mélodie de la voix qui fait partie du mot, exactement comme ses voyelles et ses consonnes. Un ton faux, et votre auditoire ne comprend plus — ou comprend autre chose. Cette section est donc un entraînement intensif : nous allons observer les tons des mots-clés de notre thème, apprendre les règles qui les modifient, puis nous entraîner avec des gestes, des cartes et une répétition en chaîne.

\courssub{Observer : les tons des mots-clés du thème}

Commençons par écouter attentivement les cinq mots les plus importants de notre leçon. Le professeur les lit lentement, deux fois, sans que vous regardiez le pinyin. Quelle mélodie entendez-vous sur chaque syllabe ?

\begin{center}
\begin{tabularx}{\linewidth}{|X|X|X|X|}\hline
\rowcolor{popblueL} Caractères & Pinyin & Nature & Traduction \\ \hline
腐败 & fǔbài & nom & corruption \\ \hline
贪污 & tānwū & verbe & détournement (de fonds) \\ \hline
受贿 & shòuhuì & verbe & accepter des pots-de-vin \\ \hline
诚实 & chéngshí & adjectif & honnête \\ \hline
法律 & fǎlǜ & nom & loi, législation \\ \hline
\end{tabularx}
\end{center}

Chaque combinaison de tons a sa propre difficulté pour un francophone :

\begin{itemize}
\item \zh{腐败} fǔbài (3 + 4) : la voix descend d'abord très bas sur \zh{腐}, puis chute franchement sur \zh{败}. Contraste fort entre les deux syllabes.
\item \zh{贪污} tānwū (1 + 1) : deux syllabes hautes, plates et tendues, comme une note de musique tenue. Ne laissez pas la voix retomber à la fin.
\item \zh{受贿} shòuhuì (4 + 4) : deux chutes successives, sèches et énergiques. C'est la combinaison la plus difficile : il faut remonter la voix \emph{entre} les deux syllabes pour pouvoir redescendre.
\item \zh{诚实} chéngshí (2 + 2) : deux montées, comme deux questions françaises qui s'enchaînent.
\item \zh{法律} fǎlǜ (3 + 4) : attention à la voyelle \zh{ü} de \zh{lǜ} : lèvres arrondies comme pour dire « ou », mais on prononce « i ».
\end{itemize}

\begin{popfigure}
\begin{tikzpicture}[scale=1.1, >=Stealth]
% Grille de référence (hauteur 4 = haut, 0 = bas)
\draw[popline, dashed] (0,0) -- (2.2,0);
\draw[popline, dashed] (0,1) -- (2.2,1);
\draw[popline, dashed] (0,2) -- (2.2,2);
\draw[popline, dashed] (0,3) -- (2.2,3);
\draw[popline, dashed] (0,4) -- (2.2,4);
\node[left, font=\small\itshape, text=popmuted] at (-0.2,4) {haut};
\node[left, font=\small\itshape, text=popmuted] at (-0.2,0) {bas};

% 1er ton : haut plat (y=3.5)
\draw[very thick, popblue] (0.1,3.5) -- (2.0,3.5);
\node[popblue, font=\small, align=center] at (-0.5,3.5) {1er ton};
\node[popblue, font=\large] at (1.05,3.9) {\zh{贪} tān};

% 2e ton : montant (y=1 -> 3.5)
\draw[very thick, popgreen] (0.1,1.0) -- (2.0,3.5);
\node[popgreen, font=\small, align=center] at (-0.5,2.2) {2e ton};
\node[popgreen, font=\large] at (1.05,2.2) {\zh{诚} chéng};

% 3e ton : descendant-remontant (y=2.5 -> 0.5 -> 2.5)
\draw[very thick, poporange] (0.1,2.5) .. controls (0.7,0.3) and (1.4,0.3) .. (2.0,2.5);
\node[poporange, font=\small, align=center] at (-0.5,1.5) {3e ton};
\node[poporange, font=\large] at (1.05,0.6) {\zh{腐} fǔ};

% 4e ton : chute (y=3.5 -> 0.5)
\draw[very thick, poppink] (0.1,3.5) -- (2.0,0.5);
\node[poppink, font=\small, align=center] at (-0.5,2.0) {4e ton};
\node[poppink, font=\large] at (1.05,2.0) {\zh{贿} huì};

% Ton neutre : point léger (y=2.0)
\fill[popmuted] (1.05,2.0) circle (2.5pt);
\node[popmuted, font=\small, align=center] at (-0.5,2.0) {ton neutre};
\node[popmuted, font=\large] at (1.6,2.0) {\zh{呢} ne};
\end{tikzpicture}
\legende{Les quatre tons et le ton neutre, avec un mot du thème placé sur chaque courbe : \zh{贪} tān (1er), \zh{诚} chéng (2e), \zh{腐} fǔ (3e), \zh{贿} huì (4e), \zh{呢} ne (neutre).}
\end{popfigure}

\begin{textebox}[Fǔbài ou fùbái ? Quand un mauvais ton change tout]
Écoutez ces deux suites de syllabes :\\
\zh{腐败} fǔbài : la corruption\\
\zh{复白} fùbái : (suite de syllabes sans sens courant)\\
Si vous prononcez \zh{腐败} avec un 4e ton puis un 2e ton, vous dites \emph{fùbái}, et votre auditoire chinois ne comprend plus rien : ces syllabes ne forment pas de mot.\\
Un mauvais ton ne fait pas seulement « accent étranger » : il peut changer le sens du mot ou le rendre totalement incompréhensible. C'est pourquoi, dans votre exposé, le mot \zh{腐败} fǔbài doit être prononcé avec soin : 3e ton bas sur \zh{腐}, 4e ton sec sur \zh{败}.
\end{textebox}

\courssub{Le sandhi du 3e ton : quand deux 3e tons se rencontrent}

Observons deux expressions que vous utiliserez dans le débat :

\zh{影响} yǐngxiǎng (influence) — écrit 3 + 3, prononcé 2 + 3 : \emph{yíngxiǎng}.

\zh{很好} hěn hǎo (très bien) — écrit 3 + 3, prononcé 2 + 3 : \emph{hén hǎo}.

\begin{propriete}[Règle du sandhi du 3e ton]
Quand deux syllabes au 3e ton se suivent, la \textbf{première} se prononce comme un 2e ton (montant). Seule l'écriture pinyin garde les deux accents d'origine.
\begin{itemize}
\item Structure : 3 + 3 → 2 + 3 (à l'oral uniquement).
\item Exemples : \zh{影响} yǐngxiǎng → \emph{yíngxiǎng} ; \zh{很好} hěn hǎo → \emph{hén hǎo} ; \zh{可以} kěyǐ → \emph{kéyǐ} ; \zh{你好} nǐ hǎo → \emph{ní hǎo}.
\item Exception : si le premier mot est un monosyllabe qu'on veut insister (« moi, je... »), on peut garder un demi-3e ton (la voix descend bas sans remonter).
\end{itemize}
\end{propriete}

Comparaison avec le français : le français fait quelque chose de semblable avec la liaison — on écrit « les amis » mais on prononce « lé-z-amis ». La forme écrite et la forme prononcée ne coïncident pas toujours. En chinois, le pinyin écrit le ton d'origine, mais la bouche prononce le ton transformé.

\begin{exemplebox}[Le sandhi dans nos phrases de débat]
\zh{腐败的影响很大。} Fǔbài de yǐngxiǎng hěn dà. (prononcé : \emph{yíngxiǎng}) — « L'influence de la corruption est grande. »\\
\zh{你的意见很好。} Nǐ de yìjiàn hěn hǎo. (prononcé : \emph{hén hǎo}) — « Ton avis est très bon. »\\
\zh{我觉得你可以这样说。} Wǒ juéde nǐ kěyǐ zhèyàng shuō. (prononcé : \emph{kéyǐ}) — « Je pense que tu peux le dire comme ça. »
\end{exemplebox}

\courssub{Le ton neutre : léger, court, sans mélodie}

Certaines syllabes perdent leur ton et se prononcent \cle{à la légère}, courtes et sans mélodie propre, un peu comme le « e » français de « je te le dis ». C'est le \cle{ton neutre}, noté sans accent dans le pinyin.

\begin{exemplebox}[Trois cas de ton neutre dans notre dialogue]
\zh{我觉得} wǒ juéde — « je pense » : \zh{得} est neutre, presque murmuré.\\
\zh{说得对} shuō de duì — « (tu) as raison » : la particule \zh{得} est neutre.\\
\zh{你呢？} nǐ ne ? — « et toi ? » : la particule \zh{呢} est toujours neutre.
\end{exemplebox}

\begin{attention}[Ne « tonifiez » pas les particules]
Les francophones, en lisant le pinyin, ont tendance à donner un ton plein à chaque syllabe. Erreur fréquente : dire \emph{nǐ né} avec un 2e ton montant sur \zh{呢}. Les particules \zh{吗}, \zh{呢}, \zh{吧}, \zh{的}, \zh{得}, \zh{了} sont presque toujours neutres : courtes, basses, légères. C'est cette légèreté qui donne le rythme naturel du chinois.
\end{attention}

\courssub{Les 4e tons successifs : descendre franchement, sans hésiter}

Notre thème contient plusieurs mots en 4 + 4 ou des suites de 4e tons :

\begin{center}
\begin{tabularx}{\linewidth}{|X|X|X|X|}\hline
\rowcolor{popblueL} Caractères & Pinyin & Nature & Traduction \\ \hline
受贿 & shòuhuì & verbe & accepter des pots-de-vin \\ \hline
但是 & dànshì & conjonction & mais \\ \hline
同意 & tóngyì & verbe & être d'accord \\ \hline
意见 & yìjiàn & nom & avis, opinion \\ \hline
\end{tabularx}
\end{center}

\begin{propriete}[Enchaîner deux 4e tons]
Le 4e ton part du haut et chute vers le bas, sec et décidé, comme un « non ! » catégorique en français. Quand deux 4e tons se suivent (\zh{受贿} shòuhuì, \zh{但是} dànshì), il faut \textbf{remonter la voix entre les deux syllabes} pour pouvoir chuter une deuxième fois. Chaque chute est complète : on ne fait pas une seule longue descente étalée sur les deux syllabes.
\end{propriete}

Entraînement : frappez deux fois dans vos mains, une fois par syllabe, en disant \emph{shòu — huì}, \emph{dàn — shì}, \emph{yì — jiàn}. Le rythme sec des mains aide la voix à redescendre franchement.

\courssub{La question avec 吗 : les tons ne changent pas}

\begin{attention}[Le piège de la mélodie montante en fin de question]
En français, on monte la voix à la fin d'une question : « Tu es \emph{d'accord} ? » monte. Les francophones transposent ce réflexe en chinois et déforment le dernier mot. Or en chinois, la question avec \zh{吗} garde les tons du mot : c'est la particule \zh{吗}, neutre et légère, qui porte à elle seule l'interrogation. La mélodie de la phrase reste celle d'une phrase affirmative.
\end{attention}

\begin{exemplebox}[La question qui ne monte pas]
\zh{你同意吗？} Nǐ tóngyì ma ? — « Es-tu d'accord ? »\\
Le mot \zh{同意} garde ses tons 2 + 4 : \emph{tóngyì} monte puis chute, exactement comme dans l'affirmative \zh{我同意} wǒ tóngyì « je suis d'accord ». Seul \zh{吗}, neutre, s'ajoute à la fin. Ne dites jamais \emph{tóngyí} en montant à la française.
\end{exemplebox}

\courssub{Mémoriser les tons avec le corps}

\begin{methode}[Le geste des tons : la main trace la mélodie]
Pour ancrer chaque mot-clé dans la mémoire, associez-lui un geste de la main qui trace la courbe du ton dans l'air :
\begin{enumerate}
\item 1er ton : la main avance à l'horizontale, bien haut, comme une baguette de chef d'orchestre (\zh{贪污} tānwū).
\item 2e ton : la main monte en diagonale, du bas vers le haut (\zh{诚实} chéngshí, deux montées).
\item 3e ton : la main descend puis remonte, comme un creux (\zh{腐败} fǔbài : creux puis chute).
\item 4e ton : la main tombe d'un coup sec, comme pour couper (\zh{受贿} shòuhuì, deux coups).
\item Ton neutre : petit point rapide du doigt, sans mélodie (\zh{呢} ne).
\end{enumerate}
Répétez chaque mot trois fois avec son geste, puis sans le geste : la mélodie reste.
\end{methode}

\begin{savaistu}[Les écoliers chinois font aussi des gestes de tons]
En Chine, les enfants de l'école primaire apprennent la prononciation des caractères en traçant les tons avec la main, exactement comme vous venez de le faire. Le professeur écrit le caractère au tableau, et toute la classe récite en chœur en dessinant la mélodie dans l'air. Même pour un locuteur natif, le ton fait partie du mot et s'apprend avec le corps : vous utilisez donc une méthode authentiquement chinoise.
\end{savaistu}

\courssub{Exercices de discrimination et répétition en chaîne}

\begin{exoresolu}[Discrimination des tons : lève la bonne carte]
Énoncé. Chaque élève prépare cinq cartes : 1, 2, 3, 4, 0 (ton neutre). Le professeur lit à voix haute les mots suivants, dans le désordre : \zh{贪污}, \zh{腐败}, \zh{诚实}, \zh{受贿}, \zh{你呢}, \zh{但是}. Pour chaque mot, levez les deux cartes correspondant aux tons entendus.
\tcblower
Solution détaillée.
\begin{itemize}
\item \zh{贪污} tānwū : cartes 1 et 1 — deux tons hauts et plats.
\item \zh{腐败} fǔbài : cartes 3 et 4 — creux puis chute.
\item \zh{诚实} chéngshí : cartes 2 et 2 — deux montées.
\item \zh{受贿} shòuhuì : cartes 4 et 4 — deux chutes sèches.
\item \zh{你呢} nǐ ne : cartes 3 et 0 — creux puis syllabe légère sans mélodie.
\item \zh{但是} dànshì : cartes 4 et 4 — deux chutes.
\end{itemize}
Piège : pour \zh{你呢}, beaucoup d'élèves lèvent la carte 2 pour \zh{呢} parce qu'une question française « monte ». Or \zh{呢} est neutre : carte 0.
\end{exoresolu}

\begin{exercice}[Répétition en chaîne avec gestuelle]
En rangées, répétez en chaîne les phrases clés du dialogue, chacun une phrase, avec le geste des tons sur chaque syllabe :
\begin{enumerate}
\item \zh{我觉得腐败是一个大问题。} Wǒ juéde fǔbài shì yí ge dà wèntí.
\item \zh{你同意吗？} Nǐ tóngyì ma ?
\item \zh{我同意，但是法律也很重要。} Wǒ tóngyì, dànshì fǎlǜ yě hěn zhòngyào.
\item \zh{你说得对。你呢？} Nǐ shuō de duì. Nǐ ne ?
\end{enumerate}
Le voisin écoute et signale tout ton faux avant de prendre la relève.
\end{exercice}

\begin{corrige}[Exercice de la répétition en chaîne]
Points de vigilance attendus :
\begin{itemize}
\item Phrase 1 : \zh{腐败} fǔbài en 3 + 4 (creux puis chute) ; \zh{得} de \zh{觉得} reste neutre.
\item Phrase 2 : \zh{同意} garde ses tons 2 + 4 ; \zh{吗} neutre, la voix ne monte pas en fin de question.
\item Phrase 3 : \zh{但是} dànshì en 4 + 4, deux chutes franches ; \zh{法律} fǎlǜ en 3 + 4 ; \zh{很} hěn devant \zh{重} zhòng (4e ton) garde son 3e ton, pas de sandhi ici.
\item Phrase 4 : \zh{得} de \zh{说得对} est neutre ; \zh{呢} de \zh{你呢} est neutre et léger.
\end{itemize}
\end{corrige}

\begin{aretenir}[L'essentiel de la prononciation pour votre débat]
\begin{itemize}
\item Le ton fait partie du mot : \zh{腐败} fǔbài mal prononcé devient incompréhensible.
\item Deux 3e tons qui se suivent : le premier devient un 2e ton (\emph{yíngxiǎng}, \emph{hén hǎo}).
\item Les particules \zh{吗}, \zh{呢}, \zh{的}, \zh{得} sont neutres : courtes et légères.
\item Deux 4e tons successifs : deux chutes complètes, on remonte la voix entre les deux.
\item La question avec \zh{吗} ne monte pas : les tons des mots ne changent jamais.
\item La main qui trace la courbe du ton est votre meilleur outil de mémorisation.
\end{itemize}
\end{aretenir}

Votre prononciation est maintenant prête. Il ne reste plus qu'à la mettre en action : passons aux jeux de rôle et au débat guidé sur la corruption.

\coursec{Jeux de rôle et débat guidé}

Dans la section précédente, vous avez travaillé les tons et la prononciation sur les phrases clés du thème de la corruption. Il est maintenant temps de mettre cette prononciation au service d'une vraie communication : jouer des rôles et débattre en chinois. C'est l'objectif final de cette leçon : parler, échanger, argumenter — en phrases courtes, mais correctes et bien prononcées.

\courssub{Jeu de rôle 1 : un journaliste interviewe un élève}

\textbf{Situation.} Un journaliste d'une chaîne de télévision vient dans votre lycée à Yaoundé. Il interviewe un élève sur la corruption à l'école (par exemple : payer pour avoir de bonnes notes, copier aux examens).

\textbf{Répartition des rôles.} Par groupes de deux : l'élève A joue le journaliste, l'élève B joue l'élève interviewé. Puis on échange les rôles.

\textbf{Formules à utiliser obligatoirement :}

\begin{center}
\begin{tabularx}{\linewidth}{|X|X|X|}\hline
\rowcolor{popblueL} Fonction & Formule chinoise + pinyin & Traduction \\ \hline
Poser la question d'opinion & \zh{你怎么看？} Nǐ zěnme kàn ? & Qu'en penses-tu ? \\ \hline
Donner son avis & \zh{我觉得……} Wǒ juéde... & Je pense que... / Je trouve que... \\ \hline
Donner un avis plus fort & \zh{我认为……} Wǒ rènwéi... & Je considère que... \\ \hline
Proposer une solution & \zh{我们应该……} Wǒmen yīnggāi... & Nous devrions... \\ \hline
\end{tabularx}
\end{center}

\begin{exemplebox}[Modèle de dialogue à imiter]
A : \zh{你好！请问，你怎么看学校里的腐败？}\\
\emph{Nǐ hǎo ! Qǐngwèn, nǐ zěnme kàn xuéxiào lǐ de fǔbài ?}\\
« Bonjour ! S'il vous plaît, que pensez-vous de la corruption à l'école ? »

B : \zh{我觉得腐败很不好。学生不应该付钱买好分数。}\\
\emph{Wǒ juéde fǔbài hěn bù hǎo. Xuésheng bù yīnggāi fù qián mǎi hǎo fēnshù.}\\
« Je pense que la corruption est très mauvaise. Les élèves ne devraient pas payer pour acheter de bonnes notes. »

A : \zh{那我们应该怎么做？}\\
\emph{Nà wǒmen yīnggāi zěnme zuò ?}\\
« Alors, que devrions-nous faire ? »

B : \zh{我们应该诚实，也应该告诉老师。}\\
\emph{Wǒmen yīnggāi chéngshí, yě yīnggāi gàosu lǎoshī.}\\
« Nous devrions être honnêtes, et nous devrions aussi le dire au professeur. »
\end{exemplebox}

\begin{attention}[Erreur fréquente : la place de 怎么]
Les francophones traduisent souvent mot à mot « comment vois-tu ? » et produisent des phrases incorrectes. En chinois, \zh{怎么} (zěnme, « comment ») se place \textbf{avant le verbe}, jamais en fin de phrase : \zh{你怎么看？} et non *\zh{你看怎么？}. De même : \zh{应该怎么做？} (yīnggāi zěnme zuò ?) = « comment devrait-on faire ? », c'est-à-dire « que devons-nous faire ? ».
\end{attention}

\courssub{Jeu de rôle 2 : deux amis face à un pot-de-vin}

\textbf{Situation.} Deux amis, Adama et Lin, discutent : un fonctionnaire demande de l'argent à leur oncle pour signer un document. Que faut-il faire ?

\begin{definition}[Lexique indispensable du jeu de rôle 2]
\begin{center}
\begin{tabularx}{\linewidth}{|X|X|X|X|}\hline
\rowcolor{popblueL} Caractères & Pinyin & Nature & Traduction \\ \hline
\zh{行贿} & xínghuì & verbe & donner un pot-de-vin, corrompre \\ \hline
\zh{受贿} & shòuhuì & verbe & recevoir un pot-de-vin, se laisser corrompre \\ \hline
\zh{腐败} & fǔbài & nom / adj. & corruption ; corrompu \\ \hline
\zh{诚实} & chéngshí & adj. & honnête \\ \hline
\zh{反对} & fǎnduì & verbe & s'opposer à, combattre \\ \hline
\zh{应该} & yīnggāi & verbe modal & devoir (il faut) \\ \hline
\zh{拒绝} & jùjué & verbe & refuser \\ \hline
\zh{举报} & jǔbào & verbe & dénoncer, signaler \\ \hline
\end{tabularx}
\end{center}
\end{definition}

\begin{exemplebox}[Modèle de dialogue à imiter]
A : \zh{那个工作人员要钱，怎么办？}\\
\emph{Nà ge gōngzuò rényuán yào qián, zěnme bàn ?}\\
« Ce fonctionnaire demande de l'argent, que faire ? »

B : \zh{我认为我们不应该行贿。行贿和受贿都不对。}\\
\emph{Wǒ rènwéi wǒmen bù yīnggāi xínghuì. Xínghuì hé shòuhuì dōu bù duì.}\\
« Je pense que nous ne devrions pas donner de pot-de-vin. Corrompre et se laisser corrompre, les deux sont mal. »

A : \zh{那应该怎么办？}\\
\emph{Nà yīnggāi zěnme bàn ?}\\
« Alors que faut-il faire ? »

B : \zh{我们应该拒绝，也应该举报。}\\
\emph{Wǒmen yīnggāi jùjué, yě yīnggāi jǔbào.}\\
« Nous devrions refuser, et nous devrions aussi le dénoncer. »
\end{exemplebox}

\begin{textebox}[Complément utile : l'hypothèse avec 如果……会……]
A : \zh{如果你是警察，你会怎么做？}\\
\emph{Rúguǒ nǐ shì jǐngchá, nǐ huì zěnme zuò ?}\\
« Si tu étais policier, que ferais-tu ? »

B : \zh{我会反对腐败。}\\
\emph{Wǒ huì fǎnduì fǔbài.}\\
« Je combattrais la corruption. »
\end{textebox}

\begin{propriete}[La structure hypothétique 如果……（就）会……]
\textbf{Structure :} 如果 + condition， + sujet + （就）会 + verbe.

\textbf{Emploi :} \zh{如果} (rúguǒ, « si ») introduit la condition ; \zh{会} (huì) marque ici le \emph{futur hypothétique} ou le \emph{conditionnel} (« ferais »). Le mot \zh{就} (jiù, « alors ») est facultatif mais très fréquent.

\textbf{Comparaison avec le français :} le français exige l'imparfait dans la subordonnée et le conditionnel dans la principale (« si tu \emph{étais}, tu \emph{ferais} »). Le chinois n'a pas de conjugaison : le verbe reste invariable, seuls les mots-outils \zh{如果} et \zh{会} portent le sens.

Exemples :
\begin{itemize}
\item \zh{如果我有时间，我会学习汉语。} \emph{Rúguǒ wǒ yǒu shíjiān, wǒ huì xuéxí Hànyǔ.} « Si j'avais du temps, j'étudierais le chinois. »
\item \zh{如果每个人都诚实，社会就会更好。} \emph{Rúguǒ měi ge rén dōu chéngshí, shèhuì jiù huì gèng hǎo.} « Si chacun était honnête, la société serait meilleure. »
\end{itemize}
\end{propriete}

\courssub{Le débat guidé : organisation et rôles}

\textbf{Sujet du débat :} \zh{腐败是我们社会最大的问题吗？} \emph{Fǔbài shì wǒmen shèhuì zuì dà de wèntí ma ?} « La corruption est-elle le plus grand problème de notre société ? »

La classe est divisée en deux équipes : l'équipe \textbf{POUR} (oui, c'est le plus grand problème) et l'équipe \textbf{CONTRE} (non, il y a des problèmes plus importants : la pauvreté, le chômage, les maladies…). Quatre rôles sont distribués :

\begin{itemize}
\item \textbf{L'animateur} (主持人 zhǔchírén) : il pose les questions et relance la parole avec \zh{你呢？} (Nǐ ne ? « Et toi ? ») et \zh{为什么？} (Wèishénme ? « Pourquoi ? »).
\item \textbf{Les orateurs} : ils donnent leur avis avec \zh{我觉得……} ou \zh{我认为……}, puis un argument avec \zh{因为……} (yīnwèi, « parce que »).
\item \textbf{L'équipe adverse} répond avec \zh{我不同意，但是……} (Wǒ bù tóngyì, dànshì... « Je ne suis pas d'accord, mais... »).
\item \textbf{Le rapporteur} (报告人 bàogàorén) : il conclut le débat avec \zh{所以……} (suǒyǐ, « donc, c'est pourquoi »).
\end{itemize}

\begin{popfigure}
\begin{tikzpicture}[
  node distance=0.55cm and 1.1cm,
  boite/.style={rectangle, rounded corners=3pt, draw=popdark, thick, align=center, font=\small, inner sep=5pt},
  fleche/.style={-{Stealth[length=2.5mm]}, thick, popdark}
]
\node[boite, fill=popgoldL, align=center] (anim){Animateur\\ \zh{你呢？为什么？}\\ \emph{Nǐ ne ? Wèishénme ?}};
\node[boite, fill=popgreenL, below left=0.9cm and 0.2cm of anim, align=center] (pour){Équipe POUR\\ \zh{我认为……因为……}\\ \emph{Wǒ rènwéi... yīnwèi...}};
\node[boite, fill=poppinkL, below right=0.9cm and 0.2cm of anim, align=center] (contre){Équipe CONTRE\\ \zh{我不同意，但是……}\\ \emph{Wǒ bù tóngyì, dànshì...}};
\node[boite, fill=popblueL, below=2.6cm of anim, align=center] (rapp){Rapporteur\\ \zh{所以……}\\ \emph{Suǒyǐ...}};
\draw[fleche] (anim) -- (pour);
\draw[fleche] (anim) -- (contre);
\draw[fleche] (pour.east) -- (contre.west) node[midway, above, font=\footnotesize\itshape] {réponse};
\draw[fleche] (pour.south) |- ($(rapp.west)+(0,0.15)$) -- (rapp.west);
\draw[fleche] (contre.south) |- ($(rapp.east)+(0,0.15)$) -- (rapp.east);
\end{tikzpicture}
\legende{Organigramme du débat guidé : circulation de la parole et formules associées à chaque rôle.}
\end{popfigure}

\begin{methode}[Réussir sa prise de parole en trois phrases]
Avant de prendre la parole, prépare mentalement (ou sur une fiche) exactement trois phrases :
\begin{enumerate}
\item \textbf{Opinion} : \zh{我觉得腐败是一个很大的问题。} \emph{Wǒ juéde fǔbài shì yí ge hěn dà de wèntí.} « Je pense que la corruption est un très grand problème. »
\item \textbf{Argument} : \zh{因为腐败让穷人更穷。} \emph{Yīnwèi fǔbài ràng qióngrén gèng qióng.} « Parce que la corruption rend les pauvres encore plus pauvres. »
\item \textbf{Exemple} : \zh{比如，有人付钱买好分数。} \emph{Bǐrú, yǒu rén fù qián mǎi hǎo fēnshù.} « Par exemple, certains paient pour acheter de bonnes notes. »
\end{enumerate}
Puis prends la parole en regardant tes camarades, pas ta fiche.
\end{methode}

\begin{attention}[Ne lis pas tes notes mot à mot]
L'erreur classique des débutants : écrire tout son texte et le lire. Résultat : les tons disparaissent, le débit devient incompréhensible, et l'on ne peut plus répondre aux relances \zh{你呢？为什么？}. La bonne méthode : \textbf{mémorise les formules fonctionnelles} (\zh{我觉得}、\zh{我认为}、\zh{因为}、\zh{所以}、\zh{我不同意，但是}) et \textbf{improvise le contenu} avec le lexique de la leçon (\zh{腐败}、\zh{行贿}、\zh{受贿}、\zh{诚实}、\zh{应该}). Une phrase courte et bien prononcée vaut mieux qu'un long texte lu.
\end{attention}

\courssub{Grille d'observation et évaluation}

Pendant le débat, deux élèves observateurs (ou le professeur) remplissent la grille suivante pour chaque orateur :

\begin{center}
\begin{tabularx}{\linewidth}{|X|X|X|X|}\hline
\rowcolor{popblueL} Critère & À observer & Bien & À améliorer \\ \hline
Prononciation des tons & Les tons de \zh{腐败} (fǔbài, 3+4), \zh{诚实} (chéngshí, 2+2), \zh{我觉得} (Wǒ juéde) sont-ils nets ? & & \\ \hline
Formules fonctionnelles & L'orateur utilise-t-il \zh{我觉得} / \zh{我认为} / \zh{因为} / \zh{所以} ? & & \\ \hline
Participation & Prend-il la parole sans être relancé ? & & \\ \hline
Respect de la parole & Écoute-t-il sans couper ? Répond-il à l'argument de l'autre ? & & \\ \hline
\end{tabularx}
\end{center}

\begin{exoresolu}[Mini-exercice : compléter la conclusion du rapporteur]
Énoncé. Le rapporteur doit conclure le débat. Complète sa phrase avec \zh{所以} : « Les deux équipes ont donné de bons arguments ; donc, chacun doit être honnête. »

\tcblower

Solution détaillée. On construit la phrase en deux parties reliées par \zh{所以} (suǒyǐ, « donc »), placé \textbf{en tête de la deuxième proposition} (et non à la fin, comme parfois en français parlé) :

\zh{两个队都说了很好的理由，所以每个人都应该诚实。}

\emph{Liǎng ge duì dōu shuō le hěn hǎo de lǐyóu, suǒyǐ měi ge rén dōu yīnggāi chéngshí.}

« Les deux équipes ont donné de très bonnes raisons, donc chacun doit être honnête. »

Points de grammaire vérifiés : classificateur \zh{个} devant \zh{队} ; \zh{都} placé après le sujet et avant le verbe ; modal \zh{应该} avant le verbe \zh{诚实} (adjectif employé comme prédicat après un modal, sans \zh{是}).
\end{exoresolu}

\begin{exemplebox}[Phrases à réutiliser dans n'importe quel débat]
\zh{我们应该怎么做？} \emph{Wǒmen yīnggāi zěnme zuò ?} « Que devons-nous faire ? »

\zh{每个人都应该诚实。} \emph{Měi ge rén dōu yīnggāi chéngshí.} « Chacun doit être honnête. »

\zh{我不同意，但是你的理由很好。} \emph{Wǒ bù tóngyì, dànshì nǐ de lǐyóu hěn hǎo.} « Je ne suis pas d'accord, mais ton argument est bon. »
\end{exemplebox}

\courssub{Complément utile à l'examen : comparer le Cameroun et la Chine}

À l'examen, on peut vous demander de produire une phrase de comparaison. Retenez ce modèle simple :

\zh{在中国，政府反腐败。} \emph{Zài Zhōngguó, zhèngfǔ fǎnfǔbài.} « En Chine, le gouvernement lutte contre la corruption. »

\zh{在喀麦隆，政府也反腐败。} \emph{Zài Kāmàilóng, zhèngfǔ yě fǎnfǔbài.} « Au Cameroun, le gouvernement lutte aussi contre la corruption. »

Structure : \zh{在} + lieu， + sujet + （也）+ verbe. Le mot \zh{也} (yě, « aussi ») se place \textbf{avant le verbe}, jamais en fin de phrase comme le français « aussi ».

\begin{savaistu}[Lutte contre la corruption : Cameroun et Chine]
Le Cameroun dispose d'une agence spécialisée dans la lutte contre la corruption, la CONAC (Commission nationale anti-corruption), créée en 2006. La Chine, de son côté, mène depuis 2012 une vaste campagne anticorruption appelée \zh{反腐败} (fǎnfǔbài, littéralement « s'opposer à la corruption »), très médiatisée : le mot \zh{反} signifie « contre, anti- » et se retrouve dans de nombreux mots utiles. Dans les deux pays, l'honnêteté — \zh{诚实} — est présentée comme une valeur essentielle de l'éducation des jeunes citoyens.
\end{savaistu}

\begin{experience}[Le grand débat de la classe (20 minutes)]
\begin{enumerate}
\item Le professeur écrit le sujet au tableau : \zh{腐败是我们社会最大的问题吗？}
\item Deux équipes de quatre à six élèves se forment ; chaque équipe prépare cinq minutes : trois phrases par orateur (opinion + argument + exemple).
\item L'animateur ouvre le débat : \zh{你怎么看？} Chaque orateur parle trente secondes maximum.
\item L'équipe adverse répond : \zh{我不同意，但是……}
\item Le rapporteur conclut : \zh{所以……}
\item Les observateurs remplissent la grille et annoncent l'équipe la plus convaincante — en chinois si possible : \zh{我觉得这个队说得很好！} \emph{Wǒ juéde zhè ge duì shuō de hěn hǎo !}
\end{enumerate}
\end{experience}

\newpage

\coursec{Activité d'intégration}

\courssub{Objectifs de l'activité}

À l'issue de cette activité d'intégration, tu seras capable de :
\begin{itemize}
\item mobiliser le \cle{lexique de la corruption} (\zh{腐败}, \zh{贪污}, \zh{受贿}, \zh{法律}, \zh{诚实}…) dans un échange oral spontané ;
\item exprimer ton \cle{opinion} avec \zh{我觉得} et \zh{我认为}, et la justifier avec \zh{首先}…\zh{然后}… ;
\item marquer l'\cle{accord} et le \cle{désaccord} de façon polie (\zh{我同意你的看法} / \zh{我不太同意}) ;
\item \cle{relancer} la parole d'un camarade avec \zh{你呢？}, \zh{为什么？}, \zh{你能举个例子吗？} ;
\item conclure un débat au nom d'une équipe avec \zh{所以，我们认为……} ;
\item prononcer correctement les \cle{tons}, en particulier les suites de troisièmes tons et le ton neutre.
\end{itemize}

\courssub{Situation}

Dans le cadre de la semaine de la citoyenneté organisée par ton lycée à Yaoundé, le club de chinois participe à un grand débat interclasses. Le thème retenu cette année est la lutte contre la corruption. Ta classe de Première A doit débattre en chinois devant un jury composé du professeur et de deux élèves de Terminale. La motion est la suivante :

\begin{textebox}[La motion du débat]
腐败是我们社会最大的问题吗？\\
\emph{Fǔbài shì wǒmen shèhuì zuì dà de wèntí ma ?}\\
« La corruption est-elle le plus grand problème de notre société ? »\\[4pt]
Équipe A (pour) : 是，腐败是最大的问题。\\
\emph{Shì, fǔbài shì zuì dà de wèntí.}\\
« Oui, la corruption est le plus grand problème. »\\[4pt]
Équipe B (contre) : 不是，还有更大的问题，比如贫困和失业。\\
\emph{Bú shì, hái yǒu gèng dà de wèntí, bǐrú pínkùn hé shīyè.}\\
« Non, il y a de plus grands problèmes, par exemple la pauvreté et le chômage. »
\end{textebox}

Chaque équipe dispose de dix minutes de préparation, puis le débat proprement dit dure dix minutes. L'animateur de chaque équipe distribue la parole, le rapporteur prend des notes et conclut. Le jury écoute et remplit une grille d'observation : prononciation des tons, emploi des formules d'opinion, participation de tous les membres.

\courssub{Dialogue modèle : donner son avis et réagir}

Avant le grand débat, entraîne-toi avec ce court dialogue entre deux élèves, Awa (A) et Marc (B). Lis-le à voix haute avec un camarade, en veillant aux tons.

\begin{textebox}[Dialogue modèle]
A：我觉得腐败是我们社会最大的问题。你呢？\\
\emph{Wǒ juéde fǔbài shì wǒmen shèhuì zuì dà de wèntí. Nǐ ne ?}\\
« Je pense que la corruption est le plus grand problème de notre société. Et toi ? »\\[4pt]
B：我不太同意。我认为贫困是更大的问题。\\
\emph{Wǒ bú tài tóngyì. Wǒ rènwéi pínkùn shì gèng dà de wèntí.}\\
« Je ne suis pas tout à fait d'accord. Je pense que la pauvreté est un problème plus grand. »\\[4pt]
A：为什么？你能举个例子吗？\\
\emph{Wèishénme ? Nǐ néng jǔ ge lìzi ma ?}\\
« Pourquoi ? Peux-tu donner un exemple ? »\\[4pt]
B：比如，很多人没有钱看病，也没有工作。\\
\emph{Bǐrú, hěn duō rén méiyǒu qián kànbìng, yě méiyǒu gōngzuò.}\\
« Par exemple, beaucoup de gens n'ont pas d'argent pour se soigner et n'ont pas de travail. »\\[4pt]
A：你说得对，但是腐败让贫困更严重。\\
\emph{Nǐ shuō de duì, dànshì fǔbài ràng pínkùn gèng yánzhòng.}\\
« Tu as raison, mais la corruption rend la pauvreté encore plus grave. »\\[4pt]
B：嗯，我同意你的看法：这两个问题都很重要。\\
\emph{Ng, wǒ tóngyì nǐ de kànfǎ : zhè liǎng ge wèntí dōu hěn zhòngyào.}\\
« Hum, je suis d'accord avec toi : ces deux problèmes sont tous les deux très importants. »
\end{textebox}

Observe : A donne son avis avec \zh{我觉得}, B nuance son désaccord avec \zh{我不太同意}, A relance avec deux questions, et le dialogue se termine par un accord partiel. C'est exactement la dynamique attendue dans ton débat.

\courssub{Consignes}

\begin{enumerate}
\item \textbf{Compréhension.} Lis la motion à voix haute avec ton groupe. Souligne les mots du lexique de la corruption que tu connais déjà. Vérifie le sens de \zh{社会} (shèhuì, « société ») et de \zh{问题} (wèntí, « problème »).
\item \textbf{Préparation individuelle.} Rédige sur ta feuille trois phrases en chinois : une opinion avec \zh{我觉得} ou \zh{我认为}, un argument introduit par \zh{首先} puis \zh{然后}, et un exemple concret (école, marché, hôpital, football…). Chaque phrase doit être courte et correcte.
\item \textbf{Répartition des rôles.} Dans chaque équipe, choisissez un \cle{animateur} (il relance avec \zh{你呢？}, \zh{为什么？}, \zh{你能举个例子吗？}) et un \cle{rapporteur} (il note les arguments et prépare la conclusion avec \zh{所以，我们认为……}).
\item \textbf{Entraînement des tons.} Répétez à voix haute les mots difficiles : \zh{腐败} (fǔbài, 3\up{e}+4\up{e} ton), \zh{我觉得} (wǒ juéde, 3\up{e}+2\up{e}+neutre), \zh{我认为} (wǒ rènwéi, 3\up{e}+4\up{e}+2\up{e} ton). Attention : deux troisièmes tons qui se suivent se prononcent 2\up{e}+3\up{e} ton (\zh{wǒ hěn} se dit presque \emph{wó hěn}).
\item \textbf{Débat.} Pendant dix minutes, chaque élève prend la parole au moins deux fois : une fois pour donner son opinion, une fois pour réagir à un camarade (accord ou désaccord poli). L'animateur veille à ce que tout le monde parle.
\item \textbf{Conclusion.} Le rapporteur de chaque équipe résume les arguments en deux ou trois phrases commençant par \zh{所以，我们认为……}.
\item \textbf{Prolongement.} À la maison, enregistre sur ton téléphone un mini-exposé de cinq phrases commençant par \zh{大家好！今天我要谈一谈腐败的问题。} (\emph{Dàjiā hǎo ! Jīntiān wǒ yào tán yi tán fǔbài de wèntí.} — « Bonjour à tous ! Aujourd'hui, je vais parler du problème de la corruption. »).
\end{enumerate}

\courssub{Ressources à mobiliser}

\begin{center}
\begin{tabularx}{\linewidth}{|X|X|X|X|}
\hline
\rowcolor{popblueL} Caractères & Pinyin & Nature & Traduction \\
\hline
腐败 & fǔbài & nom / adj. & corruption ; corrompu \\
\hline
贪污 & tānwū & verbe & détourner (de l'argent) \\
\hline
受贿 & shòuhuì & verbe & accepter des pots-de-vin \\
\hline
法律 & fǎlǜ & nom & loi, droit \\
\hline
诚实 & chéngshí & adj. & honnête \\
\hline
贫困 & pínkùn & nom / adj. & pauvreté ; pauvre \\
\hline
失业 & shīyè & verbe / nom & perdre son emploi ; chômage \\
\hline
解决 & jiějué & verbe & résoudre \\
\hline
严重 & yánzhòng & adj. & grave, sérieux \\
\hline
反对 & fǎnduì & verbe & s'opposer à, lutter contre \\
\hline
例子 & lìzi & nom & exemple \\
\hline
看法 & kànfǎ & nom & point de vue, opinion \\
\hline
\end{tabularx}
\end{center}

\begin{aretenir}[Boîte à outils du débat]
\begin{itemize}
\item Donner son avis : \zh{我觉得……} (\emph{wǒ juéde…}) / \zh{我认为……} (\emph{wǒ rènwéi…}) « je pense que… »
\item Ordonner ses arguments : \zh{首先……然后……} (\emph{shǒuxiān… ránhòu…}) « d'abord… ensuite… »
\item Accord : \zh{我同意你的看法。} (\emph{Wǒ tóngyì nǐ de kànfǎ.}) « Je suis d'accord avec ton point de vue. » ; \zh{你说得对。} (\emph{Nǐ shuō de duì.}) « Tu as raison. »
\item Désaccord poli : \zh{我不太同意。} (\emph{Wǒ bú tài tóngyì.}) « Je ne suis pas tout à fait d'accord. »
\item Relancer : \zh{你呢？} (\emph{Nǐ ne ?}) « Et toi ? » ; \zh{为什么？} (\emph{Wèishénme ?}) « Pourquoi ? » ; \zh{你能举个例子吗？} (\emph{Nǐ néng jǔ ge lìzi ma ?}) « Peux-tu donner un exemple ? »
\item Conclure : \zh{所以，我们认为……} (\emph{Suǒyǐ, wǒmen rènwéi…}) « Donc, nous pensons que… »
\end{itemize}
\end{aretenir}

\begin{attention}[Pièges de prononciation et de formulation]
\begin{itemize}
\item \zh{觉得} se prononce \emph{juéde} : le second caractère perd son ton (ton neutre). Ne dis pas \emph{juédé}.
\item Après \zh{我觉得} et \zh{我认为}, on met directement une phrase complète, sans \zh{是} devant un adjectif : \zh{我觉得腐败很严重。} et non *\zh{我觉得腐败是很严重}.
\item Pour nuancer le désaccord, préfère \zh{我不太同意} à \zh{我不同意}, qui sonne trop brutal dans un débat scolaire.
\item Attention à l'ordre des mots : le complément de temps ou de lieu se place avant le verbe : \zh{在我们学校，……} (\emph{Zài wǒmen xuéxiào,…}) « Dans notre école, … ».
\item \zh{不} devient \emph{bú} devant un quatrième ton : \zh{不是} se prononce \emph{bú shì}, \zh{不太同意} se prononce \emph{bú tài tóngyì}.
\end{itemize}
\end{attention}

\courssub{Proposition de production et corrigé commenté}

\begin{exoresolu}[Mini-exposé modèle (prolongement enregistré)]
Énoncé : enregistre un mini-exposé de cinq phrases sur le thème du débat, en commençant par la formule d'ouverture imposée.
\tcblower
\textbf{Production modèle :}\\[3pt]
大家好！今天我要谈一谈腐败的问题。\\
\emph{Dàjiā hǎo ! Jīntiān wǒ yào tán yi tán fǔbài de wèntí.}\\
« Bonjour à tous ! Aujourd'hui, je vais parler du problème de la corruption. »\\[3pt]
我认为腐败是我们社会一个很大的问题。\\
\emph{Wǒ rènwéi fǔbài shì wǒmen shèhuì yí ge hěn dà de wèntí.}\\
« Je pense que la corruption est un très grand problème de notre société. »\\[3pt]
首先，腐败让穷人更穷；然后，它破坏我们的法律。\\
\emph{Shǒuxiān, fǔbài ràng qióngrén gèng qióng ; ránhòu, tā pòhuài wǒmen de fǎlǜ.}\\
« D'abord, la corruption rend les pauvres encore plus pauvres ; ensuite, elle détruit nos lois. »\\[3pt]
比如，在一些地方，没有钱就看不到医生。\\
\emph{Bǐrú, zài yìxiē dìfang, méiyǒu qián jiù kàn bu dào yīshēng.}\\
« Par exemple, dans certains endroits, sans argent on ne peut pas voir de médecin. »\\[3pt]
所以，我认为我们应该诚实，一起反对腐败。谢谢大家！\\
\emph{Suǒyǐ, wǒ rènwéi wǒmen yīnggāi chéngshí, yìqǐ fǎnduì fǔbài. Xièxie dàjiā !}\\
« Donc, je pense que nous devons être honnêtes et lutter ensemble contre la corruption. Merci à tous ! »\\[6pt]
\textbf{Commentaire des choix de langue :}
\begin{itemize}
\item La formule d'ouverture \zh{大家好！} capte l'attention ; \zh{谈一谈} (redoublement du verbe) adoucit le ton : « parler un peu de ».
\item L'opinion est introduite par \zh{我认为}, plus soutenu que \zh{我觉得}, ce qui convient à un exposé.
\item La paire \zh{首先……然后……} structure clairement les deux arguments ; \zh{让} (ràng, « faire que ») introduit une conséquence : structure Sujet + 让 + personne + adjectif.
\item L'exemple utilise \zh{比如} (« par exemple ») et la structure \zh{没有……就……} (« sans… on ne peut pas… »), très fréquente à l'oral ; \zh{看不到} (kàn bu dào) est le potentiel négatif de \zh{看到} : « ne pas arriver à voir ».
\item La conclusion reprend la formule du débat \zh{所以，我认为……} et se termine par la politesse \zh{谢谢大家！}.
\item Tons à surveiller : \zh{fǔbài} (3+4), \zh{hěn dà} (3+4), \zh{yìqǐ} (4+3). Enchaîne les phrases lentement, en marquant bien les quatrièmes tons descendants.
\end{itemize}
\end{exoresolu}

\courssub{Bilan de l'activité}

Cette activité t'a permis de mobiliser en situation réelle l'ensemble des acquis de la leçon : le lexique de la corruption, les formules d'opinion (\zh{我觉得}, \zh{我认为}), les marqueurs d'organisation (\zh{首先}, \zh{然后}, \zh{所以}), les expressions d'accord et de désaccord poli, et les questions de relance (\zh{你呢？}, \zh{为什么？}). Tu as également travaillé la prononciation des tons dans un contexte de communication authentique, devant un public. La grille d'observation du jury (tons, formules, participation) t'aide à identifier tes points forts et tes progrès à faire. Enfin, le mini-exposé enregistré en prolongement te permet de t'auto-évaluer : réécoute-toi, vérifie tes tons et réenregistre si nécessaire. Ces compétences de prise de parole en public te seront utiles pour les débats du programme de Terminale et pour l'épreuve orale.

\newpage

\coursec{Méthodes et exercices résolus}

\courssub{Lexique du thème : la corruption}

Avant de parler, il faut connaître les mots du débat. Voici le vocabulaire indispensable de cette leçon, à apprendre avec les tons.

\begin{center}
\begin{tabularx}{\linewidth}{|X|X|X|X|}
\hline
\rowcolor{popblueL} Caractères & Pinyin & Nature & Traduction \\
\hline
腐败 & fǔbài & nom / adj. & corruption ; corrompu \\
\hline
贪污 & tānwū & verbe & détourner (de l'argent), se corrompre \\
\hline
诚实 & chéngshí & adjectif & honnête \\
\hline
法律 & fǎlǜ & nom & loi, droit \\
\hline
反对 & fǎnduì & verbe & s'opposer à, combattre \\
\hline
接受 & jiēshòu & verbe & accepter \\
\hline
严重 & yánzhòng & adjectif & grave, sérieux \\
\hline
重要 & zhòngyào & adjectif & important \\
\hline
教育 & jiàoyù & nom & éducation \\
\hline
钱 & qián & nom & argent \\
\hline
老百姓 & lǎobǎixìng & nom & le peuple, les gens ordinaires \\
\hline
政府 & zhèngfǔ & nom & gouvernement, administration \\
\hline
警察 & jǐngchá & nom & police, policier \\
\hline
问题 & wèntí & nom & problème, question \\
\hline
看法 & kànfǎ & nom & point de vue, opinion \\
\hline
\end{tabularx}
\end{center}

\begin{attention}[Homophones et quasi-homophones à ne pas confondre]
\begin{itemize}
\item \zh{诚实} \emph{chéngshí} « honnête » et \zh{城市} \emph{chéngshì} « ville » : seul le second ton change (2\textsuperscript{e} ton contre 4\textsuperscript{e} ton). À l'oral, l'erreur est très visible.
\item \zh{腐败} \emph{fǔbài} « corruption » et \zh{府} \emph{fǔ} « gouvernement, résidence officielle » (comme dans \zh{政府}) : même syllabe \emph{fǔ}, caractères et sens différents.
\item \zh{问题} \emph{wèntí} « problème » : ne pas prononcer \emph{wèntì} ; le deuxième caractère est au 2\textsuperscript{e} ton.
\end{itemize}
\end{attention}

\begin{savaistu}[La lutte contre la corruption au Cameroun et en Chine]
Au Cameroun comme en Chine, la lutte contre la corruption est un sujet de société majeur. En Chine, une grande campagne nationale contre la corruption (\zh{反腐败} \emph{fǎn fǔbài}) a été lancée dans les années 2010 ; on y parle souvent de « mouches » (\zh{苍蝇} cāngying, petits fonctionnaires corrompus) et de « tigres » (\zh{老虎} lǎohǔ, hauts responsables). Au Cameroun, des structures comme la CONAC (Commission nationale anti-corruption) sensibilisent les citoyens, et le mot d'ordre « nul n'est au-dessus de la loi » se retrouve dans les deux pays. En chinois, on peut dire : \zh{法律面前，人人平等。} \emph{Fǎlǜ miànqián, rénrén píngděng.} « Devant la loi, tous sont égaux. »
\end{savaistu}

\courssub{Fiche méthode 1 : Présenter oralement un exposé simple en chinois}

\begin{methode}[Structurer un mini-exposé en trois temps]
Un exposé oral au lycée dure 1 à 2 minutes. Il suit toujours le même plan, facile à mémoriser :
\begin{enumerate}
\item \textbf{Introduction} : saluer et annoncer le sujet.\\
\zh{大家好！今天我要谈谈腐败问题。} \emph{Dàjiā hǎo! Jīntiān wǒ yào tántan fǔbài wèntí.} « Bonjour à tous ! Aujourd'hui, je vais parler du problème de la corruption. »
\item \textbf{Développement} : 2 ou 3 idées reliées par des mots de liaison simples : \zh{首先} \emph{shǒuxiān} « d'abord », \zh{然后} \emph{ránhòu} « ensuite », \zh{最后} \emph{zuìhòu} « enfin ». Une idée = une phrase courte Sujet + Verbe + Complément.
\item \textbf{Conclusion} : donner son avis et remercier.\\
\zh{我认为我们应该反对腐败。谢谢大家！} \emph{Wǒ rènwéi wǒmen yīnggāi fǎnduì fǔbài. Xièxie dàjiā!} « Je pense que nous devons lutter contre la corruption. Merci à tous ! »
\end{enumerate}
Réflexes : parler lentement, une phrase à la fois ; préparer des fiches avec les caractères et le pinyin ; ne jamais lire mot à mot, regarder le public entre deux phrases.
\end{methode}

\begin{exoresolu}[Un exposé à compléter et à présenter]
\textbf{Énoncé.} Voici le squelette d'un exposé. Complétez-le avec : \zh{首先}、\zh{最后}、\zh{我认为}、\zh{腐败}，puis présentez-le oralement.\\
\zh{大家好！今天我要谈谈\_\_\_\_问题。\_\_\_\_，腐败对国家不好。然后，腐败让老百姓的生活更困难。\_\_\_\_，\_\_\_\_我们应该诚实。谢谢大家！}
\tcblower
\textbf{Solution détaillée.}
\begin{enumerate}
\item Premier blanc : le sujet annoncé est la corruption, donc \zh{腐败} (fǔbài). La phrase \zh{今天我要谈谈腐败问题} suit la structure Sujet (\zh{我}) + Verbe (\zh{谈谈}, verbe redoublé pour adoucir) + Complément (\zh{腐败问题}).
\item Deuxième blanc : on ouvre le développement avec \zh{首先} « d'abord ». Attention : \zh{首先} se place en tête de phrase, avant le sujet.
\item Troisième blanc : après \zh{首先} et \zh{然后}, la série se termine par \zh{最后} « enfin ».
\item Quatrième blanc : pour introduire une opinion personnelle, on utilise \zh{我认为} « je pense que », suivi d'une proposition complète : \zh{我认为我们应该诚实} (Wǒ rènwéi wǒmen yīnggāi chéngshí). \zh{应该} est un verbe modal placé avant le verbe principal \zh{诚实}.
\end{enumerate}
\textbf{Texte complet à présenter :}\\
\zh{大家好！今天我要谈谈腐败问题。首先，腐败对国家不好。然后，腐败让老百姓的生活更困难。最后，我认为我们应该诚实。谢谢大家！}\\
\emph{Dàjiā hǎo! Jīntiān wǒ yào tántan fǔbài wèntí. Shǒuxiān, fǔbài duì guójiā bù hǎo. Ránhòu, fǔbài ràng lǎobǎixìng de shēnghuó gèng kùnnan. Zuìhòu, wǒ rènwéi wǒmen yīnggāi chéngshí. Xièxie dàjiā!}\\
« Bonjour à tous ! Aujourd'hui, je vais parler du problème de la corruption. D'abord, la corruption est mauvaise pour le pays. Ensuite, elle rend la vie des gens plus difficile. Enfin, je pense que nous devons être honnêtes. Merci à tous ! »
\end{exoresolu}

\courssub{Fiche méthode 2 : Échanger dans un petit débat}

\begin{methode}[Donner son avis, être d'accord, relancer]
Un débat guidé repose sur des formules toutes faites à mémoriser par cœur :
\begin{enumerate}
\item \textbf{Donner son avis} : \zh{我认为……} \emph{wǒ rènwéi} « je pense que » ; \zh{我觉得……} \emph{wǒ juéde} « je trouve que » ; \zh{在我看来，……} \emph{zài wǒ kànlái} « à mon avis ».
\item \textbf{Être d'accord} : \zh{我同意你的看法。} \emph{Wǒ tóngyì nǐ de kànfǎ.} « Je suis d'accord avec ton point de vue. » ; \zh{你说得对！} \emph{Nǐ shuō de duì!} « Tu as raison ! »
\item \textbf{N'être pas d'accord} (toujours poliment) : \zh{我不同意。} \emph{Wǒ bù tóngyì.} ; \zh{对不起，可是我觉得……} \emph{Duìbuqǐ, kěshì wǒ juéde...} « Pardon, mais je trouve que... »
\item \textbf{Relancer la parole} : \zh{你呢？} \emph{Nǐ ne?} « Et toi ? » ; \zh{你怎么看？} \emph{Nǐ zěnme kàn?} « Qu'en penses-tu ? » ; \zh{为什么？} \emph{Wèishénme?} « Pourquoi ? »
\item \textbf{Conclure} : \zh{我们都同意，……} \emph{Wǒmen dōu tóngyì,...} « Nous sommes tous d'accord pour dire que... »
\end{enumerate}
Réflexe d'examen : dans un jeu de rôle, utilise au moins une formule de chaque catégorie (avis, accord ou désaccord, relance) pour montrer que tu maîtrises l'interaction.
\end{methode}

\begin{textebox}[Dialogue modèle : à la sortie du lycée à Yaoundé]
\zh{甲：在我看来，腐败是一个严重的问题。你怎么看？}\\
\emph{Jiǎ: Zài wǒ kànlái, fǔbài shì yí ge yánzhòng de wèntí. Nǐ zěnme kàn?}\\
« A : À mon avis, la corruption est un problème grave. Qu'en penses-tu ? »\\
\zh{乙：我同意你的看法。腐败让老百姓的生活更困难。}\\
\emph{Yǐ: Wǒ tóngyì nǐ de kànfǎ. Fǔbài ràng lǎobǎixìng de shēnghuó gèng kùnnan.}\\
« B : Je suis d'accord avec toi. La corruption rend la vie des gens plus difficile. »\\
\zh{甲：我认为我们应该从小诚实。你呢？}\\
\emph{Jiǎ: Wǒ rènwéi wǒmen yīnggāi cóng xiǎo chéngshí. Nǐ ne?}\\
« A : Je pense que nous devons être honnêtes dès le plus jeune âge. Et toi ? »\\
\zh{乙：你说得对！诚实比钱重要。}\\
\emph{Yǐ: Nǐ shuō de duì! Chéngshí bǐ qián zhòngyào.}\\
« B : Tu as raison ! L'honnêteté est plus importante que l'argent. »
\end{textebox}

\begin{exoresolu}[Dialogue à reconstituer]
\textbf{Énoncé.} Remettez ce mini-débat dans l'ordre (phrases A à D), puis jouez-le à deux.\\
A. \zh{我不同意。我觉得教育最重要。} \emph{Wǒ bù tóngyì. Wǒ juéde jiàoyù zuì zhòngyào.}\\
B. \zh{你怎么看？} \emph{Nǐ zěnme kàn?}\\
C. \zh{我认为法律最重要。你呢？} \emph{Wǒ rènwéi fǎlǜ zuì zhòngyào. Nǐ ne?}\\
D. \zh{你说得也有道理。} \emph{Nǐ shuō de yě yǒu dàolǐ.}
\tcblower
\textbf{Solution détaillée.}
\begin{enumerate}
\item Le débat s'ouvre par une opinion : c'est la phrase C (\zh{我认为……} + relance \zh{你呢？}).
\item La phrase C se termine par une relance ; le partenaire répond par un désaccord argumenté : phrase A (\zh{我不同意} + \zh{我觉得……}).
\item Pour faire avancer l'échange, on relance à nouveau : phrase B (\zh{你怎么看？}).
\item On clôt poliment en reconnaissant le point de vue de l'autre : phrase D (\zh{你说得也有道理} « ce que tu dis a aussi du sens »).
\end{enumerate}
\textbf{Ordre correct : C → A → B → D.}\\
Dialogue complet :\\
\zh{甲：我认为法律最重要。你呢？} \emph{Jiǎ: Wǒ rènwéi fǎlǜ zuì zhòngyào. Nǐ ne?}\\
\zh{乙：我不同意。我觉得教育最重要。} \emph{Yǐ: Wǒ bù tóngyì. Wǒ juéde jiàoyù zuì zhòngyào.}\\
\zh{甲：你怎么看？} \emph{Jiǎ: Nǐ zěnme kàn?}\\
\zh{乙：你说得也有道理。} \emph{Yǐ: Nǐ shuō de yě yǒu dàolǐ.}\\
« A : Je pense que la loi est le plus important. Et toi ? — B : Je ne suis pas d'accord. Je trouve que l'éducation est le plus important. — A : Qu'en penses-tu (donc) ? — B : Ce que tu dis a aussi du sens. »
\end{exoresolu}

\courssub{Fiche méthode 3 : Prononcer correctement les tons}

\begin{methode}[Sécuriser les tons du lexique « corruption »]
\begin{enumerate}
\item \textbf{Repérer les tons à risque} du thème : \zh{腐败} \emph{fǔbài} (3\textsuperscript{e} + 4\textsuperscript{e} ton), \zh{诚实} \emph{chéngshí} (2\textsuperscript{e} + 2\textsuperscript{e} ton), \zh{法律} \emph{fǎlǜ} (3\textsuperscript{e} + 4\textsuperscript{e} ton), \zh{贪污} \emph{tānwū} (1\textsuperscript{er} + 1\textsuperscript{er} ton).
\item \textbf{Règle des deux 3\textsuperscript{es} tons} : quand deux 3\textsuperscript{es} tons se suivent, le premier devient un 2\textsuperscript{e} ton à l'oral : \zh{你好} se prononce \emph{ní hǎo} (mais s'écrit \emph{nǐ hǎo}) ; de même \zh{很好} se prononce \emph{hén hǎo} (mais s'écrit \emph{hěn hǎo}).
\item \textbf{Le 4\textsuperscript{e} ton} tombe franchement, comme un ordre en français : \zh{不！} \emph{bù!} Ne pas le faire monter.
\item \textbf{Le ton neutre} : dans \zh{我们} \emph{wǒmen}, la deuxième syllabe est légère et courte.
\item \textbf{Entraînement} : lire le mot seul lentement, puis dans une phrase courte, puis à vitesse normale. Enregistrer sa voix et comparer avec le professeur.
\end{enumerate}
\end{methode}

\begin{exoresolu}[Dictée de tons et lecture]
\textbf{Énoncé.} a) Ajoutez les marques de tons sur le pinyin : fu bai (腐败), cheng shi (诚实), fa lü (法律), fan dui (反对). b) Lisez à voix haute : \zh{我们反对腐败，我们要诚实。} en respectant les tons.
\tcblower
\textbf{Solution détaillée.}
\begin{enumerate}
\item \zh{腐败} = \emph{fǔbài} : 3\textsuperscript{e} ton descendant-remontant sur \emph{fǔ}, 4\textsuperscript{e} ton sec sur \emph{bài}. Erreur typique : dire \emph{fùbái} — le sens devient incompréhensible.
\item \zh{诚实} = \emph{chéngshí} : deux 2\textsuperscript{es} tons montants. Attention à ne pas confondre avec \emph{chéngshì} \zh{城市} « ville » : seul le ton change !
\item \zh{法律} = \emph{fǎlǜ} : 3\textsuperscript{e} + 4\textsuperscript{e} ton. Le \emph{ü} de \emph{lǜ} se prononce les lèvres arrondies (comme le « u » français).
\item \zh{反对} = \emph{fǎnduì} : 3\textsuperscript{e} + 4\textsuperscript{e} ton.
\end{enumerate}
\textbf{Lecture de la phrase :} \emph{Wǒmen fǎnduì fǔbài, wǒmen yào chéngshí.} « Nous nous opposons à la corruption, nous voulons être honnêtes. » Points de vigilance : \emph{wǒmen} avec ton neutre sur \emph{men} ; enchaîner \emph{fǎnduì fǔbài} sans hésitation ; \emph{chéngshí} en fin de phrase, bien faire monter les deux tons.
\end{exoresolu}

\courssub{Fiche méthode 4 : Construire des phrases d'opinion correctes}

\begin{methode}[Grammaire de l'opinion : modaux et comparatif]
\begin{enumerate}
\item \textbf{Verbe modal avant le verbe} : \zh{应该} \emph{yīnggāi} « devoir », \zh{要} \emph{yào} « vouloir/devoir », \zh{能} \emph{néng} « pouvoir ». Structure : Sujet + modal + Verbe. \zh{我们应该反对腐败。} « Nous devons combattre la corruption. »
\item \textbf{La négation} se place avant le modal : \zh{我们不应该贪污。} \emph{Wǒmen bù yīnggāi tānwū.} « Nous ne devons pas détourner. »
\item \textbf{Le comparatif avec} \zh{比} : A + 比 + B + adjectif. \zh{诚实比钱重要。} \emph{Chéngshí bǐ qián zhòngyào.} « L'honnêteté est plus importante que l'argent. »
\item \textbf{Le superlatif avec} \zh{最} : \zh{教育最重要。} \emph{Jiàoyù zuì zhòngyào.} « L'éducation est la plus importante. »
\item \textbf{« De plus en plus »} : \zh{越来越} + adjectif : \zh{腐败问题越来越严重。} \emph{Fǔbài wèntí yuèláiyuè yánzhòng.} « Le problème de la corruption devient de plus en plus grave. »
\end{enumerate}
\end{methode}

\begin{center}
\begin{tabularx}{\linewidth}{|X|X|X|}
\hline
\rowcolor{popblueL} Structure & Exemple (caractères + pinyin) & Traduction \\
\hline
Sujet + modal + Verbe & \zh{我们应该反对腐败。} Wǒmen yīnggāi fǎnduì fǔbài. & Nous devons combattre la corruption. \\
\hline
Sujet + 不 + modal + Verbe & \zh{我们不应该贪污。} Wǒmen bù yīnggāi tānwū. & Nous ne devons pas détourner. \\
\hline
A + 比 + B + adjectif & \zh{诚实比钱重要。} Chéngshí bǐ qián zhòngyào. & L'honnêteté est plus importante que l'argent. \\
\hline
Sujet + 最 + adjectif & \zh{教育最重要。} Jiàoyù zuì zhòngyào. & L'éducation est la plus importante. \\
\hline
Sujet + 越来越 + adjectif & \zh{腐败问题越来越严重。} Fǔbài wèntí yuèláiyuè yánzhòng. & Le problème devient de plus en plus grave. \\
\hline
\end{tabularx}
\end{center}

\begin{exoresolu}[Traduire et transformer]
\textbf{Énoncé.} a) Traduisez en chinois : « Nous ne devons pas accepter la corruption. » b) Traduisez : « L'honnêteté est plus importante que l'argent. » c) Transformez \zh{腐败很严重。} avec \zh{越来越}.
\tcblower
\textbf{Solution détaillée.}
\begin{enumerate}
\item a) Sujet \zh{我们} + négation \zh{不} + modal \zh{应该} + verbe \zh{接受} + complément \zh{腐败} : \zh{我们不应该接受腐败。} \emph{Wǒmen bù yīnggāi jiēshòu fǔbài.} Piège : ne pas mettre \zh{不} après le modal (\zh{我们应该不接受} est fautif).
\item b) Structure du comparatif : A + \zh{比} + B + adjectif : \zh{诚实比钱重要。} \emph{Chéngshí bǐ qián zhòngyào.} Piège : jamais \zh{比} + adjectif seul sans les deux termes comparés, et jamais \zh{很} avec \zh{比} (\zh{诚实比钱很重要} est fautif).
\item c) On remplace l'intensif \zh{很} par \zh{越来越} : \zh{腐败越来越严重。} \emph{Fǔbài yuèláiyuè yánzhòng.} « La corruption devient de plus en plus grave. » Règle : avec \zh{越来越}, on ne met jamais \zh{很} devant l'adjectif.
\end{enumerate}
Ces trois phrases peuvent être réutilisées telles quelles dans un exposé ou un débat sur la corruption.
\end{exoresolu}

\courssub{Fiche méthode 5 : Réussir un jeu de rôle guidé}

\begin{methode}[Préparer et jouer une scène en 5 étapes]
\begin{enumerate}
\item \textbf{Lire la consigne} : qui parle ? où ? quel est le problème ? (ex. : deux élèves discutent de la corruption à l'école ou dans l'administration à Yaoundé).
\item \textbf{Répartir les rôles} et noter 3 phrases obligatoires par personnage : une opinion, une réaction (accord/désaccord), une relance.
\item \textbf{Écrire un canevas} (pas un texte complet) : mots-clés en caractères + pinyin.
\item \textbf{Répéter à voix haute} en travaillant les tons des mots-clés (\zh{腐败} fǔbài, \zh{诚实} chéngshí, \zh{反对} fǎnduì).
\item \textbf{Jouer la scène} en regardant son partenaire, avec gestuelle naturelle, et conclure par une phrase de synthèse : \zh{我们都认为，诚实很重要。} \emph{Wǒmen dōu rènwéi, chéngshí hěn zhòngyào.} « Nous pensons tous que l'honnêteté est très importante. »
\end{enumerate}
\end{methode}

\begin{exoresolu}[Jeu de rôle : « À la mairie »]
\textbf{Énoncé.} À deux, jouez la scène suivante : un citoyen (甲) explique qu'on lui a demandé de l'argent pour un papier administratif ; son ami (乙) réagit et propose de refuser et de signaler. Produisez au moins 4 répliques.
\tcblower
\textbf{Solution détaillée (modèle de dialogue).}\\
\zh{甲：昨天我去市政府，有人跟我要钱。} \emph{Jiǎ: Zuótiān wǒ qù shì zhèngfǔ, yǒu rén gēn wǒ yào qián.} « Hier, je suis allé à la mairie, quelqu'un m'a demandé de l'argent. »\\
\zh{乙：这是腐败！你给了吗？} \emph{Yǐ: Zhè shì fǔbài! Nǐ gěi le ma?} « C'est de la corruption ! Tu as donné ? »\\
\zh{甲：没有。我觉得我们不应该给钱。} \emph{Jiǎ: Méiyǒu. Wǒ juéde wǒmen bù yīnggāi gěi qián.} « Non. Je trouve que nous ne devons pas donner d'argent. »\\
\zh{乙：我同意。我们应该告诉老师或者警察。} \emph{Yǐ: Wǒ tóngyì. Wǒmen yīnggāi gàosu lǎoshī huòzhě jǐngchá.} « Je suis d'accord. Nous devons le dire au professeur ou à la police. »\\
\zh{甲：你说得对！诚实比钱重要。} \emph{Jiǎ: Nǐ shuō de duì! Chéngshí bǐ qián zhòngyào.} « Tu as raison ! L'honnêteté est plus importante que l'argent. »\\
Points de grammaire vérifiés : \zh{跟我要钱} (préposition \zh{gēn} + personne + verbe) ; \zh{给了吗} avec \zh{了} pour l'action accomplie et \zh{吗} pour la question ; réponse négative au passé avec \zh{没有} (jamais \zh{不} pour un passé accompli) ; \zh{或者} « ou » dans une phrase affirmative.
\end{exoresolu}

\begin{experience}[Débat guidé de classe : « Faut-il punir sévèrement la corruption ? »]
Par groupes de quatre : deux élèves sont « pour » (\zh{我同意}), deux sont « contre » (\zh{我不同意}). Chaque élève doit prononcer au moins trois phrases : une opinion (\zh{我认为……} ou \zh{我觉得……}), une réaction au partenaire (\zh{你说得对！} ou \zh{对不起，可是……}), une relance (\zh{你呢？} / \zh{你怎么看？}). Le groupe conclut ensemble : \zh{我们都认为，……} « Nous pensons tous que... ». Le professeur évalue la justesse des tons, l'ordre des mots et l'utilisation des formules fonctionnelles.
\end{experience}

\begin{attention}[Les 5 erreurs qui coûtent des points au Bac]
\begin{enumerate}
\item \textbf{Tons négligés à l'oral} : confondre \zh{诚实} \emph{chéngshí} « honnête » et \zh{城市} \emph{chéngshì} « ville ». Un ton faux peut changer le sens et faire perdre tous les points de prononciation.
\item \textbf{Caractères confondus à l'écrit} : \zh{腐败} (fǔbài, corruption) contient l'élément 肉 « chair » (écrit 月) dans \zh{腐} ; ne pas l'écrire \zh{府} (fǔ, gouvernement, comme dans \zh{政府}). Même pinyin, sens différent !
\item \textbf{Place de la négation} : on dit \zh{我们不应该贪污} (négation AVANT le modal), jamais \zh{我们应该不贪污}. Calque du français « nous devons ne pas » à bannir.
\item \textbf{\zh{很} interdit avec \zh{比} et \zh{越来越}} : \zh{诚实比钱很重要} et \zh{越来越很严重} sont fautifs. Le comparatif et \zh{越来越} portent déjà l'intensité.
\item \textbf{Passé accompli} : la négation d'une action passée se fait avec \zh{没有} : \zh{我没有给钱。} « Je n'ai pas donné d'argent. » Dire \zh{我不给钱} signifie « je ne donne pas d'argent » (présent/habitude), pas un refus passé.
\end{enumerate}
\end{attention}

\begin{aretenir}[L'essentiel de la leçon]
\begin{itemize}
\item Lexique clé : \zh{腐败} fǔbài « corruption », \zh{贪污} tānwū « détourner », \zh{诚实} chéngshí « honnête », \zh{反对} fǎnduì « combattre », \zh{法律} fǎlǜ « loi ».
\item Exposé en trois temps : saluer et annoncer (\zh{大家好！今天我要谈谈……}), développer (\zh{首先}、\zh{然后}、\zh{最后}), conclure (\zh{我认为……谢谢大家！}).
\item Débat : avis (\zh{我认为} / \zh{我觉得} / \zh{在我看来}), accord (\zh{我同意你的看法}), désaccord poli (\zh{对不起，可是……}), relance (\zh{你呢？} / \zh{你怎么看？}).
\item Grammaire : modal après le sujet, négation avant le modal ; comparatif A + \zh{比} + B + adjectif ; superlatif avec \zh{最} ; \zh{越来越} + adjectif sans \zh{很}.
\item Prononciation : soigner les tons de \zh{腐败} fǔbài et \zh{诚实} chéngshí ; deux 3\textsuperscript{es} tons qui se suivent : le premier devient un 2\textsuperscript{e} ton.
\end{itemize}
\end{aretenir}

\newpage

\coursec{Exercices}

\courssub{Je vérifie mes connaissances}

\begin{exercice}[QCM : le bon choix]
Pour chaque question, entoure la bonne réponse.
\begin{enumerate}
\item \zh{腐败} (fǔbài) signifie :
\begin{itemize}
\item a) l'honnêteté \quad b) la corruption \quad c) la loi
\end{itemize}
\item Pour donner son avis, on dit :
\begin{itemize}
\item a) \zh{我觉得……} \quad b) \zh{再见} \quad c) \zh{多少钱}
\end{itemize}
\item \zh{我同意你的看法。} (Wǒ tóngyì nǐ de kànfǎ.) signifie :
\begin{itemize}
\item a) Je ne suis pas d'accord. \quad b) Je suis d'accord avec toi. \quad c) Je ne comprends pas.
\end{itemize}
\item Pour relancer la conversation, on peut dire :
\begin{itemize}
\item a) \zh{你呢？} \quad b) \zh{谢谢！} \quad c) \zh{对不起！}
\end{itemize}
\end{enumerate}
\end{exercice}

\begin{exercice}[Vrai ou faux ? Justifie ta réponse]
Dis si chaque affirmation est vraie ou fausse, puis justifie en une phrase en français.
\begin{enumerate}
\item \zh{诚实} (chéngshí) signifie « malhonnête ».
\item Pour exprimer un désaccord poli, on peut dire \zh{你说得有道理，但是……} (Nǐ shuō de yǒu dàolǐ, dànshì...).
\item \zh{法律} (fǎlǜ) signifie « la loi ».
\item Dans un exposé, \zh{最后} (zuìhòu) sert à introduire la première idée.
\end{enumerate}
\end{exercice}

\begin{exercice}[Vocabulaire : associe chaque mot à sa traduction]
Relie chaque mot chinois à sa traduction française.
\begin{center}
\begin{tabularx}{\linewidth}{|X|X|}\hline
\rowcolor{popblueL} Mot chinois & Traduction \\ \hline
\zh{贪污} (tānwū) & pot-de-vin reçu \\ \hline
\zh{受贿} (shòuhuì) & honnête \\ \hline
\zh{诚实} (chéngshí) & détournement (de fonds) \\ \hline
\zh{法律} (fǎlǜ) & impact, influence \\ \hline
\zh{影响} (yǐngxiǎng) & loi \\ \hline
\end{tabularx}
\end{center}
\end{exercice}

\begin{exercice}[Pinyin et caractères : à toi d'associer]
Associe chaque caractère ou mot au bon pinyin avec tons.
\begin{enumerate}
\item \zh{腐败} \hspace{1cm} a) tóngyì
\item \zh{诚实} \hspace{1cm} b) dànshì
\item \zh{同意} \hspace{1cm} c) fǔbài
\item \zh{但是} \hspace{1cm} d) chéngshí
\end{enumerate}
Puis écris la traduction française de chaque mot.
\end{exercice}

\courssub{Je m'entraîne}

\begin{exercice}[Traduis en français]
Traduis les phrases suivantes en français.
\begin{enumerate}
\item \zh{我觉得腐败是一个严重的问题。} (Wǒ juéde fǔbài shì yí ge yánzhòng de wèntí.)
\item \zh{我不同意你的看法。} (Wǒ bù tóngyì nǐ de kànfǎ.)
\item \zh{首先，腐败影响经济的发展。} (Shǒuxiān, fǔbài yǐngxiǎng jīngjì de fāzhǎn.)
\item \zh{我们应该诚实，也应该遵守法律。} (Wǒmen yīnggāi chéngshí, yě yīnggāi zūnshǒu fǎlǜ.)
\end{enumerate}
\end{exercice}

\begin{exercice}[Dialogue à trous : les formules fonctionnelles]
Complète le dialogue avec les mots de la boîte : \zh{我觉得} — \zh{我同意} — \zh{但是} — \zh{你呢} — \zh{所以}.
\begin{textebox}[Dialogue à compléter]
甲：\zh{你觉得腐败是一个严重的问题吗？}\\
(Jiǎ : Nǐ juéde fǔbài shì yí ge yánzhòng de wèntí ma ?)\\
乙：\zh{…………，腐败对我们的国家很不好。…………？}\\
(Yǐ : ..., fǔbài duì wǒmen de guójiā hěn bù hǎo. ... ?)\\
甲：\zh{…………你的看法，…………我觉得年轻人也可以反对腐败。}\\
(Jiǎ : ... nǐ de kànfǎ, ... wǒ juéde niánqīngrén yě kěyǐ fǎnduì fǔbài.)\\
乙：\zh{对！…………我们应该从小诚实。}\\
(Yǐ : Duì ! ... wǒmen yīnggāi cóng xiǎo chéngshí.)
\end{textebox}
\end{exercice}

\begin{exercice}[Remets les mots dans l'ordre]
Remets les mots dans le bon ordre pour faire une phrase correcte, puis traduis-la.
\begin{enumerate}
\item \zh{腐败 / 是 / 问题 / 一个 / 严重的}
\item \zh{我 / 同意 / 不 / 你 / 的 / 看法}
\item \zh{影响 / 腐败 / 经济 / 的 / 发展}
\item \zh{应该 / 我们 / 法律 / 遵守}
\end{enumerate}
\end{exercice}

\begin{exercice}[Réagis à une affirmation]
Ton camarade dit : \zh{腐败不是一个严重的问题。} (Fǔbài bú shì yí ge yánzhòng de wèntí.) — « La corruption n'est pas un problème grave. »

Exprime ton désaccord poliment en deux phrases en chinois : commence par \zh{你说得有道理，但是……} (Nǐ shuō de yǒu dàolǐ, dànshì...), puis donne ton avis avec \zh{我觉得……}. Écris le pinyin sous chaque phrase.
\end{exercice}

\courssub{Je me prépare au Bac}

\begin{exercice}[Texte de compréhension : la corruption au lycée]
Lis le texte suivant, puis réponds aux questions.

\begin{textebox}[腐败与年轻人]
\zh{腐败是一个严重的问题。在很多国家，有些人贪污、受贿，这对经济和社会的发展很不好。年轻人应该怎么做呢？首先，我们应该诚实，不说谎。然后，我们应该遵守法律。最后，如果我们看到腐败，我们应该说出来。我相信，年轻人可以改变世界！}\\
(Fǔbài shì yí ge yánzhòng de wèntí. Zài hěn duō guójiā, yǒuxiē rén tānwū, shòuhuì, zhè duì jīngjì hé shèhuì de fāzhǎn hěn bù hǎo. Niánqīngrén yīnggāi zěnme zuò ne ? Shǒuxiān, wǒmen yīnggāi chéngshí, bù shuōhuǎng. Ránhòu, wǒmen yīnggāi zūnshǒu fǎlǜ. Zuìhòu, rúguǒ wǒmen kàndào fǔbài, wǒmen yīnggāi shuō chūlái. Wǒ xiāngxìn, niánqīngrén kěyǐ gǎibiàn shìjiè !)\\
« La corruption est un problème grave. Dans beaucoup de pays, certaines personnes détournent de l'argent et reçoivent des pots-de-vin, ce qui nuit au développement de l'économie et de la société. Que doivent faire les jeunes ? D'abord, nous devons être honnêtes et ne pas mentir. Ensuite, nous devons respecter la loi. Enfin, si nous voyons de la corruption, nous devons la dénoncer. Je crois que les jeunes peuvent changer le monde ! »
\end{textebox}

\textbf{Questions de compréhension} (réponds en français) :
\begin{enumerate}
\item Quels sont les deux actes de corruption cités dans le texte ?
\item Quelles sont les trois actions que les jeunes doivent faire, selon l'auteur ?
\item Que croit l'auteur à la fin du texte ?
\end{enumerate}

\textbf{Questions de langue} :
\begin{enumerate}
\item Relève dans le texte les trois mots qui organisent l'exposé (\zh{首先}, …) et donne leur traduction.
\item Trouve dans le texte un mot qui signifie « loi » et un mot qui signifie « honnête ».
\item Traduis en français : \zh{如果我们看到腐败，我们应该说出来。}
\end{enumerate}
\end{exercice}

\begin{exercice}[Thème et version]
\textbf{Thème} — Traduis en chinois (caractères et pinyin) :
\begin{enumerate}
\item La corruption est un problème grave.
\item Je ne suis pas d'accord avec toi.
\item Nous devons être honnêtes et respecter la loi.
\end{enumerate}

\textbf{Version} — Traduis en français :
\begin{enumerate}
\item \zh{我觉得年轻人应该反对腐败。} (Wǒ juéde niánqīngrén yīnggāi fǎnduì fǔbài.)
\item \zh{贪污和受贿对社会很不好。} (Tānwū hé shòuhuì duì shèhuì hěn bù hǎo.)
\item \zh{你说得有道理，但是我不完全同意。} (Nǐ shuō de yǒu dàolǐ, dànshì wǒ bù wánquán tóngyì.)
\end{enumerate}
\end{exercice}

\begin{exercice}[Expression orale et écrite : exposé et jeu de rôle]
\textbf{Partie A — Exposé oral (préparation écrite).}

Prépare un exposé d'une minute sur le sujet : \zh{腐败的影响} (Fǔbài de yǐngxiǎng) — « Les effets de la corruption ». Rédige ta fiche de préparation en chinois (entre 60 et 80 caractères) en utilisant obligatoirement les mots \zh{首先}, \zh{然后}, \zh{最后} et \zh{所以}, ainsi qu'au moins deux mots du lexique de la leçon (\zh{贪污}, \zh{受贿}, \zh{诚实}, \zh{法律}, \zh{影响}).

\textbf{Partie B — Jeu de rôle noté.}

Par groupes de deux : l'un joue le journaliste d'une radio de Yaoundé, l'autre un élève de Première. Le journaliste interviewe l'élève sur la corruption au Cameroun. L'élève doit :
\begin{itemize}
\item utiliser au moins deux formules d'opinion (\zh{我觉得……}, \zh{我认为……}) ;
\item faire au moins trois relances ou réactions (\zh{你呢？}, \zh{我同意你的看法。}, \zh{你说得有道理，但是……}) ;
\item prononcer correctement les tons des mots \zh{fǔbài}, \zh{chéngshí}, \zh{tóngyì}, \zh{dànshì}.
\end{itemize}
Rédige d'abord le dialogue complet en chinois avec pinyin (au moins 8 répliques), puis jouez-le devant la classe.
\end{exercice}

\newpage

\coursec{Corrigés des exercices}

\begin{corrige}[Exercice 1 — QCM : le bon choix]
\begin{enumerate}
\item \textbf{b) la corruption.} \zh{腐败} (fǔbài) signifie « la corruption ». L'honnêteté se dit \zh{诚实} (chéngshí) et la loi \zh{法律} (fǎlǜ).
\item \textbf{a)} \zh{我觉得……} (Wǒ juéde...) signifie « je pense que..., je trouve que... » : c'est la formule de base pour donner son avis. \zh{再见} (zàijiàn) = « au revoir » ; \zh{多少钱} (duōshao qián) = « combien ça coûte ? ».
\item \textbf{b) Je suis d'accord avec toi.} \zh{同意} (tóngyì) = « être d'accord » ; \zh{你的看法} (nǐ de kànfǎ) = « ton avis, ton point de vue ».
\item \textbf{a)} \zh{你呢？} (Nǐ ne ?) = « et toi ? » : c'est la relance la plus simple pour rendre la parole à son interlocuteur. \zh{谢谢} = « merci » ; \zh{对不起} = « pardon ».
\end{enumerate}
\end{corrige}

\begin{corrige}[Exercice 2 — Vrai ou faux ? Justifie ta réponse]
\begin{enumerate}
\item \textbf{Faux.} \zh{诚实} (chéngshí) signifie « honnête ». « Malhonnête » se dirait \zh{不诚实} (bù chéngshí) : la négation \zh{不} change le sens.
\item \textbf{Vrai.} \zh{你说得有道理，但是……} (Nǐ shuō de yǒu dàolǐ, dànshì...) signifie « ce que tu dis est sensé, mais... » : on reconnaît d'abord le point de vue de l'autre, puis on introduit le désaccord avec \zh{但是} (dànshì, « mais »). C'est la formule de désaccord poli.
\item \textbf{Vrai.} \zh{法律} (fǎlǜ) signifie bien « la loi ». « Respecter la loi » se dit \zh{遵守法律} (zūnshǒu fǎlǜ).
\item \textbf{Faux.} \zh{最后} (zuìhòu) signifie « enfin, pour finir » : il introduit la dernière idée. La première idée est introduite par \zh{首先} (shǒuxiān, « d'abord »).
\end{enumerate}
\end{corrige}

\begin{corrige}[Exercice 3 — Vocabulaire : associe chaque mot à sa traduction]
\begin{center}
\begin{tabularx}{\linewidth}{|X|X|}\hline
\rowcolor{popblueL} Mot chinois & Traduction correcte \\ \hline
\zh{贪污} (tānwū) & détournement (de fonds) \\ \hline
\zh{受贿} (shòuhuì) & pot-de-vin reçu \\ \hline
\zh{诚实} (chéngshí) & honnête \\ \hline
\zh{法律} (fǎlǜ) & loi \\ \hline
\zh{影响} (yǐngxiǎng) & impact, influence \\ \hline
\end{tabularx}
\end{center}
\textbf{Astuce mémoire :} \zh{贪} (tān) contient l'idée d'« avidité » ; \zh{受} (shòu) signifie « recevoir » : \zh{受贿} = « recevoir un pot-de-vin ».
\end{corrige}

\begin{corrige}[Exercice 4 — Pinyin et caractères : à toi d'associer]
\begin{enumerate}
\item \zh{腐败} — \textbf{c) fǔbài} — « la corruption » (3\ieme{} ton + 4\ieme{} ton).
\item \zh{诚实} — \textbf{d) chéngshí} — « honnête » (2\ieme{} ton + 2\ieme{} ton).
\item \zh{同意} — \textbf{a) tóngyì} — « être d'accord » (2\ieme{} ton + 4\ieme{} ton).
\item \zh{但是} — \textbf{b) dànshì} — « mais » (4\ieme{} ton + 4\ieme{} ton).
\end{enumerate}
\textbf{Point de vigilance :} ne pas confondre \zh{tóngyì} (être d'accord) et \zh{注意} (zhùyì, faire attention) ; les tons changent tout le sens.
\end{corrige}

\begin{corrige}[Exercice 5 — Traduis en français]
\begin{enumerate}
\item \zh{我觉得腐败是一个严重的问题。} — « Je pense que la corruption est un problème grave. » (\zh{觉得} = penser, trouver ; \zh{严重} = grave, sérieux.)
\item \zh{我不同意你的看法。} — « Je ne suis pas d'accord avec ton avis. » (La négation \zh{不} se place directement avant le verbe \zh{同意}.)
\item \zh{首先，腐败影响经济的发展。} — « D'abord, la corruption nuit au développement de l'économie. » (\zh{影响} = influencer, nuire à ; \zh{的} relie \zh{经济} et \zh{发展} : « le développement de l'économie ».)
\item \zh{我们应该诚实，也应该遵守法律。} — « Nous devons être honnêtes, et nous devons aussi respecter la loi. » (\zh{应该} = devoir ; \zh{也} = aussi, placé avant \zh{应该}.)
\end{enumerate}
\end{corrige}

\begin{corrige}[Exercice 6 — Dialogue à trous : les formules fonctionnelles]
Dialogue complété :
\begin{itemize}
\item 乙：\zh{\textbf{我觉得}，腐败对我们的国家很不好。\textbf{你呢}？}\\
(Yǐ : Wǒ juéde, fǔbài duì wǒmen de guójiā hěn bù hǎo. Nǐ ne ?)\\
« Je pense que la corruption est très mauvaise pour notre pays. Et toi ? »
\item 甲：\zh{\textbf{我同意}你的看法，\textbf{但是}我觉得年轻人也可以反对腐败。}\\
(Jiǎ : Wǒ tóngyì nǐ de kànfǎ, dànshì wǒ juéde niánqīngrén yě kěyǐ fǎnduì fǔbài.)\\
« Je suis d'accord avec ton avis, mais je pense que les jeunes peuvent aussi lutter contre la corruption. »
\item 乙：\zh{对！\textbf{所以}我们应该从小诚实。}\\
(Yǐ : Duì ! Suǒyǐ wǒmen yīnggāi cóng xiǎo chéngshí.)\\
« C'est vrai ! C'est pourquoi nous devons être honnêtes dès le plus jeune âge. »
\end{itemize}
\textbf{Justification :} \zh{我觉得} introduit l'avis ; \zh{你呢} est la relance ; \zh{我同意} marque l'accord ; \zh{但是} introduit une nuance ; \zh{所以} conclut logiquement.
\end{corrige}

\begin{corrige}[Exercice 7 — Remets les mots dans l'ordre]
\begin{enumerate}
\item \zh{腐败是一个严重的问题。} (Fǔbài shì yí ge yánzhòng de wèntí.) — « La corruption est un problème grave. » Structure : Sujet + \zh{是} + classificateur \zh{个} + adjectif + \zh{的} + nom.
\item \zh{我不同意你的看法。} (Wǒ bù tóngyì nǐ de kànfǎ.) — « Je ne suis pas d'accord avec ton avis. » La négation \zh{不} précède le verbe.
\item \zh{腐败影响经济的发展。} (Fǔbài yǐngxiǎng jīngjì de fāzhǎn.) — « La corruption nuit au développement de l'économie. »
\item \zh{我们应该遵守法律。} (Wǒmen yīnggāi zūnshǒu fǎlǜ.) — « Nous devons respecter la loi. » L'auxiliaire \zh{应该} se place avant le verbe \zh{遵守}.
\end{enumerate}
\textbf{Point de vigilance :} en chinois, l'ordre est Sujet – Verbe – Complément ; on ne déplace jamais le complément avant le verbe comme en français accentué.
\end{corrige}

\begin{corrige}[Exercice 8 — Réagis à une affirmation]
Proposition de réponse :
\begin{itemize}
\item \zh{你说得有道理，但是我不这样认为。}\\
(Nǐ shuō de yǒu dàolǐ, dànshì wǒ bù zhèyàng rènwéi.)\\
« Ce que tu dis est sensé, mais je ne le vois pas ainsi. »
\item \zh{我觉得腐败是一个很严重的问题，因为它影响经济的发展。}\\
(Wǒ juéde fǔbài shì yí ge hěn yánzhòng de wèntí, yīnwèi tā yǐngxiǎng jīngjì de fāzhǎn.)\\
« Je pense que la corruption est un problème très grave, parce qu'elle nuit au développement de l'économie. »
\end{itemize}
\textbf{Méthode :} 1) commencer par la formule de politesse \zh{你说得有道理，但是……} ; 2) donner son avis avec \zh{我觉得……} ; 3) si possible, justifier avec \zh{因为} (yīnwèi, « parce que »). Toute variante correcte et polie est acceptée.
\end{corrige}

\begin{corrige}[Exercice 9 — Texte de compréhension : la corruption au lycée]
\textbf{Questions de compréhension}
\begin{enumerate}
\item Les deux actes de corruption cités sont le détournement de fonds (\zh{贪污}, tānwū) et les pots-de-vin reçus (\zh{受贿}, shòuhuì).
\item Selon l'auteur, les jeunes doivent : d'abord être honnêtes et ne pas mentir ; ensuite respecter la loi ; enfin dénoncer la corruption quand ils la voient.
\item À la fin, l'auteur croit que les jeunes peuvent changer le monde (\zh{年轻人可以改变世界}).
\end{enumerate}
\textbf{Questions de langue}
\begin{enumerate}
\item Les trois mots organisateurs sont : \zh{首先} (shǒuxiān, « d'abord »), \zh{然后} (ránhòu, « ensuite »), \zh{最后} (zuìhòu, « enfin »).
\item « Loi » = \zh{法律} (fǎlǜ) ; « honnête » = \zh{诚实} (chéngshí).
\item \zh{如果我们看到腐败，我们应该说出来。} — « Si nous voyons de la corruption, nous devons la dénoncer (la dire à voix haute). » Structure : \zh{如果} + proposition, + \zh{应该} + verbe.
\end{enumerate}
\textbf{Barème indicatif (sur 10) :} compréhension 5 pts (2 + 2 + 1) ; questions de langue 5 pts (1,5 + 1,5 + 2).
\end{corrige}

\begin{corrige}[Exercice 10 — Thème et version]
\textbf{Thème}
\begin{enumerate}
\item \zh{腐败是一个严重的问题。} (Fǔbài shì yí ge yánzhòng de wèntí.)
\item \zh{我不同意你的看法。} (Wǒ bù tóngyì nǐ de kànfǎ.)
\item \zh{我们应该诚实，也应该遵守法律。} (Wǒmen yīnggāi chéngshí, yě yīnggāi zūnshǒu fǎlǜ.)
\end{enumerate}
\textbf{Version}
\begin{enumerate}
\item « Je pense que les jeunes doivent lutter contre la corruption. » (\zh{反对} = s'opposer à, lutter contre.)
\item « Le détournement de fonds et les pots-de-vin sont très mauvais pour la société. » (\zh{和} = et ; \zh{对……不好} = être mauvais pour.)
\item « Ce que tu dis est sensé, mais je ne suis pas tout à fait d'accord. » (\zh{完全} = complètement ; \zh{不完全同意} = ne pas être entièrement d'accord.)
\end{enumerate}
\textbf{Barème indicatif (sur 10) :} thème 5 pts (1,5 + 1,5 + 2 : caractères 1 pt, pinyin avec tons 0,5 pt, grammaire 0,5 pt par phrase) ; version 5 pts. \textbf{Points de vigilance :} place de \zh{不} avant le verbe ; classificateur \zh{个} obligatoire dans \zh{一个问题} ; \zh{也} avant \zh{应该}.
\end{corrige}

\begin{corrige}[Exercice 11 — Expression orale et écrite : exposé et jeu de rôle]
\textbf{Partie A — Proposition de fiche d'exposé (modèle)}

\zh{腐败的影响很大。首先，腐败影响经济的发展：有些人贪污、受贿，国家就没有钱建设学校和医院。然后，腐败影响社会：人们不再相信法律。最后，腐败伤害年轻人，因为他们看不到未来。所以，我们应该诚实，遵守法律，反对腐败。}\\
(Fǔbài de yǐngxiǎng hěn dà. Shǒuxiān, fǔbài yǐngxiǎng jīngjì de fāzhǎn : yǒuxiē rén tānwū, shòuhuì, guójiā jiù méiyǒu qián jiànshè xuéxiào hé yīyuàn. Ránhòu, fǔbài yǐngxiǎng shèhuì : rénmen bú zài xiāngxìn fǎlǜ. Zuìhòu, fǔbài shānghài niánqīngrén, yīnwèi tāmen kàn bu dào wèilái. Suǒyǐ, wǒmen yīnggāi chéngshí, zūnshǒu fǎlǜ, fǎnduì fǔbài.)\\
« Les effets de la corruption sont grands. D'abord, la corruption nuit au développement de l'économie : certains détournent des fonds et reçoivent des pots-de-vin, et l'État n'a plus d'argent pour construire des écoles et des hôpitaux. Ensuite, elle nuit à la société : les gens ne croient plus en la loi. Enfin, elle blesse les jeunes, car ils ne voient pas d'avenir. C'est pourquoi nous devons être honnêtes, respecter la loi et lutter contre la corruption. »

\textbf{Partie B — Proposition de dialogue modèle}
\begin{itemize}
\item 记者：\zh{你好！你觉得腐败是喀麦隆的一个严重问题吗？} (Jìzhě : Nǐ hǎo ! Nǐ juéde fǔbài shì Kāmàilóng de yí ge yánzhòng wèntí ma ?)
\item 学生：\zh{我觉得是的。腐败影响我们国家的发展。} (Xuésheng : Wǒ juéde shì de. Fǔbài yǐngxiǎng wǒmen guójiā de fāzhǎn.)
\item 记者：\zh{你能举个例子吗？} (Nǐ néng jǔ ge lìzi ma ?)
\item 学生：\zh{比如，有些人受贿，这对社会很不好。} (Bǐrú, yǒuxiē rén shòuhuì, zhè duì shèhuì hěn bù hǎo.)
\item 记者：\zh{我同意你的看法。但是年轻人能做什么呢？} (Wǒ tóngyì nǐ de kànfǎ. Dànshì niánqīngrén néng zuò shénme ne ?)
\item 学生：\zh{我认为我们应该诚实，遵守法律。你呢？} (Wǒ rènwéi wǒmen yīnggāi chéngshí, zūnshǒu fǎlǜ. Nǐ ne ?)
\item 记者：\zh{你说得有道理，但是只靠年轻人够吗？} (Nǐ shuō de yǒu dàolǐ, dànshì zhǐ kào niánqīngrén gòu ma ?)
\item 学生：\zh{不够，所以大家都应该反对腐败。谢谢你的采访！} (Bù gòu, suǒyǐ dàjiā dōu yīnggāi fǎnduì fǔbài. Xièxie nǐ de cǎifǎng !)
\end{itemize}
\textbf{Barème indicatif (sur 20) :} richesse du lexique et exactitude des formules 6 pts ; grammaire et ordre des mots 5 pts ; prononciation et tons (\zh{fǔbài}, \zh{chéngshí}, \zh{tóngyì}, \zh{dànshì}) 5 pts ; fluidité, interaction et relances 4 pts.
\textbf{Points de vigilance :} bien marquer le 3\ieme{} ton descendant-remontant de \zh{fǔ} ; ne pas dire \zh{*我同意你} seul sans \zh{的看法} ; garder \zh{但是} après la virgule, jamais en début de réplique sans formule de politesse.
\end{corrige}

\newpage

\coursec{Fiche bilan}

\begin{popfigure}
\begin{center}\tcbox[colback=popink,colframe=popink,coltext=white,arc=9pt,boxsep=4pt,left=6pt,right=6pt]{\bfseries\small Expression orale : la corruption \zh{腐败}}\end{center}
\vspace{-2pt}
\begin{tcbraster}[raster columns=3,raster equal height=rows,raster column skip=6pt,raster row skip=6pt,enhanced,blankest]
\begin{tcolorbox}[enhanced,colback=poporange,colframe=poporange,coltext=white,boxrule=0.9pt,arc=3pt,left=4pt,right=4pt,top=3pt,bottom=3pt,boxsep=1pt,halign=left]\bfseries\scriptsize Lexique clé \zh{腐败 受贿 诚实}\end{tcolorbox}
\begin{tcolorbox}[enhanced,colback=popgreen,colframe=popgreen,coltext=white,boxrule=0.9pt,arc=3pt,left=4pt,right=4pt,top=3pt,bottom=3pt,boxsep=1pt,halign=left]\bfseries\scriptsize Donner son avis \zh{我认为…}\end{tcolorbox}
\begin{tcolorbox}[enhanced,colback=poppurple,colframe=poppurple,coltext=white,boxrule=0.9pt,arc=3pt,left=4pt,right=4pt,top=3pt,bottom=3pt,boxsep=1pt,halign=left]\bfseries\scriptsize Accord / désaccord \zh{我同意 / 我不同意}\end{tcolorbox}
\begin{tcolorbox}[enhanced,colback=popgold,colframe=popgold,coltext=white,boxrule=0.9pt,arc=3pt,left=4pt,right=4pt,top=3pt,bottom=3pt,boxsep=1pt,halign=left]\bfseries\scriptsize Relancer le débat \zh{你觉得呢？}\end{tcolorbox}
\begin{tcolorbox}[enhanced,colback=poppink,colframe=poppink,coltext=white,boxrule=0.9pt,arc=3pt,left=4pt,right=4pt,top=3pt,bottom=3pt,boxsep=1pt,halign=left]\bfseries\scriptsize Présenter un exposé \zh{大家好，今天…}\end{tcolorbox}
\begin{tcolorbox}[enhanced,colback=popblue!70,colframe=popblue!70,coltext=white,boxrule=0.9pt,arc=3pt,left=4pt,right=4pt,top=3pt,bottom=3pt,boxsep=1pt,halign=left]\bfseries\scriptsize Tons et prononciation\end{tcolorbox}
\end{tcbraster}

\legende{Carte mentale de la leçon : parler de la corruption en chinois}
\end{popfigure}

\begin{bilan}
\textbf{1. Lexique clé à connaître par cœur}
\begin{center}
\begin{tabularx}{\linewidth}{|X|X|X|X|}
\hline
\rowcolor{popblueL} Caractères & Pinyin & Nature & Traduction \\
\hline
腐败 & fǔbài & n./adj. & corruption ; corrompu \\
\hline
受贿 & shòuhuì & v. & accepter des pots-de-vin \\
\hline
行贿 & xínghuì & v. & offrir des pots-de-vin \\
\hline
诚实 & chéngshí & adj. & honnête \\
\hline
法律 & fǎlǜ & n. & loi \\
\hline
反对 & fǎnduì & v. & s'opposer à, combattre \\
\hline
影响 & yǐngxiǎng & n./v. & influence ; influencer \\
\hline
发展 & fāzhǎn & n./v. & développement ; se développer \\
\hline
\end{tabularx}
\end{center}

\textbf{2. Formules fonctionnelles du débat}
\begin{center}
\begin{tabularx}{\linewidth}{|X|X|X|}
\hline
\rowcolor{popblueL} Fonction & Formule chinoise + pinyin & Traduction \\
\hline
Donner son avis & 我认为… (Wǒ rènwéi…) & Je pense que… \\
\hline
Être d'accord & 我同意你的看法。(Wǒ tóngyì nǐ de kànfǎ.) & Je suis d'accord avec toi. \\
\hline
Ne pas être d'accord & 我不同意，因为… (Wǒ bù tóngyì, yīnwèi…) & Je ne suis pas d'accord, parce que… \\
\hline
Relancer & 你觉得呢？(Nǐ juéde ne?) & Et toi, qu'en penses-tu ? \\
\hline
Ouvrir un exposé & 大家好！今天我谈谈… (Dàjiā hǎo! Jīntiān wǒ tántan…) & Bonjour à tous ! Aujourd'hui je parle de… \\
\hline
Conclure & 我的讲话完了，谢谢！(Wǒ de jiǎnghuà wán le, xièxie!) & Mon exposé est terminé, merci ! \\
\hline
\end{tabularx}
\end{center}

\textbf{3. Règles de grammaire à connaître par cœur}
\begin{itemize}
\item \cle{因为…，所以…} (\emph{yīnwèi… suǒyǐ…}) : « parce que… donc… ». Ex. : \zh{因为腐败影响国家的发展，所以我们要反对腐败。} (\emph{Yīnwèi fǔbài yǐngxiǎng guójiā de fāzhǎn, suǒyǐ wǒmen yào fǎnduì fǔbài.}) « Parce que la corruption nuit au développement du pays, nous devons la combattre. »
\item \cle{应该 + verbe} (\emph{yīnggāi}) : « devoir » (conseil moral). Ex. : \zh{我们应该诚实。} (\emph{Wǒmen yīnggāi chéngshí.}) « Nous devons être honnêtes. »
\item \cle{对…有影响} : « avoir une influence sur… ». Ex. : \zh{腐败对经济有影响。} (\emph{Fǔbài duì jīngjì yǒu yǐngxiǎng.}) « La corruption a un impact sur l'économie. »
\item Ordre des mots : Sujet + verbe d'opinion (\zh{认为 / 觉得}) + proposition complète, sans « que » en chinois.
\end{itemize}

\textbf{4. Méthode express pour un exposé simple}
\begin{enumerate}
\item Saluer : \zh{大家好！}
\item Annoncer le sujet : \zh{今天我谈谈腐败问题。}
\item Donner 2 ou 3 arguments avec \zh{因为…所以…} et \zh{我认为…}.
\item Conclure et remercier : \zh{我的讲话完了，谢谢大家！}
\end{enumerate}
\end{bilan}

\begin{aretenir}[Checklist avant le Bac]
\begin{itemize}
\item[$\square$] Je connais le lexique : \zh{腐败、受贿、行贿、诚实、法律、反对、影响}.
\item[$\square$] Je sais donner mon avis avec \zh{我认为…} et \zh{我觉得…}.
\item[$\square$] Je sais exprimer l'accord : \zh{我同意你的看法。}
\item[$\square$] Je sais exprimer le désaccord poliment : \zh{我不同意，因为…}.
\item[$\square$] Je sais relancer le débat : \zh{你觉得呢？}
\item[$\square$] Je maîtrise la structure \zh{因为…，所以…}.
\item[$\square$] Je sais utiliser \zh{应该 + verbe} pour donner un conseil.
\item[$\square$] Je prononce correctement les tons de \zh{fǔbài} (3\textsuperscript{e} + 4\textsuperscript{e} ton) et \zh{chéngshí} (2\textsuperscript{e} + 2\textsuperscript{e} ton).
\item[$\square$] Je sais ouvrir et conclure un exposé en chinois.
\item[$\square$] Je peux parler une minute sans lire mes notes.
\item[$\square$] Je respecte l'ordre des mots : sujet + verbe + complément.
\item[$\square$] Je n'oublie pas le \zh{的} dans \zh{国家的发展} et \zh{你的看法}.
\end{itemize}
\end{aretenir}

\findecours

\end{document}
